% concatenated from Without_symmetry.tex
\documentclass[a4paper, 11pt]{article}
\usepackage [a4paper,left=2.5cm,bottom=2.5cm,right=2.5cm,top=2.5cm]{geometry}

% Language setting
% Replace `english' with e.g. `spanish' to change the document language
\usepackage[english]{babel}

% Set page size and margins
% Replace `letterpaper' with `a4paper' for UK/EU standard size
%\usepackage[letterpaper,top=2cm,bottom=2cm,left=3cm,right=3cm,marginparwidth=1.75cm]{geometry}

% Useful packages
%\documentclass[11pt]{amsart}

%% Style dans MesMacros alias dans Library/texmf/tex

%Begin0
\usepackage[utf8]{inputenc} 	% LaTeX, comprend les accents !
%\usepackage[latin1]{inputenc}
\usepackage[T1]{fontenc}
\usepackage{lmodern}
\usepackage{yhmath}
\usepackage{comment}
\usepackage{bm}
\usepackage{enumerate}   

\usepackage{amsthm,amsmath,amssymb,amsbsy,bbm,mathrsfs,supertabular,
eurosym,graphicx,enumitem,xcolor, dsfont}
\usepackage{mathtools}
\usepackage{ulem}

\newtheoremstyle{BBstyle0}  {}{}{\itshape}{}{\bfseries}{}{6pt}{}
\newtheoremstyle{BBstyle1}  {3pt}{3pt}{\rmfamily}{}{\itshape}{: }{3pt}{}
\newtheoremstyle{BBstyle2}  {3pt}{3pt}{\itshape}{}{\bfseries\large}{}{0pt}{}
\newtheoremstyle{BBstyle3}  {}{}{\itshape}{}{\bfseries}{: }{3pt}{}
\newtheoremstyle{BBstyle4}  {}{}{\rmfamily}{}{\bfseries}{}{6pt}{}
  
%\usepackage[normalem]{ulem} %pour hachurer, souligner etc. 
% \uline{texte} pour souligner, \uuline{texte}  pour souligner 2 fois, \sout{texte} pour overlinerer un texte. 
%
% Style biblio pour avoir autheur(date)



%\usepackage[authoryear]{natbib}


%\bibliography{biblio}
% Style de la bibliographie
%\bibliographystyle{apalike}
% Utiliser \cite{autheur}
%end begin0


%Begin AN
%%%% \'ENONC\'ES Anglais
%%% Les styles des \'enonc\'es ne sont pas toujours ceux des thm.
\newtheorem{thm}{Theorem}
\newtheorem{lem}{Lemma}
\newtheorem{lemma}{Lemma}
%\newtheorem{proposition}{Proposition}
\newtheorem{prop}{Proposition}
\newtheorem{df}{Definition}
\newtheorem{cor}{Corollary}
\newtheorem{ass}{Assumption}
\newtheorem{pb}{Problem}
\newtheorem{prty}{Property}
\newtheorem{rem}{Remark}
%\renewcommand{\therem}{}

% Chgt de style: nom en gras et contenu en ecriture normale Roman (non en italique)
\theoremstyle{definition}
\newtheorem{exa}{Example}
\newtheorem{exe}{Exercise}

% Pour un titre en italique et contenu en \'ecriture normale Roman (comme pour 'proof')
%\theoremstyle{remark}

\usepackage[english]{babel}
%End AN

%\usepackage{showkeys}
\usepackage{enumitem}
\usepackage{hyperref}  % doit imp\'erativement \^etre apr\`es "showkeys" ! et avant Core pour avoir les \L_{p} \'ecrit correctement
%%%%%%%
%%%
%\input{Core}	

%\usepackage{refcheck}


% Parenth\`eses et autres
\newcommand{\pa}[1]{\left({#1}\right)}
\newcommand{\norm}[1]{\left\|{#1}\right\|}
\newcommand{\cro}[1]{\left[{#1}\right]}
\newcommand{\ab}[1]{\left|{#1}\right|}
\newcommand{\ac}[1]{\left\{{#1}\right\}}
\newcommand{\gui}[1]{\og #1\fg}
%\newcommand{\id}[2]{[\![#1,#2]\!]} 
\newcommand{\id}[2]{\{#1,\dots,#2\}} 

%%% fonctions
\newcommand{\argmin}{\mathop{\rm argmin}}
\newcommand{\argmax}{\mathop{\rm argmax}}
\newcommand{\esssup}{\mathop{\rm ess\ sup}}
\newcommand{\essinf}{\mathop{\rm ess\ inf}}

%%% operateurs
\newcommand{\pen}{\mathop{\rm pen}\nolimits}
\newcommand{\Var}{\mathop{\rm Var}\nolimits}
\newcommand{\Cov}{\mathop{\rm Cov}\nolimits}
\newcommand{\Tr}{\mathop{\rm Tr}\nolimits}
\newcommand{\Card}{\mathop{\rm Card}\nolimits}
\newcommand{\Diag}{\mathop{\rm Diag}\nolimits}
\newcommand{\Gui}[1]{\textcolor{black}{Gui: #1}}
\newcommand{\Hel}[1]{\textcolor{blue}{H: #1}}

% fleches
\newcommand{\CL}{\mathop{\leadsto}}
\newcommand{\fleche}{\mathop{\longrightarrow}}
\newcommand{\dfleche}[1]{\,\displaystyle{\mathop{\longrightarrow}_{#1}}\,}
\newcommand{\CP}[1]{\stackrel{\mathrm{P}}{\dfleche{#1}}}
\newcommand{\CPS}[1]{\stackrel{\mathrm{p.s.}}{\dfleche{#1}}}
\newcommand{\CV}[1]{\dfleche{#1}}
\newcommand{\CVL}[1]{\stackrel{\sL}{\dfleche{#1}}}
\newcommand{\pp}[1]{\stackrel{{#1}\textrm{-}\mathrm{p.p.}}{\,=\,}}

% Couleurs 
%\RequirePackage[colorlinks,citecolor=blue,urlcolor=blue]{hyperref}
\newcommand{\corr}[1]{\textcolor{black}{#1}}
\newcommand{\corb}[1]{\textcolor{blue}{#1}}
\newcommand{\corv}[1]{\textcolor{green}{#1}}
\newcommand{\corm}[1]{\textcolor{magenta}{#1}}
\newcommand{\coro}[1]{\textcolor{black}{#1}}
\newcommand{\corbr}[1]{\textcolor{brown}{#1}}
\newcommand{\corlg}[1]{\textcolor{lightgray}{#1}}
\newcommand{\corvi}[1]{\textcolor{violet}{#1}}
\newcommand{\cort}[1]{\textcolor{teal}{#1}}
\newcommand{\corol}[1]{\textcolor{olive}{#1}}
\newcommand{\corj}[1]{\textcolor{yellow}{#1}}
\newcommand{\corc}[1]{\textcolor{cyan}{#1}}
\newcommand{\corg}[1]{\textcolor{gray}{#1}}


% Mathbb
\newcommand{\B}{{\mathbb{B}}}
\newcommand{\C}{{\mathbb{C}}}
\newcommand{\E}{{\mathbb{E}}}
\newcommand{\F}{{\mathbb{F}}} 
\renewcommand{\L}{{\mathbb{L}}}
\newcommand{\N}{{\mathbb{N}}}
\renewcommand{\P}{{\mathbb{P}}}
\newcommand{\Q}{{\mathbb{Q}}} 
\newcommand{\R}{{\mathbb{R}}}
\renewcommand{\S}{{\mathbb{S}}} 
\newcommand{\T}{{\mathbb{T}}} 
\newcommand{\X}{{\mathbb{X}}}
\newcommand{\Z}{{\mathbb{Z}}}


% Mathcal
\newcommand{\cA}{{\mathcal{A}}}
\newcommand{\cB}{{\mathcal{B}}}
\newcommand{\cC}{{\mathcal{C}}}
\newcommand{\cD}{{\mathcal{D}}}
\newcommand{\cE}{{\mathcal{E}}}
\newcommand{\cF}{{\mathcal{F}}}
\newcommand{\cG}{{\mathcal{G}}} 
\newcommand{\cH}{{\mathcal{H}}}
\newcommand{\cI}{{\mathcal{I}}}
\newcommand{\cJ}{{\mathcal{J}}}
\newcommand{\cK}{{\mathcal{K}}}
\newcommand{\cL}{{\mathcal{L}}} 
\newcommand{\cM}{{\mathcal{M}}}
\newcommand{\cN}{{\mathcal{N}}}
\newcommand{\cO}{{\mathcal{O}}}
\newcommand{\cP}{{\mathcal{P}}}
\newcommand{\cQ}{{\mathcal{Q}}} 
\newcommand{\cR}{{\mathcal{R}}}
\newcommand{\cS}{{\mathcal{S}}} 
\newcommand{\cT}{{\mathcal{T}}}
\newcommand{\cU}{{\mathcal{U}}}
\newcommand{\cV}{{\mathcal{V}}}
\newcommand{\cW}{{\mathcal{W}}}
\newcommand{\cX}{{\mathcal{X}}}
\newcommand{\cY}{{\mathcal{Y}}} 
\newcommand{\cZ}{{\mathcal{Z}}}

% Mathcal en gras
\newcommand{\cbA}{\boldsymbol{\mathcal{A}}}
\newcommand{\cbB}{\boldsymbol{\mathcal{B}}}
\newcommand{\cbC}{\boldsymbol{\mathcal{C}}}
\newcommand{\cbD}{\boldsymbol{\mathcal{D}}}
\newcommand{\cbE}{\boldsymbol{\mathcal{E}}}
\newcommand{\cbF}{\boldsymbol{\mathcal{F}}}
\newcommand{\cbG}{\boldsymbol{\mathcal{G}}}
\newcommand{\cbH}{\boldsymbol{\mathcal{H}}}
\newcommand{\cbI}{\boldsymbol{\mathcal{I}}}
\newcommand{\cbJ}{\boldsymbol{\mathcal{J}}}
\newcommand{\cbK}{\boldsymbol{\mathcal{K}}}
\newcommand{\cbL}{\boldsymbol{\mathcal{L}}}
\newcommand{\cbM}{\boldsymbol{\mathcal{M}}}
\newcommand{\cbN}{\boldsymbol{\mathcal{N}}}
\newcommand{\cbO}{\boldsymbol{\mathcal{O}}}
\newcommand{\cbP}{\boldsymbol{\mathcal{P}}}
\newcommand{\cbQ}{\boldsymbol{\mathcal{Q}}}
\newcommand{\cbR}{\boldsymbol{\mathcal{R}}}
\newcommand{\cbS}{\boldsymbol{\mathcal{S}}}
\newcommand{\cbT}{\boldsymbol{\mathcal{T}}}
\newcommand{\cbU}{\boldsymbol{\mathcal{U}}}
\newcommand{\cbV}{\boldsymbol{\mathcal{V}}}
\newcommand{\cbW}{\boldsymbol{\mathcal{W}}}
\newcommand{\cbX}{\boldsymbol{\mathcal{X}}}
\newcommand{\cbY}{\boldsymbol{\mathcal{Y}}}
\newcommand{\cbZ}{\boldsymbol{\mathcal{Z}}}


%%%%% LETTRES en GRAS
%% minuscules
\newcommand{\ga}{{\mathbf{a}}}
\newcommand{\gb}{{\mathbf{b}}}
\newcommand{\gc}{{\mathbf{c}}}
\newcommand{\gd}{{\mathbf{d}}}
\newcommand{\gee}{{\mathbf{e}}} 
\newcommand{\gf}{{\mathbf{f}}}
\renewcommand{\gg}{{\mathbf{g}}} 
\newcommand{\gh}{{\mathbf{h}}}
\newcommand{\gi}{{\mathbf{i}}}
\newcommand{\gj}{{\mathbf{j}}}
\newcommand{\gk}{{\mathbf{k}}}
\newcommand{\gl}{{\mathbf{l}}} 
\newcommand{\gm}{{\mathbf{m}}}
\newcommand{\gn}{{\mathbf{n}}}
\newcommand{\go}{{\mathbf{o}}}
\newcommand{\gp}{{\mathbf{p}}}
\newcommand{\gq}{{\mathbf{q}}} 
\newcommand{\gr}{{\mathbf{r}}}
\newcommand{\gs}{{\mathbf{s}}} 
\newcommand{\gt}{{\mathbf{t}}}
\newcommand{\gu}{{\mathbf{u}}}
\newcommand{\gv}{{\mathbf{v}}}
\newcommand{\gw}{{\mathbf{w}}}
\newcommand{\gx}{{\mathbf{x}}}
\newcommand{\gy}{{\mathbf{y}}} 
\newcommand{\gz}{{\mathbf{z}}}

%% majuscules
\newcommand{\gA}{{\mathbf{A}}}
\newcommand{\gB}{{\mathbf{B}}}
\newcommand{\gC}{{\mathbf{C}}}
\newcommand{\gD}{{\mathbf{D}}}
\newcommand{\gE}{{\mathbf{E}}}
\newcommand{\gF}{{\mathbf{F}}}
\newcommand{\gG}{{\mathbf{G}}} 
\newcommand{\gH}{{\mathbf{H}}}
\newcommand{\gI}{{\mathbf{I}}}
\newcommand{\gJ}{{\mathbf{J}}}
\newcommand{\gK}{{\mathbf{K}}}
\newcommand{\gL}{{\mathbf{L}}} 
\newcommand{\gM}{{\mathbf{M}}}
\newcommand{\gN}{{\mathbf{N}}}
\newcommand{\gO}{{\mathbf{O}}}
\newcommand{\gP}{{\mathbf{P}}}
\newcommand{\gQ}{{\mathbf{Q}}} 
\newcommand{\gR}{{\mathbf{R}}}
\newcommand{\gS}{{\mathbf{S}}} 
\newcommand{\gT}{{\mathbf{T}}}
\newcommand{\gU}{{\mathbf{U}}}
\newcommand{\gV}{{\mathbf{V}}}
\newcommand{\gW}{{\mathbf{W}}}
\newcommand{\gX}{{\mathbf{X}}}
\newcommand{\gY}{{\mathbf{Y}}} 
\newcommand{\gZ}{{\mathbf{Z}}}
%\newcommand{\modch}{\color{red}}

%% symboles math\'ematiques en gras
\newcommand{\bs}[1]{\boldsymbol{#1}}
\newcommand{\bsA}{{\bs{A}}}
\newcommand{\bsB}{{\bs{B}}}
\newcommand{\bsC}{{\bs{C}}}
\newcommand{\bsD}{{\bs{D}}}
\newcommand{\bsE}{{\bs{E}}}
\newcommand{\bsF}{{\bs{F}}}
\newcommand{\bsG}{{\bs{G}}} 
\newcommand{\bsH}{{\bs{H}}}
\newcommand{\bsI}{{\bs{I}}}
\newcommand{\bsJ}{{\bs{J}}}
\newcommand{\bsK}{{\bs{K}}}
\newcommand{\bsL}{{\bs{L}}} 
\newcommand{\bsM}{{\bs{M}}}
\newcommand{\bsN}{{\bs{N}}}
\newcommand{\bsO}{{\bs{O}}}
\newcommand{\bsP}{{\bs{P}}}
\newcommand{\bsQ}{{\bs{Q}}} 
\newcommand{\bsR}{{\bs{R}}}
\newcommand{\bsS}{{\bs{S}}} 
\newcommand{\bsT}{{\bs{T}}}
\newcommand{\bsU}{{\bs{U}}}
\newcommand{\bsV}{{\bs{V}}}
\newcommand{\bsW}{{\bs{W}}}
\newcommand{\bsX}{{\bs{X}}}
\newcommand{\bsY}{{\bs{Y}}} 
\newcommand{\bsZ}{{\bs{Z}}}
%
\newcommand{\biX}{{\bs{X}}}

\newcommand{\modar}{\color{cyan}}
\newcommand{\modarn}{\color{black}}
\newcommand{\modch}{\color{black}}
\newcommand{\modchi}{\color{black}}
\newcommand{\revar}{\color{black}}
\newcommand{\rev}{\color{black}}

% Notation
\newcommand{\theconvex}{C}


\usepackage[textwidth=1.8cm, textsize=tiny]{todonotes}
\newcommand{\chiara}[1]{\todo[color=blue!20]{{ #1}}}
\newcommand{\arnaud}[1]{\todo[color=green!20]{{ #1}}}



\newcommand{\indi}[1]{\mathds{1}_{#1}}
\DeclarePairedDelimiter{\abs}{\lvert}{\rvert}


\title{Factorization by extremal privacy mechanisms: new insights into efficiency}
\author{Chiara Amorino\thanks{ Universitat Pompeu Fabra and Barcelona School of Economics, Department of Economics and Business, Ram\'on Trias Fargas 25-27, 08005, Barcelona, Spain. The author gratefully acknowledges financial support of PID2022-138268NB-I00/AEI/10.13039/501100011033.} \qquad Arnaud Gloter \thanks{Laboratoire de Math\'ematiques et Mod\'elisation d'Evry, CNRS, Univ Evry, Universit\'e Paris-Saclay, 91037, Evry, France.}} 

\begin{document}
\maketitle
%\footnote{
%CA gratefully acknowledge financial support of ERC Consolidator Grant 815703 “STAMFORD: Statistical Methods for High Dimensional Diffusions”.

%$^{*}$Unit\'e de Recherche en Math\'ematiques, Universit\'e du Luxembourg%, %\url{chiara.amorino@uni.lu}. \\

%{$^\dagger$ }

%{$^\ddagger$ }}



\begin{abstract}
\noindent We study the problem of efficiency under $\alpha$-local differential privacy ($\alpha$-LDP) in both discrete and continuous settings. Building on a factorization lemma—which shows that any privacy mechanism can be decomposed into an extremal mechanism followed by additional randomization—we reduce the Fisher information maximization problem to a search over extremal mechanisms. The representation of extremal mechanisms requires working in infinite-dimensional spaces and invokes advanced tools from convex and functional analysis, such as Choquet’s theorem. Our analysis establishes matching upper and lower bounds on the Fisher information in the high-privacy regime ($\alpha \to 0$), and proves that the maximization problem always admits a solution for any $\alpha$.
As a concrete application, we consider the problem of estimating the parameter of a uniform distribution on $[0, \theta]$ under $\alpha$-LDP. Guided by our theoretical findings, we design an extremal mechanism that yields a consistent and asymptotically efficient estimator in high privacy regime. Numerical experiments confirm our theoretical results.
\\
\\
\noindent
 \textit{Keywords: local differential privacy, Fisher information, efficiency, staircase mechanism, Choquet's theorem} \\ 
 
\noindent
\textit{MSC2020 subject classifications: Primary 62F12, 68P27; secondary 62B15, 46A55}
\end{abstract}

\tableofcontents


\section{Introduction}{\label{s: intro}
With the generalization of large-scale data collection and the ever-increasing computational power of modern computers, statisticians now have access to unprecedented amounts of data. This revolution is not merely quantitative but also qualitative, as it fundamentally changes the way knowledge is built. Statisticians constructs models based on data arising from the natural and social sciences (e.g., physics, chemistry, sociology). However, this data is often sensitive. For example, several recent studies have relied on sensitive medical data, %\cite{Dub20, Jun21, TruOud22}, 
where particular care must be taken to ensure privacy, leading to a necessary trade-off between statistical utility and privacy protection.

Classical privacy-preserving techniques, such as data permutation and traditional anonymization procedures, have proven to be largely ineffective. In fact, numerous privacy breaches have been documented using these methods (see, for example, \cite{Nar08}). One of the most famous examples is the case of the Netflix Prize competition, where anonymized movie ratings were successfully de-anonymized, forcing Netflix to cancel its second competition in 2010 \cite{KeaRot19}. For a detailed and engaging account of this story, we also recommend the introduction of the PhD thesis \cite{Lal23}, which provides an in-depth, well-written, and slightly romanticized (but very enjoyable) description of the events.

In response to such challenges, researchers in statistics, databases, and computer science have worked to formalize privacy as a mathematical framework for limiting disclosure risks. This formalization aims to protect individual data while still allowing for the statistical analysis of aggregated information. The foundational work in this direction is that of Dwork et al. \cite{Dwo06}, who introduced the notion of global differential privacy (also known as central differential privacy). In this setting, a trusted aggregator has access to the entire dataset and produces privatized outputs—while the raw data remains hidden from external users. A practical example of this model would be a centralized data center that stores all user information but only releases privatized statistics.

Research on central privacy has since flourished, but alternative models have also emerged. Notably, local differential privacy (LDP) does not assume a trusted aggregator. Instead, each user privatizes their own data before sending it to a central aggregator. At first glance, local privacy may appear more desirable since it offers stronger privacy guarantees. However, this increased protection often comes at a significant cost to statistical utility, leading to lower-quality estimates compared to those obtained under central privacy.

In recent years, several important extensions and refinements of differential privacy have been developed. These include concentrated differential privacy \cite{BunSte16, DwoRot16}, pufferfish privacy \cite{puff}, and componentwise local differential privacy \cite{AmoGlo23}, among others.



Major technology companies such as Apple \cite{5LDP, 49LDP} and Google \cite{1LDP, 23LDP} have already adopted local differential privacy mechanisms, reflecting the significant real-world impact of privacy-preserving measures on billions of devices. As a result, understanding the fundamental trade-off between local differential privacy and statistical utility has become crucial, leading to a surge of academic research focused on statistical inference under privacy constraints.

A key feature of local differential privacy, as introduced in Section \ref{s: pb formulation and main} and specifically formalized in Equation \eqref{eq: def LDP}, is that the privacy level is directly controlled by the parameter $\alpha \in [0, \infty)$. When $\alpha = 0$, perfect privacy is achieved, while as $\alpha \to \infty$, the privacy constraint progressively disappears. A particularly active area of current research focuses on fixing a desired privacy level $\alpha$ and, among all random mechanisms that satisfy this privacy constraint, identifying those that optimize statistical utility according to a specified criterion. This is precisely the approach adopted in the present manuscript, which specifically aims to investigate privacy mechanisms that are optimal in the sense of maximizing the Fisher information.



The foundational steps in the study of statistical inference under local differential privacy constraints can be found in early works such as \cite{Martin} and \cite{WassermanZhou10}. Since then, the field has experienced rapid and extensive development, with an increasing number of contributions spanning various areas of statistical methodology. Recent advancements include investigations into hypothesis testing \cite{4CLDP, 29CLDP}, M-estimation \cite{Avella}, robustness \cite{31CLDP}, change-point detection \cite{6CLDP}, as well as estimation of means and medians \cite{Martin}. Moreover, significant progress has been made in nonparametric estimation \cite{10CLDP, 11CLDP, Kro24} and parametric inference \cite{AmoGloHal, Asi22, DucRog19, Jos19, KalSte}. Particularly within the parametric setting, considerable attention has been devoted to the pursuit of efficient estimators under privacy constraints—a central theme in recent research.

In the centralized (global) differential privacy paradigm, Smith’s early work \cite{Smi08, Smi11} laid the foundations for efficient estimation under privacy constraints. By contrast, in the local model, \cite{BarChe20} derived upper bounds on the Fisher information achievable under LDP but did not construct mechanisms that attain them. Nam and Lee \cite{NamLee22} then investigated the maximization of Fisher information in the LDP setting, albeit restricted to one-bit communication schemes.

More recently, Duchi and Ruan \cite{DucRua24} introduced the notion of $L_1$-information to precisely characterize the local minimax risk of regular parametric models under LDP, up to universal constants. The closest precursor to our work is the elegant treatment by Steinberger \cite{steinbergerEfficiencyLocalDifferent}, which establishes both the existence of asymptotically efficient locally private estimators and explicit constructions achieving minimum variance. A detailed comparison between his work and ours will be presented in the next subsection, after rigorously describing our findings.

\subsection{Our contribution}
Following the line described above, we pose the question: among all mechanisms in \(\mathcal{Q}_\alpha:= \{ q : q \text{ satisfies } \alpha\text{-LDP} \}\), which one maximizes the Fisher information?
Equivalently, given a regular parametric model \(\mathcal{P}\), the perturbed model \(q\circ\mathcal{P}\) remains regular (see \cite{steinbergerEfficiencyLocalDifferent} and recap in Section ~\ref{ss:Parametric models and Fisher information}), with Fisher information \(\mathcal{I}_{\theta_0}(q\circ\mathcal{P})\). Thus our core problem becomes looking for solutions of the optimization problem
\begin{equation}\label{eq: max info intro}
\sup_{q\in\mathcal{Q}_\alpha}\; \mathcal{I}_{\theta_0}\bigl(q\circ\mathcal{P}\bigr).
\end{equation}
Our factorization lemma (Lemma~\ref{lem : factorisation cas fini}) offers a new vantage point: it shows that any \(\alpha\)-LDP mechanism can be decomposed into two mechanisms, one of which is extremal—that is, it saturates the privacy constraint (see Eq.~\eqref{eq : extremal constrainte discret}). We also refer to these extremal mechanisms as staircase mechanisms.

Staircase mechanisms have become a central object of study in the finite-alphabet setting. Let us mention some key references on the topic. Geng and Viswanath \cite{GenVis12} introduced the one-dimensional staircase mechanism, showing both theoretically and empirically that its piecewise-constant noise outperforms the Laplace mechanism. They later extended it to \(d\) dimensions \cite{Gen15}, again proving superiority over Gaussian or Laplace noise for \(d\ge2\). More recently, \cite{Kul23} proposed a staircase-like scheme for one-dimensional mean estimation, while Kairouz et al.\ \cite{kairouzExtremalMechanismsLocal2015} formulated the privacy–utility trade-off as a constrained optimization and proved that, on finite spaces, the family of staircase mechanisms contains the exact optima for a wide class of information-theoretic utilities.

Leveraging our factorization lemma, we reduce the Fisher‐information maximization in \eqref{eq: max info intro} to searching over extremal (or staircase) mechanisms.  This additional structure lets us show that the optimum admits the closed form
\[
\max_{q\in\mathcal{Q}_\alpha}\mathcal{I}_{\theta_0}(q\circ\mathcal{P})
=(e^\alpha -1)^2\,M^*,
\]
where \(M^*\) is the solution to the linear program
\begin{equation*}
M^*=\max_{w\in\mathbb{R}_+^{2d}}w^\top i
\quad\text{s.t.}\quad R\,w=\mathbf1.
\end{equation*}
Here \(i\) is explicit and given in \eqref{eq : def i beta} in the main and \(R\) is the \(2^d\times d\) staircase‐pattern matrix (see Definition~\ref{def: staircase matrix}). This constrained optimization problem parallels the privacy–utility trade‐off of Kairouz et al.\ \cite{kairouzExtremalMechanismsLocal2015} (see their Eq.~(8) and Theorem 4). 

By establishing matching upper and lower bounds (Theorem~\ref{th : info cas fini}), we not only characterize \(M^*\) but also explicitly construct a staircase mechanism that attains it. Then, to address non discrete spaces,
%infinite-dimensional domains, 
we introduce a continuous analogue of the staircase mechanism (Section~\ref{ss: infinite}). Beyond our specific setting, we believe the continuous staircase mechanism holds broader interest, as its potential applications extend well beyond our Fisher information maximization problem and it opens the way for optimal LDP estimation in infinite domains. 



 It is worth noting that  
 \cite{KalSte} analyzes the Fisher information in a continuous space \(\mathcal{X}\). However, his solution to the (continuous) optimization problem is obtained by approximating it through a discretization procedure, whereas our approach provides an exact continuous solution, without the need for any discrete approximation.

Extending our findings to the continuous case is a highly challenging task from a purely theoretical and mathematical perspective. It requires several advanced tools from functional and convex analysis, 
%as well as concepts from topology and set theory (such as Zorn's lemma), 
which make the proofs rather involved and interdisciplinary. %A deep understanding of several distinct mathematical topics is needed to rigorously carry out the analysis. 
The main tools of functional analysis used in the paper are recalled in Section \ref{S: tools}.  


The key idea is as follows. In the finite case, proving the factorization lemma relies on the fact that every point in the hyperrectangle \([1, e^\alpha]^d\), where \(d = |\mathcal{X}|\), can be expressed as a convex combination of its extreme points \(\{1, e^\alpha\}^d\). This property is guaranteed by Carathéodory’s theorem, which applies because the hyperrectangle is a compact convex subset of \(\mathbb{R}^d\).

When extending this result to the continuous case, the hyperrectangle is replaced by the set
\[
\theconvex := \left\{ v : \mathcal{X} \rightarrow \mathbb{R} \ \text{measurable} \ \middle| \ 1 \leq v(x) \leq e^\alpha \ \text{a.e.} \right\} \subset L^\infty(\mathcal{X}, \mathbb{R}).
\]
The extremal points of the finite-dimensional hyperrectangle translate, in the continuous setting, to the set
\[
\mathcal{E} := \left\{ v : \mathcal{X} \rightarrow \mathbb{R} \ \text{measurable} \ \middle| \ v(x) \in \{1, e^\alpha\} \ \text{a.e.} \right\}.
\]
In this framework, Carathéodory’s theorem is replaced by Choquet’s theorem. To apply it, we equip \(\theconvex\) with the distance \(d_\star\), which metrizes the weak-\(\star\) topology, and we prove that \(\theconvex\) is convex and compact with respect to this topology.

Another significant challenge arises from the fact that \(\theconvex \subset L^\infty(\mathcal{X}, \mathbb{R})\). For our results, it is crucial to evaluate elements of \(\theconvex\) at specific points \(x \in \mathcal{X}\), which leads to the introduction of the evaluation operator. The key difficulty is that this operator is not continuous with respect to the weak-\(\star\) topology. Providing a rigorous interpretation of pointwise evaluation of integrals defined over the domain \(\mathcal{E}\) is therefore particularly delicate, as detailed in Subsection~\ref{ss:point eval}. The solution to this problem is provided in Proposition~\ref{p: evaluation}, whose proof is quite technical and relies on the introduction of Bochner integrals and on proving the separability of \(\theconvex\) when endowed with the weak-\(\star\) topology.

These results allow us to extend the proof of the factorization lemma to the continuous setting. Consequently, we are able to establish, even in the continuous %infinite-dimensional 
case, both upper and lower bounds for the Fisher information and these bounds coincide as \(\alpha \to 0\), which corresponds to the regime with a significant level of privacy. More precisely, we prove in Theorem \ref{thm : equivalent alpha petit ctn} that the best Fisher information under $\alpha$-LDP constraint in a regular model $(p_\theta)_\theta$ is equivalent to $\frac{\alpha^2}{4} \left( \int \abs{\dot{p}_{\theta}(x)} dx \right)^2$ as $\alpha\to0$.




More generally, we show that, without imposing any restriction on \(\alpha\), the maximization problem introduced in \eqref{eq: max info intro} admits a solution. Specifically, there exists a Radon sub-probability measure \(\bar{\mu}\) on \(\mathcal{E}\) such that the corresponding (continuous version of the) staircase privacy mechanism \(q^{(\bar{\mu})}\) satisfies
\[
\mathcal{I}_{\theta}(q^{(\bar{\mu})} \circ \mathcal{P}) = \max_{q \in \mathcal{Q}_\alpha} \mathcal{I}_{\theta}(q \circ \mathcal{P}),
\]
see Theorem~\ref{thm: main opt cont} for details.

It is important to emphasize that the proof of this result is far from trivial. Although the compactness of \(\theconvex\) in the weak-\(\star\) topology guarantees the existence of a solution in \(\theconvex\), the fact that \(\mathcal{E}\) is not closed prevents us from directly concluding that the solution is supported on \(\mathcal{E}\). Instead, the proof relies on more advanced techniques, combining convex analysis and set-theory arguments, %(in particular the application of Zorn's lemma), 
to establish the existence of an extremal solution supported on \(\mathcal{E}\).


It is also important to highlight that our theoretical framework provides valuable intuition about the structure of optimal mechanisms, which enables us to effectively apply our results even in models that are not regular. This is precisely the focus of Section~\ref{s: unif}, where we consider the case of a private variable following a uniform distribution on the interval \([0, \theta]\), with unknown \(\theta\). In this scenario, we naturally operate in the continuous setting, since \(\mathcal{X} = [0, \infty)\). The factorization lemma continues to apply in this context, allowing us to establish an upper bound on the Fisher information that closely resembles the one obtained for regular models. What is quite intriguing is that, even if we start from a model that is not regular, once privacy is enforced we effectively recover a model that becomes regular again (see Proposition \ref{prop : Fisher unif}).


In Section~\ref{Ss : efficient est unif}, we further assume the availability of a preliminary estimator of the parameter, which does not need to be particularly accurate or close to the true value (note that two-stage estimators are commonly employed in this setting, as illustrated, for instance, in \cite{KalSte} and \cite{steinbergerEfficiencyLocalDifferent}). Guided by the intuition developed through our earlier theoretical results, we introduce an \(\alpha\)-LDP extremal mechanism corresponding to a measure supported on two points in \(\mathcal{E}\). We show that, when the preliminary estimate is smaller than the true parameter value \(\theta_0\), the resulting estimator is consistent, asymptotically Gaussian, and, as \(\alpha \to 0\), its variance matches the previously derived upper bound on the Fisher information. This establishes the efficiency of the proposed estimator in this privacy regime.

We conclude our analysis with numerical experiments to assess the practical performance of the estimation procedure. The simulations confirm our theoretical findings: when the preliminary estimate exceeds \(\theta_0\), the estimator is significantly biased, while when the preliminary estimate is below \(\theta_0\), the estimator performs well and aligns with the theoretical predictions.

\vspace{0.3 cm}

The structure of the paper is as follows. In Section~\ref{s: pb formulation and main}, we rigorously introduce the notion of \(\alpha\)-LDP and formulate the estimation problem. Our main theoretical results are split between Section~\ref{ss: finite}, which presents the findings for the finite case, and Section~\ref{ss: infinite}, which addresses the  continuous setting. The application to the uniform model is developed in Section~\ref{s: unif}. Section~\ref{S: tools} provides the functional analysis tools required for the continuous framework. In Section~\ref{s: proof main}, we present the proofs of the main results, while the proofs of more technical supporting results are deferred to Section~\ref{s: proof technical}.






\section{Problem formulation}{\label{s: pb formulation and main}}
We begin by describing the mathematical formulation of the notion of privacy that we will work with. The process of privatizing raw data and transforming it into public data is modeled through a conditional distribution, known as a privacy mechanism or channel distribution. Let us now formalize this setup.

Let $\mathcal{X}$ and $\mathcal{Z}$ be two separable, complete metric spaces. Equipped with their respective Borel $\sigma$-algebras, they define the measurable spaces $(\mathcal{X}, \Sigma_X)$ and $(\mathcal{Z}, \Sigma_Z)$. In particular, $(\mathcal{X}, \Sigma_X)$ represents the space of sensitive (raw) data, while $(\mathcal{Z}, \Sigma_Z)$ corresponds to the space of privatized (public) data.

Let us formally describe the mechanism that transforms raw data into its public counterpart. Consider $X_1, \dots, X_n$, i.i.d. raw data, distributed as $X \in \mathcal{X}$. We begin by introducing the sequentially interactive privacy mechanism, in which the privatization of the $i$-th observation depends not only on the corresponding sensitive datum $X_i$, but also on the previously released public outputs $Z_1, \dots, Z_{i-1}$. This results in the following conditional independence structure:
$$
X_i, Z_1, \dots, Z_{i-1} \rightarrow Z_i, \qquad Z_i \perp X_k \mid X_i, Z_1, \dots, Z_{i-1} \quad \text{for } k \neq i.
$$
In other words, the privatized output $Z_i$ is drawn according to
$$
Z_i \sim Q_i(\cdot \mid X_i = x_i, Z_1 = z_1, \dots, Z_{i-1} = z_{i-1}),
$$
for a collection of Markov kernels $Q_i : \Sigma_Z \times \mathcal{X} \times \mathcal{Z}^{i-1} \rightarrow [0,1]$.

A line of research concentrates on a special case of the general mechanism described above: the non-interactive setting. In this setting, the value of the public variable $Z_i$ depends solely on the corresponding raw value $X_i$, and is independent of the previously generated variables $Z_1, \dots, Z_{i-1}$. Consequently, the Markov kernel $Q$ does not depend on $i$, and the same mechanism is applied uniformly across all observations. Formally, for each $i \in \{1, \dots, n\}$, the privatized output is given by
\begin{equation}{\label{eq: drawn Z nonint}}
 Z_i \sim Q(\cdot \mid X_i = x_i).   
\end{equation}
On the one hand, non-interactive mechanisms yield i.i.d. privatized samples, which are often easier to analyze and work with. On the other hand, interactive mechanisms can incorporate more information about the data and, as a result, sometimes allow for improved statistical performance compared to non-interactive approaches (see, e.g., \cite{ButKlo25} and \cite{Kro24}).

Later, we will search for privacy mechanisms that are optimal in a sense that will be extensively detailed. %We will see that certain mechanisms, such as the staircase mechanism (which will be discussed in detail below), play a central role in our analysis. 
To study such notions of optimality, it is crucial to quantify the level of privacy provided by a mechanism, allowing for meaningful comparisons between different mechanisms that guarantee the same degree of privacy.

Privacy is measured using the notion of local differential privacy. Given a parameter $\alpha \ge 0$, we say that a random variable $Z_i$ is an $\alpha$-locally differentially private version of $X_i$ if, for all $z_1, \dots, z_{i-1} \in \mathcal{Z}$ and all $x, x' \in \mathcal{X}$, the following holds:
$$
\sup_{A \in \Sigma_Z} \frac{Q_i(A \mid X_i = x, Z_1 = z_1, \dots, Z_{i-1} = z_{i-1})}{Q_i(A \mid X_i = x', Z_1 = z_1, \dots, Z_{i-1} = z_{i-1})} \le \exp(\alpha).
$$
We denote by $\mathcal{Q}_\alpha$ the set of all Markov kernels satisfying this local $\alpha$-differential privacy constraint.

As mentioned earlier, the parameter $\alpha$ quantifies the strength of privacy: the smaller the value, the harder it is to infer sensitive information from the released data. In particular, $\alpha = 0$ corresponds to perfect privacy, while letting $\alpha \to \infty$ gradually removes the privacy constraint. An analogous condition holds in the non-interactive case:
$$
\sup_{A \in \Sigma_Z} \frac{Q(A \mid X = x)}{Q(A \mid X = x')} \le \exp(\alpha).
$$
Under the local differential privacy condition, the kernels $Q(\cdot \mid X = x)$ are mutually absolutely continuous for all $x \in \mathcal{X}$. This allows us to assume the existence of a dominating measure $\mu$ on $(\mathcal{Z}, \Sigma_Z)$ with respect to which each kernel $Q$ admits a density, denoted $q_x(z)$. The $\alpha$-local differential privacy constraint can then be rewritten in terms of these densities as:
\begin{equation}{\label{eq: def LDP}}
\sup_{z \in \mathcal{Z}} \frac{q_x(z)}{q_{x'}(z)} \le \exp(\alpha), \quad \forall x, x' \in \mathcal{X}.
\end{equation}
From now on, we will write $q_x(z)$ to denote the density $q(z \mid X = x)$.

Finally, note that if $q_x(z) = 0$ for some $x \in \mathcal{X}$ and $z \in \mathcal{Z}$, then the $\alpha$-LDP constraint implies that $q_{x'}(z) = 0$ for all $x' \in \mathcal{X}$. In other words, such a value $z$ could be removed from $\mathcal{Z}$ without altering the randomization mechanism.

With this background in place, we can now present our main results, which crucially depend on whether the spaces under consideration are discrete or continuous. We begin with the finite setting in Section \ref{ss: finite}, as it allows us to fix the notation and highlight the key ideas in a simpler framework. We then extend our analysis in Section \ref{ss: infinite}, focusing on the scenario where the space $\mathcal{X}$ is an open subset of $\mathbb{R}^d$. %continuous. %infinite dimensional.


\section{Main results: Discrete case}{\label{ss: finite}}
Let us begin by considering the case where both $\mathcal{X}$ and $\mathcal{Z}$ are finite sets. Specifically, we take $\mathcal{X} = \{x_1, \dots, x_d\}$ and $\mathcal{Z} = \{z_1, \dots, z_l\}$.

In this setting, one of the most studied and widely used privacy mechanisms is the staircase mechanism. A key reference on this topic is \cite{kairouzExtremalMechanismsLocal2015}, where it is shown that, when $\mathcal{X}$ is finite, the optimal mechanism that maximizes utility functions expressible as sums of sublinear functions has an extremal structure.

Following \cite{kairouzExtremalMechanismsLocal2015}, we say that a randomization mechanism is extremal if, for all $z \in \mathcal{Z}$ and all $(x, x') \in \mathcal{X}^2$, the log-likelihood ratios satisfy
$$
\left| \ln \frac{q_{x'}(z)}{q_x(z)} \right| \in \{0, \alpha\},
$$
which is equivalent to the condition:
\begin{equation}\label{eq : extremal constrainte discret}
\forall z \in \mathcal{Z}, \ \forall (x,x') \in \mathcal{X}^2, \quad \frac{q_{x'}(z)}{q_x(z)} \in \{e^{-\alpha}, 1, e^\alpha\}.
\end{equation}
As before, if $q_x(z) = 0$ for some $z \in \mathcal{Z}$ and $x \in \mathcal{X}$, this constraint is interpreted as requiring $q_{x'}(z) = 0$ for all $x' \in \mathcal{X}$.

We refer to any mechanism satisfying this constraint as a staircase mechanism. In the context of global differential privacy, a similar definition can be extended to apply to all neighboring database queries $x, x'$ (or, equivalently, pairs differing within a given sensitivity), thus recovering many known optimal mechanisms. Notably, the mechanisms studied in \cite{GenVis12, GenVis13} and \cite{Ghosh12} are particular instances of the staircase mechanism defined in \eqref{eq : extremal constrainte discret}.

There are several reasons why staircase mechanisms deserve particular attention. First, as noted above, they arise naturally in the context of optimal mechanisms and have been extensively studied across different notions of privacy and utility. Second, extremal mechanisms saturate the $\alpha$-LDP constraint. Third, and most importantly for our purposes, one of our main results (see Lemma \ref{lem : factorisation cas fini}) shows that any privacy mechanism can be decomposed into the composition of two mechanisms, one of which is an extremal mechanism.

Given the central role staircase mechanisms play in our analysis, and the frequency with which we will refer to them throughout the paper, we now provide a more detailed description.

\begin{df}\label{def: staircase matrix}
For $\beta \in \{0, \dots, 2^d - 1\}$, let $d_j(\beta)$ denote the $j$-th component of the $d$-dimensional binary vector corresponding to the binary representation of $\beta$ in the dyadic basis, i.e., $\beta = \sum_{j=0}^{d-1} d_j(\beta) 2^j$, with $d_j(\beta) \in \{0, 1\}$.
A matrix is called a staircase pattern matrix if its $\beta$-th column is a vector $r_\beta \in \{1, e^\alpha\}^d$ for each $\beta \in \{0, \dots, 2^d - 1\}$, where the $j$-th component of $r_\beta$ is defined as
$(r_\beta)_j = \indi{\{d_j(\beta) = 0\}} + e^\alpha \, \indi{\{d_j(\beta) = 1\}}.$
Each column vector $r_\beta$ is called a staircase pattern.
\end{df}

To clarify the construction, consider the case $d = 2$. There are $2^d = 4$ staircase patterns, and the corresponding staircase pattern matrix is:
$$
\begin{bmatrix}
1 & 1 & e^\alpha & e^\alpha \\
1 & e^\alpha & 1 & e^\alpha
\end{bmatrix}.
$$
In general, for arbitrary $d$, the staircase patterns are given by:
\begin{equation}\label{eq : R matrix by slice}
	r_0=\begin{bmatrix} 1\\1\\ 1\\ \vdots\\ 1\end{bmatrix},~
	r_1=\begin{bmatrix} e^\alpha\\1 \\ 1\\ \vdots\\ 1\end{bmatrix},~
	r_2=\begin{bmatrix} 1 \\ e^\alpha\\ 1\\ \vdots\\ 1\end{bmatrix},~\dots~,~
	r_{2^d-1}=\begin{bmatrix} e^\alpha \\ e^\alpha\\ e^\alpha\\ \vdots\\ e^\alpha\end{bmatrix}.
\end{equation}
Equivalently, the vectors $(r_\beta)_{\beta \in \{0, \dots, 2^d - 1\}}$ are the extremal points of the hyper-rectangle $[1, e^\alpha]^d$.
We denote by $\mathcal{E} = \{0, \dots, 2^d - 1\}$ the index set of these extremal points ${r_0, \dots, r_{2^d - 1}}$.
Interpreting each $r_\beta$ as a function on $\mathcal{X}$, we write $r_\beta(x_j) = (r_\beta)_j$ for $x_j \in \mathcal{X}= \{ x_1, ... , x_d \}$. 

For each $\beta \in \{0, \dots, 2^d - 1\}$, we define the level sets:
\begin{equation}\label{eq: Fbeta+ 4.8}
F^+_\beta = \{ x \in \mathcal{X} \mid r_\beta(x) = e^\alpha \}, \quad
F^-_\beta = \{ x \in \mathcal{X} \mid r_\beta(x) = 1 \}.
\end{equation}
This implies that for all $x \in \mathcal{X}$, the function $r_\beta$ satisfies
\begin{equation}{\label{eq: r beta x}}
r_\beta(x) = e^\alpha \indi{F^+_\beta}(x) + \indi{F^-_\beta}(x) = 1 + (e^\alpha - 1) \indi{F^+_\beta}(x),
\end{equation}
i.e., the function $r_\beta$ takes values in $\{1, e^\alpha\}$, with the structure fully determined by $F^+_\beta$.

Note that as $\beta$ ranges from $0$ to $2^d - 1$, the sets $F^+_\beta$ range over all possible subsets of $\mathcal{X}$. For instance, $F^+_0 = \emptyset$ because $r_0 = (1, \dots, 1)^T$ by \eqref{eq : R matrix by slice}, and $F^+_{2^d - 1} = \mathcal{X}$ since $r_{2^d - 1} = (e^\alpha, \dots, e^\alpha)^T$.


\subsection{A factorization lemma}

We are now ready to present our first main result: the factorization lemma, which shows that any privacy mechanism can be decomposed into the composition of two mechanisms, one of which is extremal in the sense of Definition \ref{def: staircase matrix}.

In order to do so, let us introduce some notation. If $q^{(1)}: \mathcal{X} \to \mathcal{E}$ and $q^{(2)}: \mathcal{E} \to \mathcal{Z}$ are privacy mechanisms between finite sets, we denote their composition by $q = q^{(2)} \circ q^{(1)}$, representing the sequential application of the two randomization steps. The resulting mechanism satisfies:
$$
q_x(z) = \sum_{y \in \mathcal{E}} q^{(1)}_x(y) \, q^{(2)}_y(z).
$$
\begin{lemma}\label{lem : factorisation cas fini}
	Let $q$ be an $\alpha$-LDP mechanism between the finite spaces $\cX=\{x_1,\dots,x_d\}$ and $\cZ=\{z_1,\dots,z_l\}$. %There exists $\cY$ a finite space, 
	There exist $q^{(1)}$ an extremal mechanism from $\cX$ to $\cE$, and $q^{(2)}$ a randomization mechanism for $\cE$ to $\cZ$ such that $q=q^{(2)}\circ q^{(1)}$.
	
	Moreover, we have the representation $q^{(1)}_x(\beta)= \omega_\beta r_\beta(x)$, where
	$(\omega_\beta)_\beta$ defines a sub-probability measure on $\cE$, satisfying $e^{-\alpha} \le \sum_{\beta=0}^{2^d-1} \omega_{\beta}\le 1$.
\end{lemma}
\begin{proof}
We define the minimal mass function
$$
\underline{q}(z) = \min_{x \in \cX} q_x(z).
$$
Note that, from the definition of $\alpha$-LDP in \eqref{eq: def LDP} and the following remark, it follows that if $\underline{q}(z) = 0$, then $q_x(z) = 0$ for all $x \in \cX$.

Let us now introduce a normalized version of the mechanism $q$. For each $(x, z) \in \cX \times \cZ$, define
$$
u_x(z) = 
\begin{cases}
\displaystyle \frac{q_x(z)}{\underline{q}(z)}, & \text{if } \underline{q}(z) > 0, \\
1, & \text{if } \underline{q}(z) = 0.
\end{cases}
$$
By construction and by the $\alpha$-LDP constraint in \eqref{eq: def LDP}, we have that for all $z \in \cZ$, the vector $u(z) := (u_x(z))_{x \in \cX} \in [1, e^\alpha]^d$, where we recall that $d = |\cX|$.

The hyperrectangle $[1, e^\alpha]^d \subset \mathbb{R}^d$ is compact and convex, so by the Carathéodory theorem, the vector $u(z)$ belongs to the convex hull of the extremal points of $[1, e^\alpha]^d$. Therefore, for all $z \in \cZ$, we may write
$$
u(z) = \sum_{\beta=0}^{2^d - 1} c_{z, \beta} r_\beta,
$$
where $c_{z, \beta} \ge 0$, $\sum_{\beta=0}^{2^d - 1} c_{z, \beta} = 1$, and the $r_\beta$'s are the staircase pattern vectors defined in \eqref{eq : R matrix by slice}. Accordingly, for all $x \in \cX$,
$$
u_x(z) = \sum_{\beta=0}^{2^d - 1} c_{z, \beta} r_\beta(x),
$$
and thus
\begin{equation} \label{eq : qxz convex hull}
q_x(z) = \underline{q}(z) u_x(z) = \sum_{\beta=0}^{2^d - 1} \underline{q}(z) c_{z, \beta} r_\beta(x).
\end{equation}
Define, for each $\beta \in \mathcal{E} = \{0, \dots, 2^d - 1\}$,
\begin{equation} \label{eq: def omega beta 4.5}
\omega_\beta := \sum_{z \in \cZ} \underline{q}(z) c_{z, \beta} \ge 0.
\end{equation}
Since $q_x(z)$ is a probability distribution over $\cZ$, we have $\sum_{z \in \cZ} q_x(z) = 1$, which yields
\begin{equation} \label{eq: constraint factorisation 4.55}
1 = \sum_{\beta=0}^{2^d - 1} \sum_{z \in \cZ} \underline{q}(z) c_{z, \beta} r_\beta(x) = \sum_{\beta=0}^{2^d - 1} \omega_\beta r_\beta(x).
\end{equation}
This suggests defining the first-stage mechanism $q^{(1)} : \cX \to \cE$ as
\begin{equation} \label{eq: cond q1 4.6}
q^{(1)}_x(\beta) := \omega_\beta r_\beta(x),
\end{equation}
which forms a probability distribution over $\cE = \{0, \dots, 2^d - 1\}$. Moreover, it satisfies the extremality constraint in \eqref{eq : extremal constrainte discret}, since
$$
\frac{q^{(1)}_x(\beta)}{q^{(1)}_{x'}(\beta)} = \frac{r_\beta(x)}{r_\beta(x')} \in \{e^{-\alpha}, 1, e^\alpha\},
$$
using the fact that each $r_\beta$ takes values in $\{1, e^\alpha\}$.
Now fix any $x \in \cX$. Since $r_\beta(x) \in \{1, e^\alpha\}$, we obtain the inequalities
\begin{equation}{\label{eq: mass omega beta}}
 1 = \sum_{\beta} \omega_\beta r_\beta(x) \le \sum_\beta \omega_\beta e^\alpha, \quad \text{and} \quad 1 = \sum_\beta \omega_\beta r_\beta(x) \ge \sum_\beta \omega_\beta,   
\end{equation}
implying that the total mass $\sum_\beta \omega_\beta$ lies in $[e^{-\alpha}, 1]$. Thus, $\omega$ is a sub-probability measure.

We now define the second-stage mechanism $q^{(2)} : \cE \to \cZ$. For all $\beta$ such that $\omega_\beta > 0$, set
$$
q^{(2)}_\beta(z) := \frac{\underline{q}(z) c_{z, \beta}}{\omega_\beta},
$$
which defines a probability distribution on $\cZ$, due to the normalization in \eqref{eq: def omega beta 4.5}. If instead $\omega_\beta = 0$, we may define $q^{(2)}_\beta(z)$ arbitrarily as any probability distribution on $\cZ$ — for example, set
$$
q^{(2)}_\beta(z) := \frac{1}{|\cZ|} \quad \text{for all } z \in \cZ.
$$
This definition is consistent, since for such $\beta$, $\underline{q}(z) c_{z, \beta} = 0$ for all $z$.

In all cases, we have the identity
\begin{equation} \label{eq: 4.75}
\omega_\beta q^{(2)}_\beta(z) = \underline{q}(z) c_{z, \beta}, \quad \text{for all } (\beta, z) \in \cE \times \cZ.
\end{equation}
Substituting \eqref{eq : qxz convex hull} and \eqref{eq: 4.75} into the expression for $q_x(z)$, and using \eqref{eq: cond q1 4.6}, we obtain the desired factorization:
$$
q_x(z) = \sum_{\beta=0}^{2^d - 1} \omega_\beta q^{(2)}_\beta(z) r_\beta(x) = \sum_{\beta=0}^{2^d - 1} q^{(1)}_x(\beta) q^{(2)}_\beta(z),
$$
which shows that $q = q^{(2)} \circ q^{(1)}$.
\end{proof}

\begin{rem}
The factorization lemma above shows that any $\alpha$-LDP mechanism can be decomposed into an extremal $\alpha$-LDP mechanism, followed by an additional randomization step. Since any such randomization further reduces the statistical information available from the data—without contributing to the satisfaction of the $\alpha$-LDP constraint— applying a second randomization is, from a statistical perspective, detrimental. The first (extremal) step alone suffices to ensure privacy, so only extremal mechanisms need to be considered in theory.

Building on this insight, we will derive an expression for the maximal Fisher information achievable by any $\alpha$-LDP mechanism, under suitable regularity assumptions on a parametric model.
\end{rem}


\begin{rem}
The factorization lemma \ref{lem : factorisation cas fini} applies to the case where the input space $\mathcal{X}$ is finite. In this setting, the extremal mechanism takes values in the set $\mathcal{E}$, whose cardinality is $2^{\Card(\mathcal{X})}$—significantly larger than $\Card(\mathcal{X})$ itself. However, the actual support of the randomization $q^{(1)}$ is determined by the support of the sub-probability measure $\omega$, which can be much smaller than the full set $\mathcal{E}$ (see Point (a) of Theorem 2 in \cite{kairouzExtremalMechanismsLocal2015}).

One of the main findings of this paper consists in extending Lemma \ref{lem : factorisation cas fini} to the case where the input space $\mathcal{X}$ is continuous, for example $\mathcal{X} = \mathbb{R}$. As already mentioned in the introduction and as we will see in Section \ref{ss: infinite}, in such a setting the set of extremal points $\mathcal{E}$ becomes infinite-dimensional.
\end{rem}



\subsection{Parametric models and Fisher information}\label{ss:Parametric models and Fisher information}
In this subsection, we explain how the factorization lemma can be used to derive an expression for the maximal Fisher information achievable by any privacy mechanism satisfying the $\alpha$-LDP constraint, in the case where the data space $\mathcal{X}$ is discrete. To do so, we first introduce some background.

Consider i.i.d. data $X_1, \ldots, X_n$ drawn from a distribution $P_\theta$ on the finite space $\mathcal{X} = \{x_1, \ldots, x_d\}$, where $P_\theta$ belongs to a parametric model $\mathcal{P} = (P_\theta)_{\theta \in \Theta}$ defined on an open set $\Theta \subset \mathbb{R}^{d_\Theta}$. We assume that the model is regular in a sense to be better specified later (see in particular Definition \ref{def: DQM}), and denote the corresponding Fisher information by $\mathcal{I}_\theta(\mathcal{P}) \in \mathbb{R}^{d \times d}$. We write $p_\theta(x) = P_\theta(X = x)$ for the probability mass function.

If data must be protected through local differential privacy, the observed variables $Z_1, \ldots, Z_n$ are sanitized via a Markov kernel, as described in Section \ref{s: pb formulation and main}. Let $q : \mathcal{X} \to \mathcal{Z}$ be such an $\alpha$-LDP mechanism with output space $\mathcal{Z} = \{z_1, \ldots, z_l\}$. The distribution of the public variable $Z$ is then given by $\widetilde{p}_\theta(z) = P_\theta(Z = z) = \int_{\mathcal{X}} p_\theta(x) q_x(z) \mu(dx)$, where $\mu$ (resp. $\nu$) denotes the counting measure on $\mathcal{X}$ (resp. $\mathcal{Z}$).

If the map $\theta \mapsto p_\theta(x)$ is continuously differentiable for all $x \in \mathcal{X}$, then $\theta \mapsto \widetilde{p}_\theta(z)$ is also continuously differentiable for all $z \in \mathcal{Z}$. We will recall a more precise result on the preservation of regularity under privacy constraints, as established in \cite{steinbergerEfficiencyLocalDifferent} (see \eqref{eq : score public} below). Thus, the resulting statistical model under the privacy mechanism $q$ remains regular and has Fisher information denoted by $\mathcal{I}_\theta(q \circ \mathcal{P})$.

Under standard regularity assumptions, the maximum likelihood estimator in the privatized model $q \circ \mathcal{P}$ is asymptotically normal with distribution $N(0, \mathcal{I}_\theta(q \circ \mathcal{P})^{-1})$. It is also efficient in the sense that no regular estimator in $q \circ \mathcal{P}$ can have smaller asymptotic variance. However, since the mechanism $q$ is freely chosen subject to the $\alpha$-LDP constraint \eqref{eq: def LDP}, it is natural to ask which mechanism maximizes the Fisher information. This leads to the optimization problem
\begin{equation}{\label{eq: opt Fisher}}
\sup_{q \in \mathcal{Q}_\alpha} \mathcal{I}_\theta(q \circ \mathcal{P}),
\end{equation}
where we recall that $\mathcal{Q}_\alpha$ denotes the set of $\alpha$-LDP channels from $\mathcal{X}$ to a finite set $\mathcal{Z}$.
This optimization problem is challenging and has been thoroughly investigated in the elegant work \cite{steinbergerEfficiencyLocalDifferent}, which develops a general theory of statistical efficiency under local differential privacy.

The factorization lemma proven earlier allows us to revisit the problem from a novel perspective. It implies that we can restrict attention to the class of extremal (staircase) mechanisms, significantly reducing the search space. This additional structure enables us to derive an explicit expression for the maximal Fisher information under $\alpha$-LDP constraints, as presented in Theorem \ref{th : info cas fini}.

To formalize the notion of regularity and rigorously study asymptotic efficiency, we adopt the classical framework of differentiability in quadratic mean (DQM) at an interior point $\theta_0 \in \overset{\circ}{\Theta}$.

\begin{df}\label{def: DQM}
A statistical model $\mathcal{P} = (P_\theta)_{\theta \in \Theta}$, with $\Theta \subset \R^{d_\Theta}$, sample space $(\mathcal{X}, \mathcal{F})$, and dominating measure $\mu$, is said to be \textit{differentiable in quadratic mean (DQM) at $\theta_0 \in \overset{\circ}{\Theta}$} if the $\mu$-densities $p_\theta = \frac{dP_\theta}{d\mu}$ satisfy, for $h \in \mathbb{R}^{d_\Theta}$,
$$
\int_\cX \left(\sqrt{p_{\theta_0+h}(x)} - \sqrt{p_{\theta_0}(x)} - \frac{1}{2} h^{\top}s_{\theta_0}(x) \sqrt{p_{\theta_0}(x)} \right)^2 \mu(dx) = o(\|h\|^2)
\quad \text{as } h \to 0,
$$
for some measurable vector-valued function $s_{\theta_0} : \cX \to \mathbb{R}^{d_\Theta}$. This function is called the score at $\theta_0$.
\end{df}

This definition is independent of the choice of the dominating measure $\mu$. To see this, consider two dominating measures $\mu_1$ and $\mu_2$, and note that both corresponding integrals can be written with respect to the measure $\mu = \mu_1 + \mu_2$.

We say that a model is DQM (or simply regular) if it is DQM at every point $\theta \in \overset{\circ}{\Theta}$. Let us also recall some properties of the score function, which will be crucial in the sequel. According to Theorem 7.2 in \cite{Van07}, if the model $\mathcal{P}$ is DQM at $\theta_0$, then the score satisfies 
\begin{equation}{\label{eq: score centered}}
\mathbb{E}_{\theta_0}[s_{\theta_0}] = 0,
\end{equation}
and the Fisher information
$$
\mathcal{I}_{\theta_0}(\mathcal{P}) = \mathbb{E}_{\theta_0}[s_{\theta_0} s_{\theta_0}^{\top}]
$$
exists and is finite. Moreover, in a $\mathcal{C}^1$ model, the score function is explicitly given by
$$
s_{\theta_0}(x) = \frac{\nabla_\theta p_{\theta_0}(x)}{p_{\theta_0}(x)}.
$$
Regarding transformations of DQM models, Proposition A.5 in \cite{Bic93} states that if $\mathcal{P} = (P_\theta)_{\theta \in \Theta}$ is DQM and $T: \mathcal{X} \to \mathcal{Z}$ is measurable, then the transformed model $T\mathcal{P} = (P_\theta \circ T^{-1})_{\theta \in \Theta}$ is also DQM, with score function
$$
t_{\theta_0}(z) = \mathbb{E}[s_{\theta_0}(X) \mid Z = z],
$$
where $(X, Z) \sim (X, T(X))$ under $P_{\theta_0}$.

For local differential privacy, one must consider randomized transformations—i.e., Markov kernels. An extension of the result in \cite{Bic93} to this context is provided in Lemma 3.1 of \cite{steinbergerEfficiencyLocalDifferent}. It shows that if the original model $(P_\theta)_\theta$ is DQM at $\theta_0$, then the privatized model $(q \circ P_\theta)_\theta$ is also DQM at $\theta_0$, with score function
\begin{equation} \label{eq : score public}
t_{\theta_0}(z) = \mathbb{E}\left[s_{\theta_0}(X) \mid Z = z\right] =
\frac{\int_\cX s_{\theta_0}(x) q_x(z) p_{\theta_0}(x) \mu(dx)}{\int_\cX q_x(z) p_{\theta_0}(x) \mu(dx)}.
\end{equation}

In what follows, given a parametric model $\mathcal{P} = (P_\theta)_\theta$ on $\mathcal{X}$, we denote by $q \circ \mathcal{P} = (\tilde{p}_\theta)_\theta$ the model induced on $\mathcal{Z}$ via the channel $q$. We denote by
$$
\mathcal{I}_{\theta_0}(\mathcal{P}) = \mathbb{E}\left[s_{\theta_0}(X)s_{\theta_0}(X)^{\top}\right]
$$
the Fisher information of the private variable $X$, and by
$$
\mathcal{I}_{\theta_0}(q \circ \mathcal{P}) = \mathbb{E}\left[t_{\theta_0}(Z)t_{\theta_0}(Z)^{\top}\right]
$$
that of the public variable $Z$.

From Lemma 3.1 in \cite{steinbergerEfficiencyLocalDifferent}, we know that the Fisher information of the public variable is always smaller than or equal to that of the private one. More precisely,
\begin{equation} \label{eq : decroissance Fisher}
\mathcal{I}_{\theta_0}(q \circ \mathcal{P}) = \mathcal{I}_{\theta_0}(\mathcal{P}) -
\mathbb{E}\left[\left(s_{\theta_0}(X) - t_{\theta_0}(Z)\right)\left(s_{\theta_0}(X) - t_{\theta_0}(Z)\right)^{\top} \right].
\end{equation}

The factorization lemma provides a valuable structural insight into any randomized channel $q$. Specifically, it guarantees the existence of a decomposition of the form $q = q^{(2)} \circ q^{(1)}$, where $q^{(1)}$ is an extremal $\alpha$-LDP mechanism. This yields the chain of inequalities:
\begin{equation}{\label{eq: fact on Fisher}}
\mathcal{I}_{\theta_0}\big((\widetilde{p}_\theta)_\theta\big)
= \mathcal{I}_{\theta_0}(q \circ \mathcal{P}) 
= \mathcal{I}_{\theta_0}((q^{(2)} \circ q^{(1)}) \circ \mathcal{P}) 
= \mathcal{I}_{\theta_0}(q^{(2)} \circ (q^{(1)} \circ \mathcal{P})) 
\le \mathcal{I}_{\theta_0}(q^{(1)} \circ \mathcal{P}). 
\end{equation}
Therefore, to study the optimization problem in \eqref{eq: opt Fisher}, it suffices to restrict our attention to extremal mechanisms of the form $q = q^{(1)}$.

Moreover, the fact that \( q^{(1)} \) is an extremal staircase mechanism introduces additional structure that enables us to reformulate the optimization problem \eqref{eq: opt Fisher}. In the proof of Proposition \ref{prop : borne sup Fisher fini}, we show that if the parameter $\theta$ is one dimensional, the optimization problem \eqref{eq: opt Fisher} is the same as
\[
\max_{q \in \mathcal{Q}_\alpha} \, \mathcal{I}_{\theta_0}(q \circ \mathcal{P}) = (e^\alpha - 1)^2 M^\star,
\]
where \( M^\star \) denotes the optimal value of the following linear program:
\begin{equation} \label{eq: opt 14.75}
M^\star = \max_{\omega \in \mathbb{R}_+^{2d}} \, w^T i \quad \text{subject to} \quad Rw = \bm{1},
\end{equation}
with \( i = (i_\beta)_{\beta \in \mathcal{E}} \) defined by
\begin{equation} \label{eq : def i beta}
i_\beta := \frac{\left( \int_\mathcal{X} s_\theta(x) \indi{F^+_\beta}(x) p_\theta(x) \, \mu(dx) \right)^2}{1 + (e^\alpha - 1) \int_{F^+_\beta} p_{\theta_0}(x) \, dx},
\end{equation}
with $F_\beta^+$ as in \eqref{eq: Fbeta+ 4.8}. The expression of $i_\beta$ has been derived explicitly in our proof, starting from the use of staircase mechanisms. Here, \( R \) is the staircase pattern matrix introduced in Definition~\ref{def: staircase matrix}.


Readers familiar with \cite{kairouzExtremalMechanismsLocal2015} will notice a strong analogy between this optimization problem and the fundamental privacy–utility tradeoff discussed in their Equation (8), reformulated in their Theorem 4.

This leads directly to the statement of the next theorem. 

\begin{thm}\label{th : info cas fini}
	Assume that $\Theta \subset \R$ and that $\cP=(p_{\theta})_\theta$ 
	is DQM at $\theta_0 \in \overset{\circ}{\Theta}$ with score function $s_{\theta_0}$.
	 Assume that $s_{\theta_0}(x)\neq 0$ for all $x\in \cX$.
	Then, there exists $\overline{\alpha}$ such that if $\alpha \le \overline{\alpha}$,
	\begin{equation*}
		\max_{q \in \cQ_\alpha}  \cI_{\theta_0}(q \circ \cP) = 	\cI^{\text{max},\alpha}_{\theta_0}((p_\theta)_\theta)		
	\end{equation*}
		where
	\begin{equation}\label{eq : def I max caf fini}
		\cI^{\text{max},\alpha}_{\theta_0}((p_\theta)_\theta)
		:=\frac{(e^\alpha-1)^2}{4}% 
		\times \frac{\E\left[\abs{s_{\theta_0}(X)}\right]^2}{\left[ (1-n_\text{max})+e^\alpha n_\text{max}  \right]
			\left[n_\text{max}+(1-n_\text{max})e^\alpha   \right]}
	\end{equation}
	with 
    \begin{equation}{\label{eq: def nmax}}
   n_\text{max}:= \int_{x: s_{\theta_0}(x)>0} p_{\theta_0}(x) \mu(dx)
	 \in (0,1).  
    \end{equation}
     The value of $\overline{\alpha}$ is independent of the channel $q$, but depends on the model $(p_{\theta})_\theta$.
\end{thm}

The proof of Theorem \ref{th : info cas fini} is provided in Section \ref{s: proof main}. It follows directly from Propositions \ref{prop : borne sup Fisher fini} and \ref{prop : borne inf Fisher fini} in Section \ref{s: proof main}, which establish upper and lower bounds, respectively. We emphasize that in the lower bound, we explicitly construct an $\alpha$-LDP mechanism that attains the Fisher information $\cI^{\text{max},\alpha}_{\theta_0}((p_\theta)_\theta)$. Hence, this proposition gives an exact characterization of the maximal Fisher information achievable under the $\alpha$-LDP constraint, at least for small values of $\alpha$.

\begin{rem}{\label{rk: choice Fmax}}
The key idea in the proof of the upper bound is to introduce the sets
\[
F_{\max}^+ := \left\{ x \in \mathcal{X} : s_{\theta_0}(x) > 0 \right\} \quad \text{and} \quad F_{\max}^{\prime +} := \left\{ x \in \mathcal{X} : s_{\theta_0}(x) < 0 \right\},
\]
which form a partition of \( \mathcal{X} \), since \( s_{\theta_0}(x) \neq 0 \) by hypothesis. We then observe that for any measurable set \( F_\beta^+ \subset \mathcal{X} \), the numerator of \( i_\beta \) in \eqref{eq : def i beta} satisfies
\[
\left( \int_\mathcal{X} s_\theta(x) \indi{F^+_\beta}(x) p_\theta(x) \, \mu(dx) \right)^2 \le \left( \frac{1}{2} \int_\mathcal{X} |s_\theta(x)| p_\theta(x) \, \mu(dx) \right)^2 = \left( \frac{1}{2} \mathbb{E}[|s_\theta(X)|] \right)^2,
\]
with equality if and only if \( F_\beta^+ = F_{\max}^+ \) or \( F_\beta^+ = F_{\max}^{\prime +} \). Among the collection of sets \( (F_\beta^+)_{\beta \in \mathcal{E}} \) introduced in \eqref{eq: Fbeta+ 4.8}, we denote by \( \beta_{\max} \) and \( \beta_{\max}' \) the indices corresponding to \( F_{\max}^+ \) and \( F_{\max}^{\prime +} \), respectively. These provide natural candidates for the support of an optimal solution to the linear program \eqref{eq: opt 14.75}. The remainder of the proof of the upper bound is more technical, consisting in showing that the maximum of the optimization problem is indeed attained by a sub-probability measure \( (\omega_\beta)_\beta \) supported on the set \( \{ \beta_{\max}, \beta_{\max}' \} \).
\end{rem}

\begin{rem}
It is worth noting that, in the derivation of the upper bound, a preliminary and simple estimate shows that
$
	\cI_{\theta_0}((\widetilde{p}_\theta)_\theta) \le 
	\frac{(e^\alpha-1)^2}{4} \times \E\left[\abs{s_{\theta_0}(X)}\right]^2,
	$
and this bound holds even without assuming the additional condition $\alpha \le \bar{\alpha}$. Furthermore, when $\alpha$ is close to zero, this expression coincides with the one in \eqref{eq : def I max caf fini}, since the denominator
$
	\left[ (1-n_\text{max})+e^\alpha n_\text{max}  \right]
	\left[n_\text{max}+(1-n_\text{max})e^\alpha   \right]
	$
simply reduces to 1 in that regime. We will see in the next section that this result extends to more general settings, in particular when the sample space $\mathcal{X}$ is continuous (see Theorem \ref{thm : equivalent alpha petit ctn}).
\end{rem}



























\section{Main results: Extension to the Continuous Case}{\label{ss: infinite}}

The goal of this section is to extend the results established in Section~\ref{ss: finite} to the setting where \( \mathcal{X} \) is no longer a finite set. Specifically, we now consider \( \mathcal{X} \subset \mathbb{R}^d \) to be an open subset for some \( d \geq 1 \), \( \mathcal{B}(\mathcal{X}) \) the Borel \( \sigma \)-algebra on \( \mathcal{X} \), and \( \mathcal{Z} \) a finite or countable set.

%\arnaud{Voir sur cette remarque} We emphasize that, although we consider here the case \( \mathcal{X} \subset \mathbb{R}^d \), all our results naturally extend to a more general framework where \( \mathcal{X} \) is a topologically complete separable metric space, equipped with a Borel measure \( \mu \) that is nonatomic and normalized (see Remark~\ref{rk: extension functional space} for a more detailed discussion). This generalization is particularly interesting and useful for applications concerning the statistical analysis of stochastic processes driven by Brownian motion, under local differential privacy constraints. In such contexts, where the observations are trajectories of the process, the natural underlying space is the Wiener space equipped with the Wiener measure. It is therefore important to highlight that our theoretical results remain valid in this broader setting as well.

%\medskip

Note that the problem formulation introduced at the beginning of Section~\ref{s: pb formulation and main} was not restricted to the discrete case, and the notation used here remains fully consistent with that section. In particular, we denote by $\mathcal{Q}_\alpha$ the set of randomization mechanisms from $\mathcal{X}$ to $\mathcal{Z}$ of the form $q_x(z) = \mathbb{P}(Z = z \mid X = x)$, such that
\begin{equation*}
	\forall z \in \mathcal{Z}, \quad \frac{q_x(z)}{q_{x'}(z)} \le e^\alpha, \quad \text{for a.e. } dx \, dx'.
\end{equation*}
The key point now lies in extending the definition of extremal mechanisms to the continuous setting. As in the discrete case, such randomizations will play a central role in the study of optimal privacy mechanisms.

To this end, we introduce notation that enables us to generalize the space of extremal vectors introduced in~\eqref{eq : R matrix by slice} (or equivalently, the staircase mechanisms defined in Definition~\ref{def: staircase matrix}) to the continuous case. From this, we prove the extension of the results from Section~\ref{ss: finite}.

We define the set
\begin{align*}
	\theconvex = \left\{ v : \mathcal{X} \rightarrow \mathbb{R} \text{ measurable} \, \middle| \, 1 \leq v(x) \leq e^\alpha \text{ a.e.} \right\} \subset L^\infty(\mathcal{X}, \mathbb{R}),
\end{align*}
which plays the role of the hyperrectangle $[1, e^\alpha]^d$ in the discrete case (where $d = |\mathcal{X}|$). As the extremal points of $[1, e^\alpha]^d$ were key in the discrete setting, the extremal points of $\theconvex$ will now play an analogous role. This gives rise to the set
\begin{align*}
	\mathcal{E} = \left\{ r : \mathcal{X} \rightarrow \mathbb{R} \text{ measurable} \, \middle| \, r(x) \in \{1, e^\alpha\} \text{ a.e.} \right\}.
\end{align*}
For any $r \in \mathcal{E}$, define the sets
\[
F^+_r := r^{-1}(\{e^\alpha\}) \quad \text{and} \quad F^-_r := r^{-1}(\{1\}),
\]
so that, $dx$-almost everywhere, we can write
\begin{equation}{\label{eq: def r(x)}}
r(x) = e^\alpha \indi{F^+_r}(x) + \indi{F^-_r}(x) = 1 + (e^\alpha - 1)\indi{F^+_r}(x).
\end{equation}
These sets and expressions are the continuous analogues of the discrete objects introduced in \eqref{eq: Fbeta+ 4.8} and \eqref{eq: r beta x}.


As detailed in Section~\ref{S: tools} on functional analytic tools, we endow $\theconvex$ with a distance $d_\star$ that metrizes the weak-$\star$ topology. The set $\theconvex$ is convex, compact, and separable (see Lemmas~\ref{l: B closed} and~\ref{l: B separable}). Moreover, $\mathcal{E}$ coincides with the set of extremal points of $\theconvex$ (see Lemma \ref{lem : extremal points}). By Theorem~\ref{th: Choquet}, $\mathcal{E}$ is a $G_\delta$ set (i.e., a countable intersection of open sets), and since we equip $\theconvex$ with the Borel $\sigma$-algebra $\mathcal{B}(\theconvex)$, it follows that $\mathcal{E}$ is measurable. 

Note that $\theconvex \subset \widehat{B} \subset L^\infty(\mathcal{X}, \mathbb{R})$, 
where $\widehat{B}$ is the closed ball $B_{\norm{\cdot}_\infty}(0,e^\alpha)$,
so evaluating an element of $\theconvex$ at a point $x \in \mathcal{X}$ is essential. The evaluation operator $e$ defined and extensively studied in Section \ref{ss:point eval} plays a key role in this context:
\begin{equation*}
	e : 
	\begin{cases}
		(\mathcal{X} \times  \widehat{B}, \mathcal{B}(\mathcal{X}) \otimes \mathcal{B}( \widehat{B})) \to \mathbb{R} \\
		(x, r) \mapsto e_x(r).
	\end{cases}
\end{equation*}
%This operator is extensively studied in Subsection~\ref{ss:point eval}. 
We show in Subsection~\ref{ss:point eval} it is measurable and satisfies, for all $r \in  \widehat{B}$, $e_x(r) = r(x)$ for a.e.\ $x$. 

We are now ready to introduce extremal randomizations in the continuous setting. For any non-negative Borel measure $\mu$ on $\mathcal{E}$ satisfying
\begin{equation} \label{eq : cond norm mu continue}
	\text{for a.e.\ } x \in \mathcal{X}, \quad \int_{\mathcal{E}} e_x(r) \, \mu(dr) = 1,
\end{equation}
we define a randomization mechanism from $\mathcal{X}$ to $\mathcal{E}$ given by the kernel
\begin{equation} \label{eq : def random generic extr ctn}
	q^{(\mu)}_x(dr) = e_x(r) \mu(dr).
\end{equation}
This randomization is extremal in the sense that
\begin{align*}
	\frac{q^{(\mu)}_x(dr)}{q^{(\mu)}_{x'}(dr)} 
	= \frac{e_x(r)}{e_{x'}(r)} 
	= \frac{r(x)}{r(x')} \in \{e^{-\alpha}, 1, e^{\alpha}\}, \quad \text{for a.e.\ } dx \, dx'.
\end{align*}
Since $e_x(r) \in [1, e^\alpha]$, condition~\eqref{eq : cond norm mu continue} implies the bounds
\begin{equation}{\label{eq: bound mu E}}
e^{-\alpha} \leq \mu(\mathcal{E}) \leq 1.
\end{equation}

\subsection{A factorization lemma}{\label{ss: Fisher}}

We emphasize that the randomizations just introduced constitute a natural generalization of the extremal randomizations appearing in Lemma~\ref{lem : factorisation cas fini}. In the finite case, the randomization takes values in the discrete set $\{0, \dots, 2^{\Card(\mathcal{X})} - 1\}$, which indexes the extremal vectors appearing in~\eqref{eq : R matrix by slice}. In contrast, in the continuous setting, the randomization takes values in the set $\mathcal{E}$ of measurable functions with values in $\{1, e^\alpha\}$—the analogous of the extremal vectors from the discrete case.

The normalization condition~\eqref{eq : cond norm mu continue} is the continuous counterpart of the constraint~\eqref{eq: constraint factorisation 4.55} in the finite case. Similarly, the sub-probability measure $\mu$ on $\mathcal{E}$ plays the role of the coefficient vector $(\omega_\beta)_{\beta = 0, \dots, 2^{|\mathcal{X}|} - 1}$ used in the discrete setting.

\medskip
With this analogy in mind, we are now ready to extend the factorization result of Lemma~\ref{lem : factorisation cas fini} to the case where $\mathcal{X}$ is a continuous space.

\begin{prop} \label{prop : factor ctn}
Let $q$ be an $\alpha$-LDP mechanism from an open subset $\mathcal{X} \subset \mathbb{R}^d$ to a finite or countable space $\mathcal{Z}$. Then there exist a non-negative Radon measure $\mu$ on $\mathcal{E}$ and a randomization mechanism $q^{(2)}$ from $\mathcal{E}$ to $\mathcal{Z}$,
such that the mechanism $q$ admits the factorization
\[
	q = q^{(2)} \circ q^{(\mu)},
\]
where $q^{(\mu)}$ is defined by $q^{(\mu)}_x(dr) = e_x(r)\mu(dr)$, and the measure $\mu$ satisfies
$e^{-\alpha} \le \mu(\mathcal{E}) \le 1$.
\end{prop}




\begin{proof}
As $q$ satisfies the $\alpha$-LDP constraint, we have for all $z \in \mathcal{Z}$:
\[
\frac{q_x(z)}{q_{x'}(z)} \le e^\alpha, \quad \text{for a.e. } (x, x') \in \mathcal{X} \times \mathcal{X}.
\]
This implies that the function
\[
\underline{q}(z) := \essinf_{x \in \mathcal{X}} q_x(z),
\]
defined analogously to the minimal mass function in Lemma~\ref{lem : factorisation cas fini}, satisfies $\underline{q}(z) > 0$. Define for each $z \in \mathcal{Z}$ and $x \in \mathcal{X}$ the ratio
\[
u_x(z) := \frac{q_x(z)}{\underline{q}(z)}.
\]
By construction, for all $z \in \mathcal{Z}$, the function $x \mapsto u_x(z)$ satisfies $1 \le u_x(z) \le e^\alpha$ for almost every $x \in \mathcal{X}$. 

Let us denote $v_z : x \mapsto u_x(z)$, which defines a function in the compact convex set $\theconvex$ introduced earlier. We apply Choquet's Theorem (see Theorem~\ref{th: Choquet}), whose applicability to the set $\theconvex$ is justified in Section~\ref{S: tools}. This yields the existence of a probability measure $\mu_z$ on $\mathcal{E}$ such that
\[
v_z = \int_{\mathcal{E}} r \, \mu_z(dr),
\]
where the integral is taken in the Bochner sense. %, and the equality holds in the space $\widehat{B}$ (see Section~\ref{ss:point eval} for details). 
By Proposition~\ref{p: evaluation}, applying the evaluation operator $e_x$ gives
\[
e_x(v_z) = \int_{\mathcal{E}} e_x(r) \mu_z(dr).
\]
Note that for a.e.\ $x \in \mathcal{X}$, we have $e_x(v_z) = v_z(x) = u_x(z)$, so that
\[
u_x(z) = \int_{\mathcal{E}} e_x(r) \mu_z(dr), \quad \text{for a.e.\ } x \in \mathcal{X},\ \forall z \in \mathcal{Z}.
\]
Recalling that $q_x(z) = u_x(z) \underline{q}(z)$, we deduce
\begin{equation}
\label{eq : q x with mu r in proof}
q_x(z) = \int_{\mathcal{E}} e_x(r) \mu_z(dr)\underline{q}(z), \quad \text{for a.e.\ } x \in \mathcal{X},\ \forall z \in \mathcal{Z}.
\end{equation}
Since $q$ is a randomization mechanism, we have
\[
\sum_{z \in \mathcal{Z}} q_x(z) = 1, \quad \text{for a.e.\ } x \in \mathcal{X}.
\]
Using this and~\eqref{eq : q x with mu r in proof}, we find for a.e.\ $x \in \mathcal{X}$:
\begin{align*}
1 &= \sum_{z \in \mathcal{Z}} \int_{\mathcal{E}} e_x(r) \mu_z(dr) \underline{q}(z) = \int_{\mathcal{E}} e_x(r) \left( \sum_{z \in \mathcal{Z}} \underline{q}(z) \mu_z(dr) \right),
\end{align*}
where the second line uses Fubini-Tonelli. We now define the measure
\begin{equation}
\label{eq : def mu proof fact ctn}
\mu := \sum_{z \in \mathcal{Z}} \underline{q}(z) \mu_z,
\end{equation}
which is a non-negative measure on $\mathcal{E}$. With this definition, we have shown:
\[
\int_{\mathcal{E}} e_x(r) \mu(dr) = 1, \quad \text{for a.e.\ } x \in \mathcal{X},
\]
which is exactly the normalization condition~\eqref{eq : cond norm mu continue}. Consequently, the measure $\mu$ defines a randomization mechanism with values in $\mathcal{E}$, via
\[
q^{(\mu)}_x(dr) := e_x(r) \mu(dr).
\]
It remains to define a randomization $q^{(2)}$ from $\mathcal{E}$ to $\mathcal{Z}$ such that $q = q^{(2)} \circ q^{(\mu)}$. From~\eqref{eq : def mu proof fact ctn}, we know that for each $z \in \mathcal{Z}$, the measure $\underline{q}(z) \mu_z$ is absolutely continuous with respect to $\mu$. By the Radon-Nikodym theorem, there exists a non-negative measurable function $\tilde{q}^{(2)}_z : \mathcal{E} \to \mathbb{R}_+$ such that
\begin{equation}
\label{eq : Rad Nyk in proof}
\underline{q}(z) \mu_z(dr) = \tilde{q}^{(2)}_z(r) \mu(dr).
\end{equation}
Observe that
\begin{equation}
\label{eq: 21.5}
\left( \sum_{z \in \mathcal{Z}} \tilde{q}^{(2)}_z(r) \right) \mu(dr) = \sum_{z \in \mathcal{Z}} \underline{q}(z) \mu_z(dr) = \mu(dr),
\end{equation}
which implies that $\sum_{z \in \mathcal{Z}} \tilde{q}^{(2)}_z(r) = 1$ for $\mu$-a.e.\ $r \in \mathcal{E}$.

We now define the randomization $q^{(2)}$ from $\mathcal{E}$ to $\mathcal{Z}$ by:
\[
q^{(2)}_r(z) := 
\begin{cases}
\tilde{q}^{(2)}_z(r) & \text{if } \sum_{z \in \mathcal{Z}} \tilde{q}^{(2)}_z(r) = 1, \\
\indi{z_1}(z) & \text{otherwise},
\end{cases}
\]
for some fixed $z_1 \in \mathcal{Z}$. Note that $\mu(dr)$-a.e.,
%on the support of $\mu$, 
$q^{(2)}_r(z) = \tilde{q}^{(2)}_z(r)$ for all $z \in \mathcal{Z}$, and the definition outside the support of $\mu$ is irrelevant.

We can now return to~\eqref{eq : q x with mu r in proof}, and using~\eqref{eq : Rad Nyk in proof}, we write for all $z\in\cZ$ and $dx$-a.e. :
\begin{align*}
q_x(z) &= \int_{\mathcal{E}} e_x(r) \tilde{q}^{(2)}_z(r) \mu(dr) = \int_{\mathcal{E}} e_x(r) q^{(2)}_r(z) \mu(dr) = \int_{\mathcal{E}} q^{(2)}_r(z) q^{(\mu)}_x(dr),
\end{align*}
which establishes that
\begin{equation}
\label{eq : composition in the proof ctn}
q = q^{(2)} \circ q^{(\mu)}.
\end{equation}
This completes the factorization construction in the continuous case.
\end{proof}


\begin{rem}
\begin{enumerate}
    \item The proof of Proposition \ref{prop : factor ctn} follows the general scheme of the proof of Lemma \ref{lem : factorisation cas fini}, but in a more involved setting. A key aspect for the application of Choquet's theorem is the compactness of the set $\theconvex$, which motivates the introduction of the weak-$\star$ topology. However, as pointed out in Remark \ref{rk: E not closed} below, the set of extremal points $\cE$ is not closed in the weak-$\star$ topology, and thus not compact. This represents a fundamental difference with the finite case, where the set of extremal points is finite, and hence compact.
    
    \item The proposition above shows that any $\alpha$-LDP randomization can be factorized as the composition of two randomizations, one of which is an explicit extremal randomization $q^{(\mu)}$, fully described by a measure $\mu$. As in the discrete case, this representation will be instrumental in deriving upper bounds on the Fisher information of $\alpha$-LDP randomizations. Indeed, we can consider maximizing the Fisher information over all extremal randomizations, that is, over all possible measures $\mu$. However, since $\cE$ is not compact in general (in the weak-$\star$ topology), the set of sub-probability measures on $\cE$ is not compact either . This presents a major challenge in establishing the existence of a solution to the corresponding optimization problem, as discussed in detail before Proposition \ref{prop : maximisation pb convex}.
\end{enumerate}
\end{rem}


\subsection{Fisher information in the continuous case}{\label{Ss : Fisher cont}}
As in the discrete case, we now investigate how the factorization lemma can be leveraged to derive an expression for the maximal Fisher information. Again, addressing this problem requires some regularity assumptions on the statistical model.

Assume that $\cX$ is an open subset of $\R^d$ and that $(P_{\theta})_{\theta \in \Theta}$ is a family of probability measures on $\cX$, where $\Theta \subset \R^{d_\Theta}$. We suppose that this family is dominated by the Lebesgue measure and denote its density by $p_{\theta}(x) := \frac{dP_\theta(x)}{dx}$.

Let $\alpha > 0$ and let $\mu$ be a Radon sub-probability measure on $\cE$ satisfying \eqref{eq : cond norm mu continue}. We consider the privacy mechanism $q^{(\mu)}$ as introduced in Section \ref{ss: infinite}, given by $q^{(\mu)}_x(dr) = e_x(r)\mu(dr)$. Denote by $Z$ the associated public data, and by $(\widetilde{P}_{\theta})_{\theta \in \Theta}$ the induced model on $\cE$, i.e., the law of $Z$.

The following lemma shows that this model admits a continuous density and is differentiable in quadratic mean, with an explicit and continuous score function.
\begin{lem} \label{lem : expression tilde p et score ctn}
Let $\theta_0 \in \overset{\circ}{\Theta}$. The following statements hold:
\begin{enumerate}
	\item The probability measure $\widetilde{P}_{\theta}$ is absolutely continuous with respect to $\mu$ and its density is given by
	\begin{equation} \label{eq : tilde p ctn}
		\widetilde{p}_\theta(r) := \frac{d\widetilde{P}_{\theta}}{d\mu}(r)
		= \int_\cX r(x) p_\theta(x) dx.
	\end{equation}
	\item Assume that the model $(P_{\theta})_{\theta \in \Theta}$ is DQM at $\theta_0$ with score function $s_{\theta_0}$ such that $\E_{\theta_0}\left[\norm{s_{\theta_0}(X)}^2\right] < \infty$. Then the model $(\widetilde{P}_{\theta})_{\theta \in \Theta}$ is also DQM at $\theta_0$, and its score function is given by
	\begin{equation} \label{eq : score t ctn}
		t_{\theta_0}(r) = \frac{\int_\cX s_{\theta_0}(x) r(x)p_{\theta_0}(x) dx}{\widetilde{p}_{\theta_0}(r)} = (e^\alpha - 1) \frac{\int_{F_r^+} s_{\theta_0}(x) p_{\theta_0}(x) dx}{\widetilde{p}_{\theta_0}(r)}.
	\end{equation}
\item The functions $r \mapsto \widetilde{p}_{\theta_0}(r)$ and $r \mapsto t_{\theta_0}(r)$ are continuous on $\cE$.
\end{enumerate}
\end{lem}
\begin{proof}
\begin{enumerate}
	\item The law of $Z$ is given by
	\[
		\widetilde{P}_{\theta}(dr) = \int_\cX q^{(\mu)}_x(dr) p_\theta(x)dx = \int_\cX e_x(r) \mu(dr) p_\theta(x)dx = \left( \int_\cX r(x) p_\theta(x) dx \right) \mu(dr),
	\]
	where we used that $e_x(r) = r(x)$ $dx$-a.e. for $r \in \cE$. Hence, the density with respect to $\mu$ is as in \eqref{eq : tilde p ctn}.
	\item Applying Lemma 3.1 in \cite{steinbergerEfficiencyLocalDifferent}, we know that the model $(P_\theta)$ is DQM with score function $s_{\theta_0}$. From Lemma A1 in the same reference, we obtain:
	\[
		t_{\theta_0}(r) = \frac{1}{\widetilde{p}_{\theta_0}(r)} \int_\cX s_{\theta_0}(x) e_x(r) p_{\theta_0}(x) dx = \frac{\int_\cX s_{\theta_0}(x) r(x) p_{\theta_0}(x) dx}{\widetilde{p}_{\theta_0}(r)}.
	\]
    Recalling that $r(x)$ is as in \eqref{eq: def r(x)} and that $\int_{\mathcal{X}} s_{\theta_0}(x)p_{\theta_0}(x) dx = 0$ we deduce the second formula.
	\item The continuity of $r \mapsto \widetilde{p}_{\theta_0}(r)$ follows from \eqref{eq : tilde p ctn}, the integrability of $p_{\theta_0}$, and the definition of the weak-$\star$ topology on $\cE$. Analogously, the continuity of $r \mapsto t_{\theta_0}(r)$ follows from \eqref{eq : score t ctn} and the integrability of $s_{\theta_0}(x)p_{\theta_0}(x)$.
\end{enumerate}
\end{proof}
We now aim to use the factorization lemma to study the %maximal 
information $\cI_{\theta_0}(q\circ \cP)$ that can be extracted from an $\alpha$-LDP view of the private data through a randomisation mechanism $q$ from $\cX$ to some finite or countable output space $\cZ$.
%To this end, define
%\begin{equation*}
%	\cJ^{\text{max},\alpha}_{\theta_0} := \sup_{q \in \cQ_\alpha} \cI_{\theta_0}(q \circ \cP),
%\end{equation*}
%where the supremum runs over all $\alpha$-LDP mechanisms $q$ from $\cX$ to some finite or countable output space $\cZ$. 
Using the factorization result in Proposition \ref{prop : factor ctn}, we obtain the inequality:
\begin{align}{\label{eq: fact on Fisher cont}}
	\mathcal{I}_{\theta_0}(q \circ \mathcal{P}) 
	&= \mathcal{I}_{\theta_0}((q^{(2)} \circ q^{(\mu)}) \circ \mathcal{P}) \\
	&= \mathcal{I}_{\theta_0}(q^{(2)} \circ (q^{(\mu)} \circ \mathcal{P})) 
	\le \mathcal{I}_{\theta_0}(q^{(\mu)} \circ \mathcal{P}), \nonumber
\end{align}
where we used \eqref{eq : decroissance Fisher}.
Therefore, to study the maximization of the Fisher information, we may restrict our attention to extremal mechanisms of the form $q = q^{(\mu)}$. Thanks to Lemma \ref{lem : expression tilde p et score ctn}, we can write the Fisher information explicitly:
\begin{align}\nonumber
	\cI_{\theta_0}(q^{(\mu)}\circ \cP)
	&= \E_{\theta_0} \left[ t_{\theta_0}(Z) t_{\theta_0}(Z)^T \right] = \int_\cE t_{\theta_0}(r) t_{\theta_0}(r)^T \widetilde{p}_{\theta_0}(r) \mu(dr) \\
	&=\nonumber  \int_{\cE} 
	\frac{\left( \int_\cX s_{\theta_0}(x) r(x) p_{\theta_0}(x) dx \right)
	\left( \int_\cX s_{\theta_0}(x) r(x) p_{\theta_0}(x) dx \right)^T}
	{\widetilde{p}_{\theta_0}(r)} \mu(dr) \\
    & =\label{eq : Fisher ctn general_preli}(e^\alpha - 1)^2 \int_{\cE} 
	\frac{\left( \int_{F_r^+} s_{\theta_0}(x) p_{\theta_0}(x) dx \right)
	\left( \int_{F_r^+} s_{\theta_0}(x)p_{\theta_0}(x) dx \right)^T}
	{\widetilde{p}_{\theta_0}(r)} \mu(dr).
\end{align}
This representation is a key technical tool to derive the next theorem, which provides general lower and upper bounds for the Fisher information under $\alpha$-LDP constraints, valid for any $\alpha > 0$, and characterizes its behavior in the asymptotic regime $\alpha \rightarrow 0$.

If the parameter space is one-dimensional, the Fisher information is a scalar and we can define
\begin{equation*}
	\cJ^{\text{max},\alpha}_{\theta_0} := \sup_{q \in \cQ_\alpha} \cI_{\theta_0}(q \circ \cP),
\end{equation*}
where the supremum runs over all $\alpha$-LDP mechanisms $q$ from $\cX$ to some finite or countable output space $\cZ$.


\begin{thm} \label{thm : equivalent alpha petit ctn}
Assume that $\Theta \subset \R$ and that the model $(P_{\theta})_{\theta \in \Theta}$ is DQM at $\theta_0$ with score function $s_{\theta_0}$ such that $\E_{\theta_0}\left[\norm{s_{\theta_0}(X)}^2\right] < \infty$.
	Then we have, for any $\alpha>0$,
	\begin{equation} \label{eq :  encadrement Fisher ctn small alpha}
 \frac{(e^\alpha-1)^2}{2e^\alpha(1+e^\alpha)}
 %	\frac{(e^\alpha-1)^2}{1+e^\alpha} \frac{1}{2e^\alpha} 
\left( \int_\cX \abs{s_{\theta_0}(x)}p_{\theta_0}(x)dx \right)^2
\le
		\cJ^{\text{max},\alpha}_{\theta_0} \le \frac{(e^\alpha-1)^2}{4}
		\left( \int_\cX \abs{s_{\theta_0}(x)}p_{\theta_0}(x)dx \right)^2.	
	\end{equation}
	Consequently,
\begin{equation*}
\cJ^{\text{max},\alpha}_{\theta_0} \sim_{\alpha\to0} \frac{\alpha^2}{4}
	\left( \int_\cX \abs{s_{\theta_0}(x)}p_{\theta_0}(x)dx \right)^2.		
\end{equation*}	
\end{thm}


\begin{rem}
We can compare Theorem~\ref{th : info cas fini} in the finite case with Theorem~\ref{thm : equivalent alpha petit ctn}, for the continuous setting. Theorem~\ref{th : info cas fini} provides the exact value of the optimal Fisher information in the discrete case, whereas Theorem~\ref{thm : equivalent alpha petit ctn} gives an asymptotic characterization. These two results are consistent with each other. Let us emphasize that the $\alpha$-LDP mechanism introduced in the proof of the lower bound in Theorem~\ref{thm : equivalent alpha petit ctn} is precisely the continuous analogue of the mechanism used for the lower bound in the discrete case, in Proposition~\ref{prop : borne inf Fisher fini}.
\end{rem}

\begin{rem}
It is interesting to note that, in the case where the score function is given by $s = \frac{\dot{p}_{\theta_0}}{p_{\theta_0}}$ 
and we consider a Gaussian distribution, the quantity $\int_{\mathcal{X}} \abs{s_{\theta_0}(x)} \, p_{\theta_0}(x) \, dx$
is easily computable and equals \( \sqrt{\frac{2}{\pi}} \). We can compare this with the results in \cite{KalSte} concerning the efficient estimation of a Gaussian mean under local differential privacy. In that work, the optimal Fisher information is shown to be 
$\frac{2}{\pi}\frac{(e^\alpha - 1)^2}{(e^\alpha + 1)^2}$.
Notably, in the limit as \( \alpha \to 0 \), one recovers the same result as in our calculation. It is also interesting to compare this with the exact maximal Fisher information derived in the discrete case, in Equation \eqref{eq : def I max caf fini}. Specifically, for a Gaussian variable, it is \( n_{\max} = \frac{1}{2} \), which simplifies the denominator to $\left( \frac{e^\alpha + 1}{2} \right)^2,$
leading to the same optimal Fisher information as reported in \cite{KalSte}. When \( \alpha \) is close to zero, this denominator is asymptotically equivalent to 1, which explains why neglecting it in \eqref{eq : def I max caf fini} still provides the correct asymptotic result. %However, for general values of \( \alpha \), the denominator is no longer negligible. The difficulty is that, in the continuous setting, it is not possible to analyze the denominator in as much detail as in the discrete case (see Remark \ref{rk} for further discussion).

\end{rem}
\begin{rem}
	The representation \eqref{eq : Fisher ctn general_preli} is valid for multidimensional parameters. However, it is not clear how results like Theorems \ref{th : info cas fini}--\ref{thm : equivalent alpha petit ctn} can be extended to this setting. In particular, the definition of the set $F^+_\text{max}$ introduced in Remark \ref{rk: choice Fmax} does not appear to generalize straightforwardly in the multidimensional case.
\end{rem}


We now move to the proof that, for any regular statistical model and given $\alpha > 0$, there exists an extremal mechanism $q^{(\mu)}$ which maximizes the Fisher information.

The key observation is that the continuous version of the maximization problem given in~\eqref{eq: opt 14.75} admits a solution. This is stated formally in the proposition below, the proof of which is provided in Section~\ref{ss: proof cont}. 
%The proof relies on tools from set theory, such as Zorn's lemma, along with convexity arguments. Additional details are provided in Remark \ref{rk: E not closed}.

\begin{prop} \label{prop : maximisation pb convex}
Let $j : \begin{cases} \theconvex \to \R_+ \\ r \mapsto j(r) \end{cases}$ be a bounded, continuous, and convex function. Consider the optimization problem:
\begin{equation} \label{eq : optimisation pb ctn J}
	J^* := \sup_{\mu} \int_\cE j(r) \, \mu(dr), \quad
	\text{subject to } \int_\cE e_x(r) \, \mu(dr) = 1 \quad \text{for $dx$-almost every } x.
\end{equation}
where the supremum is taken over all Radon sub-probability measures $\mu$ on $\cE$.

Then, there exists a Radon sub-probability measure $\overline{\mu}$ on $\cE$ such that
\[
	J^* = \int_\cE j(r) \, \overline{\mu}(dr)
	\quad \text{and} \quad
	\int_\cE e_x(r) \, \overline{\mu}(dr) = 1 \quad \text{for $dx$-almost every } x.
\]
\end{prop}

\begin{rem}{\label{rk: E not closed}}
The proof of the above proposition relies on the compactness of $\theconvex$ under the weak-$\star$ topology (see Lemma~\ref{l: B closed}), which allows one to establish existence of a solution over $\theconvex$. However, the set $\mathcal{E}$ is not closed in general (see a counter example in Remark~\ref{rem : E not closed}), which prevents us from directly concluding that the solution to the maximization problem is supported in $\mathcal{E}$ via this route.

Nevertheless, we are able to construct an optimal measure supported on $\mathcal{E}$ by using convexity arguments and several tools from set theory, including techniques akin to those used in the proof of Choquet’s theorem, as developed in~\cite{Choquet}, whose main tool to derive the existence of measures supported on $\cE$ is the 
%use of 
Zorn's lemma.
\end{rem}

We now apply Proposition~\ref{prop : maximisation pb convex} to the problem of maximizing the Fisher information, in order to deduce that the corresponding optimization problem admits a solution. 

To simplify the notation, let us define, for any $r \in \theconvex$,
\begin{equation} \label{eq : def i r ctn convex}
	i(r) := \frac{ \left( \int_\cX  s_{\theta_0}(x) \, r(x) \, p_{\theta_0}(x) \, dx \right)^2 }{
		\int_{\cX} r(x) \, p_{\theta_0}(x) \, dx } .
\end{equation}
For $r \in \cE$, by using  \eqref{eq: def r(x)}, we have
\begin{equation}		\label{eq : def i r ctn extremal}
	i(r)	= (e^\alpha - 1)^2 \frac{ \left( \int_{F_r^+}  s_{\theta_0}(x) \, p_{\theta_0}(x) \, dx \right)^2 }{ 1 + (e^\alpha - 1)
		\int_{F_{r}^+} p_{\theta_0}(x) \, dx }.
\end{equation}
 Observe that the quantity multiplying \( (e^\alpha - 1)^2 \) above plays exactly the same role as \( i_\beta \) defined in \eqref{eq : def i beta}, in the finite case.
We are thus led to consider the following optimization problem:
\begin{equation} \label{eq : optimisation pb ctn}
	I^* := \sup_{\mu} \int_\cE i(r) \, \mu(dr), \quad
	\text{subject to } \int_\cE e_x(r) \, \mu(dr) = 1 \quad \text{for $dx$-almost every } x,
\end{equation}
where the supremum is taken over all Radon sub-probability measures $\mu$ on $\cE$.

By Proposition~\ref{prop : maximisation pb convex}, we deduce that the above optimization problem admits a solution. Its applicability is formalized in the following theorem, whose proof can be found in Subsection~\ref{ss: proof cont}. 

\begin{thm}{\label{thm: main opt cont}}
There exists a Radon sub-probability measure $\overline{\mu}$ on $\cE$ such that the mechanism defined by $q^{(\overline{\mu})}(dr) := e_x(r) \, \overline{\mu}(dr)$ satisfies
\begin{equation*}
	\cJ^{\text{max},\alpha}_{\theta_0} = \cI_{\theta_0}\left( q^{(\overline{\mu})} \circ \cP \right).
\end{equation*}
\end{thm}
We have therefore established that one can always find an extremal privacy mechanism that reaches the optimal Fisher information for $\Theta \subset \mathbb{R}$. Moreover, we have explicitly characterized this information in the regime where the privacy parameter $\alpha$ tends to zero. Let us stress that Proposition \ref{prop : maximisation pb convex} is a general result which can be used to maximize other criteria than the Fisher information. 
%\arnaud{C'est un peu flou, car il y a pas d'autre criteria..Cela pourrait etre la trace de la matrice d'Info par ex.}


%The preceding analysis highlights the crucial role played by the structure of the set \(\mathcal{E}\) in shaping the behavior of the Fisher information. As in the discrete setting, the extremal values \(1\) and \(e^\alpha\) appear to be particularly significant. This naturally raises the question: could the analysis be simplified by considering measures on \(\mathcal{E}\) that are, in some sense, insensitive to the distinction between these two values? This question remains open and is left for future investigation. 








\section{Application to the uniform law}{\label{s: unif}}

In this section, we consider the situation of a private variable following a uniform law on the interval $[0,\theta]$ where $\theta>0$ is an unknown parameter. We have $\cX=[0,\infty)$ and $p_{\theta}(x)=\frac{1}{\theta} \indi{[0,\theta]}(x)$. The model is not regular, and it is well known that the rate of estimation of $\theta$ is proportional to the number of data $n$ instead of the usual $\sqrt{n}$ corresponding to finite Fisher information situations.
Consequently, the results of the Section \ref{Ss : Fisher cont} do not apply in this case. However, the factorization result holds true, and so we know that any $\alpha$-LDP mechanism can be factorized through a mechanism $q^{(\mu)}$ as described in Section \ref{ss: infinite}. 
The next proposition shows that for any such channel $q^{(\mu)}$ the model is DQM for almost all values of the parameter. We also give an upper bound on the Fisher information.

It shows that under $\alpha$-LDP constraint, the best rate of estimation for the range of a uniform variable is $\sqrt{n}$, with some upper bound on the Fisher information. The proof of this proposition also elucidate which $\alpha$-LDP mechanism is optimal in the regime $\alpha \to 0$. It leads to an estimation procedure that will be described in Section \ref{Ss : efficient est unif}.

\subsection{Upper bound on the Fisher information}

Let us consider $\mu$ any Radon sub-probability on  measure $\cE$ satisfying the normalization constraint \eqref{eq : cond norm mu continue}. We denote $q^{(\mu)}$ the extremal mechanism from $\cX=[0,\infty)$ to $\cE$ given by $q^{(\mu)}=e_x(r) \mu(dr)$. Let us denote by $\tilde{p}_\theta(dr)$ the law of the public data on $\cE$ when the private one follows the uniform law on $[0,\theta]$. We have
\begin{equation*}
	\tilde{p}_\theta(dr)=\left(\int_\cX r(x)p_\theta(x)dx\right) \mu(dr)=
	\left(\frac{1}{\theta}\int_0^\theta r(x)dx \right)\mu(dr).
\end{equation*}
\begin{prop}\label{prop : Fisher unif}
	 For almost every $\theta_0 \in (0,\infty)$ the model
	$\tilde{p}_\theta(dr)$ is DQM at $\theta_0$. The following upper bound on the Fisher information holds
	\begin{equation*}
		\cI_{\theta_0}\left((\widetilde{p}_\theta)_\theta\right) \le \frac{(e^\alpha-1)^2}{\theta_0^2}.
	\end{equation*}
\end{prop}
\begin{proof}
	The density of the model is given by 
	\begin{equation}
		\label{eq : tilde p unif}
		\tilde{p}_\theta(r) =\frac{\tilde{p}_\theta(dr)}{\mu(dr)}=\frac{1}{\theta} \int_0^\theta r(x)dx.
	\end{equation}
		 We have to prove that this density admits a derivative with respect to the parameter $\theta$. Define 
	\begin{equation*}
		\cR=\left\{(r,x)\in B \times [0,\infty) \mid e_x^+(r):=\lim_{h\to0} h^{-1}\int_x^{x+h} r(y) dy \text{ exists} \right\}.
	\end{equation*}
	By similar argument as for the proof of the measurability of the operator $e_x$ defined by \eqref{eq : def evaluation operator}, we can check that set $\cR$ is measurable with respect to $\cB(B)\otimes\cB([0,\infty))$ (see Step 4 in the proof of Proposition \ref{p: evaluation}). By Fubini theorem, we deduce
	$\int_{\cE \times [0,\infty)} \indi{\cR^c}(r,\theta) \mu(dr)\times d\theta=
\int_{\cE} \left(\int_0^\infty \indi{\cR^c}(r,\theta) d\theta \right)\mu(dr)=0$, as for any $r\in \cE$, $\lim_{h\to0} h^{-1}\int_\theta^{\theta+h} r(y) dy$ exists for almost every $\theta$ and is equal to $r(\theta)$. Now, a further application of Fubini theorem entails  
$ \int_0^\infty\left(\int_{\cE} \indi{\cR^c}(r,\theta) \mu(dr) \right)d\theta=0$.
 Consequently, there exists $\overline{\Theta}$ a set of real, with full Lebesgue
  measure, such that $\int_{\cE} \indi{\cR^c}(r,\theta) \mu(dr)=0$ for all $\theta \in \overline{\Theta}$.

Now, we consider any $\theta_0 \in \overline{\Theta}$. From the definition of $\cR$, we have that $\mu(dr)$ a.e., the function $\theta \mapsto \int_0^\theta r(x)dx$ admits a derivative at $\theta=\theta_0$, denoted as $e^+_{\theta_0}(r)$. Recalling \eqref{eq : tilde p unif}, we deduce that $\tilde{p}_\theta(r)$ admits a derivative at $\theta=\theta_0$ given by
\begin{align} \nonumber
	\frac{\partial}{\partial \theta} \tilde{p}_{\theta_0}(r)
	&= -\frac{1}{\theta_0^2}\int_0^{\theta_0} r(x)dx+\frac{e^+_{\theta_0}(r)}{\theta_0} 
	\\ \label{eq : p tilde dot unif}
	&= \frac{1}{\theta_0^2}\int_0^{\theta_0} (e^+_{\theta_0}(r)-r(x))dx.
\end{align}
We recall the definition of $F^+_r = r^{-1} \left( \{ e^\alpha \} \right),$
and that \( r(x) = 1 + (e^{\alpha} - 1) \, \indi{F^+_r}(x) \).  
Substituting this expression into the definition of \( e^+_{\theta_0}(r) \), we obtain  
\[
e^+_{\theta_0}(r) = 1 + (e^{\alpha} - 1) \, e^+_{\theta_0} \left( \indi{F^+_r} \right).
\]
It implies
$e^+_{\theta_0}(r)-r(x)=(e^{\alpha}-1)(e^+_{\theta_0}(\indi{F^+_r})-\indi{F^+_r}(x))$.
%From \eqref{eq : tilde p unif}--
Inserting in \eqref{eq : p tilde dot unif}, we deduce that
\begin{equation*}
	\frac{(\frac{\partial}{\partial \theta} \tilde{p}_{\theta_0}(r))^2}{\tilde{p}_{\theta_0}(r)}
	=  \frac{(e^\alpha-1)^2}{\theta_0^4}\frac{\left( \int_0^{\theta_0} 
	(e^+_{\theta_0}(\indi{F^+_r})-\indi{F^+_r}(x))dx    \right)^2}{\widetilde{p}_{\theta_0}(r)}.
\end{equation*}
Collecting with  \eqref{eq : tilde p unif}, we deduce the expression for the Fisher information
\begin{equation*}
	\cI_{\theta_0}\left((\widetilde{p}_\theta)_\theta\right) 
	=	\frac{(e^\alpha-1)^2}{\theta_0^3}
	\int_{\cE}
\frac{\left( \int_0^{\theta_0} 
		(e^+_{\theta_0}(\indi{F^+_r})-\indi{F^+_r}(x))dx    \right)^2}{ \int_0^{\theta_0} r(x)dx } \mu(dr).
\end{equation*}	
	Since $e^+_{\theta_0}(\indi{F^+_r}) \in [0,1]$ and $\indi{F^+_r}(x)\in \{0,1\}$, we deduce that 
	\begin{equation} \label{eq : F_max unif}
	\abs*{ \int_0^{\theta_0} 
	(e^+_{\theta_0}(\indi{F^+_r})-\indi{F^+_r}(x))dx} \le \theta_0. 
	\end{equation}
	From $r\ge1$, we have $\int_0^{\theta_0} r(x)dx \ge \theta_0$, which gives
\begin{equation*}
	\cI_{\theta_0}\left((\widetilde{p}_\theta)_\theta\right) 
	\le	\frac{(e^\alpha-1)^2}{\theta_0^3}
	\int_{\cE} \frac{\theta_0^2}{\theta_0} \mu(dr) \le \frac{(e^\alpha-1)^2}{\theta_0^2},
\end{equation*}	
	where we used that $\mu$ is a sub-probability.	
\end{proof}
\begin{rem}
The upper bound given in Proposition \ref{prop : Fisher unif} is very similar
to the upper bound in Theorem \ref{thm : equivalent alpha petit ctn} for regular models. In particular, this upper bound is obtained by searching for $F_r^{+}$ achieving the maximum of the LHS 
in \eqref{eq : F_max unif}. We can see that \eqref{eq : F_max unif}
becomes an equality when $F_r^+=[0,\theta_0)$ and  $F_r^+=[\theta_0,\infty)$. It suggests that the measure $\mu$, supported on the two elements $r_1=1+(e^{\alpha}-1)\indi{[0,\theta_0)}$ and $r_2=1+(e^{\alpha}-1)\indi{[\theta_0,\infty)}$, makes the Fisher information close to the optimal one in the regime $\alpha \to 0$. %This mechanism is the analogue, in this setting, of the mechanism introduced in Thoerem \ref{thm : equivalent alpha petit ctn} for the regular case.
\end{rem}


\subsection{Efficient estimator in the regime $\alpha \to 0$} \label{Ss : efficient est unif}	
	Assume that $X_1,\dots,X_n$ is a sampling of random variables with uniform distribution on $[0,\theta]$, and $\theta>0$. We denote by $\theta_0$ the true value of the parameter.

 From the previous section, we expect that a Bernoulli randomization mechanism is efficient when $\alpha \to 0$. However, the construction of the Bernoulli mechanism necessitates the knowledge of $\theta_0$.
 
 Let us assume that we have a preliminary estimation of the parameter, which we call $\widehat{\theta}^{\text{p}}>0$. For simplicity of the exposition, we consider that 
 $\widehat{\theta}^{\text{p}}$ is a deterministic value. It can be a preliminary estimator independent from $X_1,\dots,X_n$, and our analysis is then conditional to $\widehat{\theta}^{\text{p}}$. Remark that we do not need $\widehat{\theta}^{\text{p}}$ to be close to $\theta_0$.
 
 We introduce the $\alpha$-LDP extremal mechanism corresponding to a two point supported measure $\mu$ on $\cE$, which we write in a more amenable way as a randomization valued on $\cZ=\{0,1\}$, by
 \begin{equation*}
 	P(Z=z \mid X=x)=q_x(z)=
 	\begin{cases}
 		\frac{1}{1+e^\alpha}\left(e^\alpha \indi{\{1\}}(z)+\indi{\{0\}}(z) \right),
 		\text{if $x \in F^+:=[0,\widehat{\theta}^{\text{p}})$}
 		\\
 		\frac{1}{1+e^\alpha}\left( \indi{\{1\}}(z)+e^\alpha \indi{\{0\}}(z) \right),
 		\text{if $x \in (F^+)^c=[\widehat{\theta}^{\text{p}} ,\infty)$}.
 	\end{cases}
 \end{equation*}
We draw the public views $Z_1,\dots,Z_n$ according to this channel. 
Then, we have 
\begin{align} \nonumber
	\widetilde{p}_{\theta}(1)=P_\theta(Z=1)&=\frac{1}{\theta}\int_0^{\widehat{\theta}^{\text{p}} \wedge \theta} \frac{e^\alpha}{1+e^\alpha} dx +
	\frac{1}{\theta}
	\int_{\widehat{\theta}^{\text{p}} \wedge \theta}^{\theta} \frac{1}{1+e^\alpha} dx 
	\\&= \label{eq : tilde p unif estimateur}
	\frac{1}{1+e^\alpha} + \frac{e^\alpha-1}{1+e^\alpha} \left(\frac{\widehat{\theta}^{\text{p}}}{\theta} \wedge 1 \right).
 %	\in [\frac{1}{1+e^\alpha},\frac{e^\alpha}{1+e^\alpha}].  
\end{align}
 Estimating $\widetilde{p}_{\theta}(1)$ by $\overline{Z}_n:=\frac{1}{n}\sum_{i=1}^n Z_i$ gives the following estimator for $\theta$
 \begin{equation*}
 	\widehat{\theta}_n:= \widehat{\theta}^{\text{p}} \frac{e^\alpha -1}{(1+e^\alpha) \overline{Z}_n -1}.
 \end{equation*}
 \begin{prop}
 	1) We have $\widehat{\theta}_n \xrightarrow[a.s.]{n\to\infty} \theta_0 \vee \widehat{\theta}^{\text{p}}$.
 	
 	2) If $\widehat{\theta}^{\text{p}} \le \theta_0$, the estimator is consistent and
 	\begin{equation*}
 		\sqrt{n} \left(\widehat{\theta}_n-\theta_0\right)
 		\xrightarrow[\cL]{n \to \infty} \mathcal{N}\left(0,v(\theta_0,\widehat{\theta}^{\text{p}})\right),
 	\end{equation*}
 	where 
 	\begin{equation*}
 		v(\theta_0,\widehat{\theta}^{\text{p}})=
 		\frac{\theta_0^4}{(\widehat{\theta}^{\text{p}})^2} \frac{1}{(e^\alpha-1)^2}
 		\left[1+(e^\alpha-1) \frac{\widehat{\theta}^{\text{p}}}{\theta_0}\right]
 		\left[e^\alpha-(e^\alpha-1) \frac{\widehat{\theta}^{\text{p}}}{\theta_0}\right].
 	\end{equation*}
 \end{prop}
 \begin{proof} 1) This is obtained by the Law of Large Numbers, $\overline{Z}_n
 	\xrightarrow[a.s.]{n\to\infty}  \widetilde{p}_{\theta_0}(1)$ together with the expression of $\widehat{\theta}_n$ and \eqref{eq : tilde p unif estimateur}.
 	
 2) With a few lines of algebra, we obtain:
\[
\sqrt{n}\left(\widehat{\theta}_n - \theta_0\right) = - \frac{\theta_0 (1 + e^\alpha)}{(1 + e^\alpha)\overline{Z}_n - 1} \sqrt{n}\left(\overline{Z}_n - \widetilde{p}_{\theta_0}(1)\right).
\]
Indeed, we start from the expression:
\[
\widehat{\theta}_n - \theta_0 = - \frac{\hat{\theta}^p (e^\alpha - 1)(1 + e^\alpha)}{\left((1 + e^\alpha)\overline{Z}_n - 1\right)\left((1 + e^\alpha)\widetilde{p}_{\theta_0}(1) - 1\right)} \left( \overline{Z}_n - \widetilde{p}_{\theta_0}(1) \right).
\]
Recalling that we are in the case where \( \hat{\theta}^p \leq \theta_0 \), which implies
\[
(1 + e^\alpha)\widetilde{p}_{\theta_0}(1) - 1 = (e^\alpha - 1) \frac{\hat{\theta}^p}{\theta_0},
\]
we directly obtain the desired result. Applying the Central Limit Theorem, we have:
\[
\sqrt{n} \left( \overline{Z}_n - \widetilde{p}_{\theta_0}(1) \right) \xrightarrow[\mathcal{L}]{n \to \infty} \mathcal{N} \left( 0, \widetilde{p}_{\theta_0}(1)\left(1 - \widetilde{p}_{\theta_0}(1)\right) \right).
\]
After some straightforward computations and substituting \( \widetilde{p}_{\theta_0}(1) \) as in Equation~\eqref{eq : tilde p unif estimateur}, this leads to the desired result regarding the asymptotic distribution of \( \widehat{\theta}_n \).
 \end{proof}
\begin{rem} 	
The estimation procedure is consistent as long as the preliminary estimate $\widehat{\theta}^{\text{p}}$ is smaller than the true value of the parameter. The variance of the estimator is reduced if $\widehat{\theta}^{\text{p}}$ is close to the true value.  

When $\alpha \to 0$, we have
\[
v\bigl(\theta_0,\widehat{\theta}^{\text{p}}\bigr) \sim 
\frac{\theta_0^4}{\alpha^2 \bigl(\widehat{\theta}^{\text{p}}\bigr)^2}.
\]
In particular, if $\widehat{\theta}^{\text{p}} = \theta_0$, the variance of the estimator is, for $\alpha \to 0$, equivalent to 
\[
\frac{\theta_0^2}{\alpha^2},
\]
which is exactly the best achievable variance according to Proposition~\ref{prop : Fisher unif}. In practice, it is impossible to choose $\widehat{\theta}^{\text{p}} = \theta_0$, but it is possible to replace it with some preliminary consistent estimator based on a small subset of the data.
\end{rem}

To see how well the estimation procedure works in practice, we perform a Monte Carlo experiment for the estimation of $\theta_0=1$, with a sample size of $n=10^3$. The value of $\widehat{\theta}^{\text{p}}$ ranges from $0.5* \theta_0$ to $1.3*\theta_0$. 
In Figures \ref{fig : mean}--\ref{fig : std}, we draw the value of the empirical mean and the empirical standard deviation as a function of $\widehat{\theta}^{\text{p}}$. The dashed line in
Figure \ref{fig : std} is the lower bound of the standard deviation as given by the assessment on the Fisher information of Proposition \ref{prop : Fisher unif}. The privacy level is $\alpha=0.3$ and the number of Monte Carlo iterations is $10^6$.

If the value of $\widehat{\theta}^{\text{p}}$ is too small, the standard deviation of the estimator is rather high (see Figure \ref{fig : std}). This standard deviation decreases when $\widehat{\theta}^{\text{p}}$ increases to the true value of the parameter. 
 
In practice, we see that if $\widehat{\theta}^{\text{p}} \in [0.7*\theta_0, \theta_0]$ the estimator works well and is close to the efficiency.
% We see that the estimator has small bias and standard deviation and is, in practice, very close to the efficiency if $\widehat{\theta}^{\text{guess}}$ is close to and smaller than $\theta_0$.
 If $\widehat{\theta}^{\text{p}} > \theta_0$, the estimator is completely biased as shown in Figure \ref{fig : mean}.

\begin{figure}
	\centering
	\begin{minipage}{0.50\textwidth}
		\centering
		\includegraphics[width=0.9\textwidth]{Mean} % first figure itself
		\caption{Empirical Mean \label{fig : mean}}
	\end{minipage}\hfill
	\begin{minipage}{0.50\textwidth}
		\centering
		\includegraphics[width=0.9\textwidth]{Std} % second figure itself
		\caption{Empirical Std \label{fig : std}}
	\end{minipage}
\end{figure}



\section{Functional analysis tools about Choquet's theorem}{\label{S: tools}}
To establish our main results in a continuous space \({X}\), our key tool is Choquet’s theorem, as stated in Theorem \ref{th: Choquet}. In order to apply it within our framework, a solid understanding of certain concepts in functional analysis is required. The purpose of this section is therefore to introduce the necessary background, with a particular focus on the weak-\(\star\) topology.


\subsection{Weak -$\star$ topology}
Let \({X}\) be a Banach space, and let \({X}'\) denote its topological dual—that is, the space of all continuous linear functionals on \({X}\). This space is itself a Banach space, equipped with the norm
\[
| | | f | | | := \sup_{x \neq 0} \frac{\langle f, x \rangle}{\| x \|},
\]
where \(\langle f, x \rangle\) denotes the dual pairing, meaning the image of \(x\) under the linear functional \(f\).

The weak topology on \({X}\) is defined as the coarsest topology that makes all functionals \(f \in {X}'\) continuous. In finite-dimensional settings, the weak and strong topologies coincide, but this is generally not the case in infinite dimensions. The weak topology becomes particularly useful when the dual space \({X}'\) is well understood—for example, in the case of \(L^p(\cX)\) spaces, where \(\cX\) is an open subset of \(\mathbb{R}^d\) equipped with the Lebesgue measure and \(p \in [1, \infty)\).

In this setting, the topological dual of \(L^p(\cX)\) can be identified with \(L^q(\cX)\), where \(q\) is the conjugate exponent of \(p\), i.e., \(\frac{1}{p} + \frac{1}{q} = 1\). Thanks to the Riesz representation theorem (see \cite{Bre83}, p.61), for any \(f \in (L^p(\cX))'\), there exists a unique \(v \in L^q(\cX)\) such that
\[
\langle f, u \rangle = \int_\cX v u, \quad \text{for all } u \in L^p(\cX).
\]

Given a Banach space \({X}\) and its topological dual \({X}'\), one can similarly define a weak topology on \({X}'\). However, we are now interested in introducing a topology on \({X}'\) that is even finer—strictly so when \({X}\) is not reflexive.


\begin{df}{\label{def: weakstar top}}
The weak-\(\star\) topology on \({X}'\) is the coarsest topology that makes all the maps  
\begin{align*}
J_x : \, \, & {X}' \rightarrow \R \\
& f \mapsto \langle f, x \rangle
\end{align*} 
continuous for every \(x \in {X}\).  
\end{df}

As a direct consequence of its definition, the weak-$\star$ topology satisfies the Hausdorff condition. A sequence \((f_n) \subset {X}'\) is said to converge to \(f \in {X}'\) in the weak-$\star$ topology, denoted by
\[
f_n \xrightharpoonup{\star} f \quad \Leftrightarrow \quad \forall x \in {X}, \ \langle f_n, x \rangle \to \langle f, x \rangle.
\]
It is important to observe that strong convergence implies weak convergence, which in turn implies weak-$\star$ convergence. Moreover, by the Banach–Steinhaus theorem, any sequence that converges in the weak-$\star$ topology is necessarily bounded (see p.16 of \cite{Bre83}).

Weak-$\star$ closed sets are, by definition, the complements of weak-$\star$ open sets. However, it is worth noting that sets which are closed in the strong topology — even convex ones — are not necessarily closed in the weak-$\star$ topology. For instance, if \({X}\) is not reflexive and \(\xi \in {X}'' \setminus J({X})\), where $J$ is the canonic injection, then the hyperplane \(\xi^\perp\), which is closed in both the strong and weak topologies, is not weak-$\star$ closed (see p.42 of \cite{Bre83}).

Nevertheless, closed balls in \({X}'\) are weak-$\star$ closed, and even compact in the weak-$\star$ topology (see pp.42–43 of \cite{Bre83}). This property will play a key role in our subsequent analysis. 

\begin{thm}[Banach-Alaoglu]{\label{th: BanAla}}
 The closed unit ball of the dual space of a normed vector space is compact in the weak-$\star$ topology. 
\end{thm}

Note that the classical notion of compactness, defined via the Borel–Lebesgue property, is equivalent to the Bolzano–Weierstrass property (i.e., every bounded sequence admits a convergent subsequence) only in metric spaces. Now, the dual space \({X}'\), equipped with the weak-$\star$ topology, is not metrizable if \({X}\) is infinite-dimensional. However, it is important to highlight that the closed unit ball of \({X}'\) is metrizable whenever \({X}\) is separable—that is, it admits a countable dense subset.

In that case, one can define a compatible metric on the closed unit ball of \({X}'\) by
\begin{equation}{\label{eq: d ball}}
 d(f, g) := \sum_{n = 1}^\infty 2^{-n} \left|\langle f - g, x_n \rangle\right|,   
\end{equation}
where \((x_n)_{n \in \mathbb{N}}\) is a countable dense subset of the closed unit ball of \({X}\) (see pp.48–49 of \cite{Bre83}).

The Banach–Alaoglu theorem is particularly useful in the case where \({X}\) is separable, since it implies that every bounded sequence in \({X}'\) admits a weak-$\star$ convergent subsequence. This applies, for instance, to the case \({X} = L^1(\cX)\), in which any bounded sequence in \(L^\infty(\cX)\) has a subsequence that converges in the weak-$\star$ sense.


\subsection{About Choquet's theorem}
It is well known that the two endpoints of a line segment determine all the points in between: in vector terms, the segment from \( v \) to \( w \) consists of all points of the form \( \lambda v + (1 - \lambda) w \) with \( 0 \leq \lambda \leq 1 \). A classical result by Hermann Minkowski states that in a Euclidean space, any bounded, closed, and convex set \( C \) is the convex hull of its set of extreme points \( \mathcal{E} \). This means that every point \( c \in C \) can be written as a finite convex combination of points \( e \in \mathcal{E} \). The set \( \mathcal{E} \) may be finite or infinite.

In vector terms, assigning non-negative weights \( w(e) \) to elements \( e \in \mathcal{E} \), with almost all weights being zero, any point \( c \in C \) can be represented as
\[
c = \sum_{e \in \mathcal{E}} w(e) e, \qquad \sum_{e \in \mathcal{E}} w(e) = 1,
\]
so that the weights \( w(e) \) define a probability measure supported on a finite subset of \( \mathcal{E} \). For any affine function \( f \) defined on \( C \), its value at \( c \) can then be expressed as
\[
f(c) = \int_{\mathcal{E}} f(e) \, dw(e).
\]

In infinite-dimensional settings, one aims to obtain a similar representation. This is made possible by Choquet's theorem (see p.337 of \cite{Choquet}), which we state below and will be a key tool in our approach.

\begin{thm}[Choquet]{\label{th: Choquet}}
Let \( C \) be a convex, compact, and metrizable subset of a locally convex Hausdorff topological vector space \( U \). Then:
\begin{itemize}
    \item The set \( \mathcal{E} \subset C \) of extreme points of \( C \) is a \( G_\delta \)-set (i.e., a countable intersection of open sets).
    \item Every point \( b \in C \) is the barycenter of at least one non-negative Radon measure supported on \( \mathcal{E} \); that is, there exists a Radon probability measure \( \mu \) on \( \mathcal{E} \) such that
  \begin{equation} \label{eq : Choquet representation}
  b = \int_{\mathcal{E}} r \, d\mu(r).
  	\end{equation}
  
\end{itemize}
\end{thm}

In the following, we will apply Choquet's theorem with \( U = L^\infty(\mathcal{X}) \), where \( \mathcal{X} \) denotes the space from which the privacy mechanisms originate and is an open subset of $\R^d$. %to simplify, but the results extend to a more general case (see Remark \ref{rk: extension functional space} below for details). Recall that \( L^\infty(\mathcal{X}) \) is the topological dual of \( L^1(\mathcal{X}) \), and when equipped with the weak-\(\star\) topology, it satisfies the assumptions of Choquet's theorem, since all normed vector spaces are locally convex.

As anticipated in Section \ref{ss: infinite}, in our setting it will be essential to apply Choquet’s theorem to the set
\begin{equation}{\label{eq: def B}}
  \theconvex = \left\{ v : \mathcal{X} \rightarrow \mathbb{R} \, \text{ such that } \, 1 \leq v(x) \leq e^\alpha \text{ a.e.} \right\}. 
\end{equation}
To this end, we must verify that the assumptions of Choquet's theorem are fulfilled for this subset of \( L^\infty(\mathcal{X}) \). Note that \( \theconvex \) can be interpreted as a truncated closed ball of radius \( e^\alpha \), with the unit ball removed. As previously mentioned, although the dual space is not metrizable in general if infinite-dimensional, one can construct a compatible metric on its closed unit ball using the formula in \eqref{eq: d ball}. We begin by showing that a similar construction can be applied to the set \( \theconvex \).

We introduce the notation $\widehat{B}$ for the closed ball, 
\begin{equation}\label{eq: def Ball L infty}
\widehat{B}=B_{\|\cdot\|_\infty}(0, e^{\alpha}) =\left\{ v : \mathcal{X} \rightarrow \mathbb{R} \, \text{ such that } \, \abs{ v(x)} \leq e^\alpha \text{ a.e.} \right\}.
\end{equation}

\begin{lem}{\label{l: B metrizable}}
The set $\theconvex$ defined in \eqref{eq: def B} equipped with the weak-$\star$ topology is metrizable by some metric denoted $d_\star$.  
\end{lem}
\begin{proof}
Recall that \( L^1(\mathcal{X}) \) is a separable space, meaning that it admits a countable dense subset. Let us denote such a sequence by \( (x_n)_{n \in \mathbb{N}} \). Using this sequence, one can construct a compatible metric on the closed ball of radius \( e^\alpha \) in \( L^\infty(\mathcal{X}) \). Analogously to the distance \( d \) introduced in \eqref{eq: d ball} for the closed unit ball, we define:
\begin{equation}{\label{eq: def dstar}}
d_\star(f, g) := \sum_{n = 1}^\infty 2^{-n} \left|\langle f - g, x_n \rangle_{L^1, L^\infty}\right|.   
\end{equation}
This distance metrizes the closed ball $\widehat{B}$,
%of radius \( e^\alpha \), denoted by \( B_{\|\cdot\|_\infty}(0, e^{\alpha}) \), 
with respect to the weak-\(\star\) topology. Since the set \( \theconvex \) is a subset of $\widehat{B}$, % \( B_{\|\cdot\|_\infty}(0, e^{\alpha}) \), 
it is metrizable as well.
\end{proof}

We now turn to the proof of the compactness of \( C \). In Lemma \ref{l: B closed}, whose proof can be found in Section \ref{s: proof technical}, we show that \( \theconvex \) is closed with respect to the weak-\(\star\) topology. The compactness of \( \theconvex \) then follows from compactness of $\widehat{B}$ by the Banach–Alaoglu Theorem, as stated in Theorem \ref{th: BanAla}.

\begin{lem}{\label{l: B closed}}
The set $\theconvex$ defined in \eqref{eq: def B} is a compact with respect to the weak-$\star$ topology.  
\end{lem}

\begin{rem}\label{rem : E not closed}
We show that the set $\theconvex$ is closed with respect to the weak-$\star$ topology. On the other hand, in general, we cannot expect the set $\cE$ to be closed. For instance, take $\cX = [0,1]$, and for $n \ge 1$, define  
\[
r_n = \sum_{k=0}^{n-1} \left( \indi{\left[\frac{2k}{2n}, \frac{2k+1}{2n}\right)} + e^\alpha \, \indi{\left[\frac{2k+1}{2n}, \frac{2k+2}{2n}\right)} \right).
\]
The sequence $(r_n)_n$ takes values in $\cE$. However, it can be shown that $r_n \xrightharpoonup{\star} v$, where $v(x) = \frac{1 + e^\alpha}{2}$ for all $x \in \cX$. Since $v \notin \cE$, we conclude that the set $\cE$ is not closed.

\end{rem}

The set \( C \), as defined in Equation \eqref{eq: def B}, is also convex. Therefore, by Lemmas \ref{l: B metrizable} and \ref{l: B closed}, we are in a position to apply Choquet’s theorem. In particular, this implies that any \( b \in C \subset L^\infty(\mathcal{X}) \) can be represented as  
\[
b = \int_{\mathcal{E}} r \, d\mu(r),
\]  
for some probability measure \( \mu \) supported on the extreme points of \( C \). This representation is an equality between functions of $L^\infty(\cX)$. 
Our next objective is to give a rigorous meaning to the pointwise representation  
\[
b(x) = \int_{\mathcal{E}} r(x) \, d\mu(r), \quad \text{for } x \in \mathcal{X},
\]  
which will be the focus of the upcoming subsection. Before doing so, let us provide extra details on the extremal points of $\theconvex$, as in the following lemma.

\begin{lem}\label{lem : extremal points}
The element	$r \in C$ is extremal iff $r(x)\in \{1,e^\alpha\}$ $dx$ almost everywhere. 
\end{lem}
\begin{proof}
We recall that $r$ is extremal in $\theconvex$ if any decomposition $r = \frac{r_1}{2} + \frac{r_2}{2}$ with $r_1$ and $r_2$ in $\theconvex$ implies $r_1 = r_2 = r$ a.e.. It is then clear that if $r(x) \in \{1, e^\alpha\}$ $dx$-a.e., then $r$ is extremal.  

Conversely, let us show that if $r$ is extremal, it takes a.e. only the two values $1$ and $e^\alpha$. By contradiction, if this is not the case, then there exists some $\varepsilon > 0$ and $\mathcal{N} \subset \mathcal{X}$ such that for $x \in \mathcal{N}$, $r(x) \in [1+\varepsilon, e^\alpha - \varepsilon]$ and $\mathcal{N}$ has positive Lebesgue measure. Then we can split $r$ as
\[
r = \frac{1}{2} \bigl( r \indi{\mathcal{N}^c} + (r+\varepsilon) \indi{\mathcal{N}} \bigr)
+ \frac{1}{2} \bigl( r \indi{\mathcal{N}^c} + (r-\varepsilon) \indi{\mathcal{N}} \bigr)
= \frac{r_1}{2} + \frac{r_2}{2}.
\]
This shows that $r$ is not extremal, arriving at a contradiction.
\end{proof}





\subsection{Point evaluation operator and Bochner integral}{\label{ss:point eval}}
The aim of this subsection is to give a rigorous interpretation to the pointwise representation  
\[
b(x) = \int_{\mathcal{E}} r(x) \, d\mu(r)
\quad \text{for almost every } x \in \mathcal{X} = \R^d.
\]  
To achieve this, the notion of the Bochner integral, introduced in Definition \ref{def: Bocnher} below, will play a key role. 
The measure $\mu$ is supported on $\cE\subset \widehat{B}\subset L^\infty(\cX)$, and the RHS of \eqref{eq : Choquet representation} is an integral of a function valued in the Banach space $U=L^\infty(\cX)$ (see \cite{DieUhl_Book}, \cite{Lang_Book}  for presentation of Banach valued integrals). In particular, the integral can be seen as a Dunford integral on this Banach space (see Lemma 1 p. 52 of \cite{DieUhl_Book}). However, we will precise that this integral is also a Bochner integral for the $d_\star$ distance, meaning that an approximation as given in Definition \ref{def: Bocnher} below is valid. We will deduce a rigorous statement for the pointwise representation.

%We begin by discussing the Bochner integral and the point evaluation operator in a general setting, before showing how this framework applies specifically to our set \( C \).



\begin{df}{\label{def: Bocnher}}
	Let \( (\Omega, \mathcal{F}, \mu) \) be a measure space, where $\mu$ is a sub-probability and let 
	$f : (\Omega, \mathcal{F}) \mapsto (\widehat{B},\cB(\widehat{B}))$ be a measurable map, where $\cB(\widehat{B})$ is the Borel $\sigma$-field of $(\widehat{B}, d_\star)$.
	%	\( Y \) be a Banach space. A measurable function \( f: X \rightarrow Y \) 
	The map $f$ is said to be Bochner integrable if there exists a sequence of integrable simple functions \( (s_n) \) such that
	\[
	\lim_{n \to \infty} \int_\Omega d_\star(f,s_n) d\mu = 0,
	\]
	where the integral on the left-hand side is the usual Lebesgue integral. In this case, the Bochner integral of \( f \) is defined as
	\[
	\int_\Omega f \, d\mu := \lim_{n \to \infty} \int_\Omega s_n \, d\mu,
	\]
	where the convergence holds with the $d_\star$ distance.
\end{df}
%
%\begin{df}{\label{def: Bocnher}}
%Let \( (X, \mathcal{F}, \mu) \) be a measure space, and let \( Y \) be a Banach space. A measurable function \( f: X \rightarrow Y \) is said to be Bochner integrable if there exists a sequence of integrable simple functions \( (s_n) \) such that
%\[
%\lim_{n \to \infty} \int_X \|f - s_n\|_Y \, d\mu = 0,
%\]
%where the integral on the left-hand side is the usual Lebesgue integral. In this case, the Bochner integral of \( f \) is defined as
%\[
%\int_X f \, d\mu := \lim_{n \to \infty} \int_X s_n \, d\mu.
%\]
%\end{df}
%

It can be shown that the sequence \( \left\{ \int_{\Omega} s_n \, d\mu \right\}_{n=1}^\infty \) is a Cauchy sequence for $d_\star$ distance %in the Banach space \( Y \), 
and hence the limit exists (see Remark \ref{rk: Cauchy} below). Moreover, this limit is independent of the choice of approximating simple functions \( \{s_n\}_{n=1}^\infty \), which ensures that the Bochner integral is well defined. 

%It can be shown that the sequence \( \left\{ \int_X s_n \, d\mu \right\}_{n=1}^\infty \) is a Cauchy sequence in the Banach space \( Y \), and hence the limit exists. Moreover, this limit is independent of the choice of approximating simple functions \( \{s_n\}_{n=1}^\infty \), which ensures that the Bochner integral is well defined. 

An important property of Bochner integral is that, if \( T: U \rightarrow \mathbb{R} \) is a continuous linear operator
% between Banach spaces \( Y \) and \( Z \), and if \( f: X \rightarrow Y \) is Bochner integrable, 
then \( T \circ f: \Omega \rightarrow Z \) is also Bochner integrable. In this case, integration and the application of \( T \) can be interchanged:
\begin{equation}{\label{eq: prop Bochner}}
	\int_{\Omega} T f \, d\mu = T \left( \int_{\Omega} f \, d\mu \right).   
\end{equation}
%for all measurable subsets \( \mathcal{E} \in \mathcal{F} \).\\
\\
%
%An important property of Bochner integral is that, if \( T: Y \rightarrow Z \) is a continuous linear operator between Banach spaces \( Y \) and \( Z \), and if \( f: X \rightarrow Y \) is Bochner integrable, then \( T \circ f: X \rightarrow Z \) is also Bochner integrable. In this case, integration and the application of \( T \) can be interchanged:
%\begin{equation}{\label{eq: prop Bochner}}
% \int_{\mathcal{E}} T f \, d\mu = T \left( \int_{\mathcal{E}} f \, d\mu \right)   
%\end{equation}
%for all measurable subsets \( \mathcal{E} \in \mathcal{F} \).\\
%\\
We would like to apply the result above to the evaluation operator at a fixed point $x\in \cX$, $f\in U \mapsto f(x)\in \mathbb{R}$. However, such operator is not continuous for the $d_\star$ distance. 

We formally introduce %on a general space $Y$ 
as follows the evaluation operator. For any \( x \in \cX \), consider the map
\begin{align}\nonumber %\label{eq: evaluation}
e_x : \quad & \widehat{B} \rightarrow \mathbb{R} \\
\label{eq : def evaluation operator}
& b \mapsto \limsup_{\epsilon \to 0} \frac{1}{|B(x, \epsilon)|} \int_{B(x, \epsilon)\cap\cX} b(y) \, dy,
\end{align}
where \( B(x, \epsilon) \) denotes the ball of radius \( \epsilon \) centered at \( x \), and \( |B(x, \epsilon)| := \int_{{\R}^d} \indi{B(x, \epsilon)}(y) \, dy \) is its Lebesgue measure. 

Note that \( e_x(b) = b(x) \) whenever \( x\in\cX \) is a Lebesgue point of the function \( b \), and consequently for almost every \( x \) with respect to the Lebesgue measure. The reason why the evaluation is defined though the limit \eqref{eq : def evaluation operator} is that it enables to prove the joint measurability $(x,b) \mapsto e_x(b)$ in Proposition \ref{p: evaluation} below. We will also justify that the Bochner integral property stated in \eqref{eq: prop Bochner} applies even without continuity of the evaluation operator. It leads in particular to
%  is a Lebesgue point of \( y \).  %, and that if \( y \) is essentially bounded, then almost every \( x \) (with respect to the Lebesgue measure) is a Lebesgue point of \( y \). The operator \( e_x \) is introduced in order to ensure bi-measurability. One can verify that this operator is linear. If it were also continuous, then—by virtue of the Bochner integral property stated in \eqref{eq: prop Bochner}—we would directly obtain, for any \( y \in Y \) admitting a representation of the form \( y = \int_Y f \, d\mu(f) \), and for almost every $x \in \R^d$,
\[
b(x) = e_x(b) = e_x\left( \int_\cE r \, d\mu(r) \right) = \int_\cE e_x(r) \, d\mu(r), ~ dx-a.e.,
\]
which is precisely the rigorous pointwise evaluation we aim to justify.


%However, a key difficulty arises from the fact that the evaluation operator \( e_x \) is not continuous with respect to the metric \( d_\star \) defined in \eqref{eq: def dstar}. As discussed in Lemma \ref{l: B metrizable}, this metric is crucial because it metrizes the weak-\(\star\) topology on \( \theconvex \). Nevertheless, even in the absence of continuity of \( e_x \), it is still possible to recover the desired pointwise representation of the Bochner integral, provided we have additional structural information about the space under consideration, as we will see at the end of this subsection for our set $\theconvex$. \\

To establish the desired pointwise evaluation, the first step is to show that the integral over the extreme points, as given by Choquet’s theorem, is in fact a Bochner integral. To this end, it is essential to verify that \( \widehat{B} \) is separable, as stated in the following lemma. Its proof can be found in Section \ref{s: proof technical}.



\begin{lem}{\label{l: B separable}}
The set $\widehat{B}$ defined in \eqref{eq: def Ball L infty} equipped with the weak-$\star$ topology is separable.  
\end{lem}


The separability of \( \widehat{B} \), as established above, plays a crucial role in ensuring that integrals over \( \widehat{B} \) are Bochner integrals. In particular, this means—according to Definition \ref{def: Bocnher}—that they can be approximated by integrals of simple functions.



\begin{lem}{\label{l: Bochner}}
	The set $\widehat{B}$ be the set defined in \eqref{eq: def Ball L infty}
%Let $B$ be the set defined in \eqref{eq: def B}, 
equipped with the weak-$\star$ topology, and corresponding Borel $\sigma$-field. Let \( \mu \) be a sub-probability measure on \( \widehat{B} \). Then,  \( \int_{\widehat{B}} f  d\mu(f) \) is a Bochner integral.
\end{lem}


The separability and metrizability of the space \( \widehat{B} \), as established
above
 %in Lemma \ref{l: B metrizable}, 
enable us to obtain the desired pointwise evaluation, as stated in Proposition \ref{p: evaluation} below. Its proof can be consulted in Section \ref{s: proof technical}.

\begin{prop}{\label{p: evaluation}}
	1) The map 
	\begin{align*}
		&\widehat{B}\times \cX \to \R
		\\
		&(f,x) \mapsto e_x(f)
	\end{align*}
	is measurable with respect to the $\sigma$-field $\cB(\widehat{B}) \otimes \cB(\cX)$.
	
	2) Let $\mu$ be a sub-probability on the space $\widehat{B}$ and assume that $b=\int_{\widehat{B}} f \mu(df)$. Then,
	%	Let $B$ be the set in Equation \ref{eq: def B}. Let us consider $\mathcal{E} \subset B$ the set of extreme points of $B$. Then, for any $b \in B$ and 
	for almost every $x \in \cX$, 
	\begin{equation}\label{eq : evalutation rigorous}
		b(x) =e_x(b) =\int_{\widehat{B}} e_x(f) d\mu (f). 
	\end{equation}
	3) The $\mu \otimes dx$ measure of the complement of the following set is zero:
	\begin{equation} \label{eq : neg by Fubini}
		\big\{(f,x)\in\widehat{B}\times\cX : e_x(f)  = \lim_{\epsilon \to 0} \frac{1}{|B(x, \epsilon)|} \int_{B(x, \epsilon)\cap\cX} f(y) \, dy\big\}.
	\end{equation}
	
	
\end{prop}


%\begin{prop}{\label{p: evaluation}}
%Let $B$ be the set in Equation \ref{eq: def B}. Let us consider $\mathcal{E} \subset B$ the set of extreme points of $B$. Then, for any $b \in B$ and almost every $x \in \R^d$, 
%$$b(x) = \int_{\mathcal{E}} f(x) d\mu (f).$$
%\end{prop}

\begin{rem}{\label{rk: extension functional space}}
It would be both natural and useful for applications to generalize the results stated for $\mathcal{X} \subset \mathbb{R}^d$ to a broader framework where $\mathcal{X}$ is a topologically complete separable metric space, equipped with a nonatomic, normalized Borel measure $\mu$. This extension is especially relevant for problems involving the statistical analysis of stochastic processes driven by Brownian motion under local differential privacy constraints.

While we leave this question for future work, we conjecture that our main results extend to this more general setting. A key technical ingredient in our proofs involving the evaluation operator is the property that almost every point in $\mathcal{X}$ is a Lebesgue point for functions in a suitable class $B \subset \mathcal{X}$. This is straightforward when $\mathcal{X} = \mathbb{R}^d$ (see Step 1 in the proof of Proposition~\ref{p: evaluation}), but is less immediate in an abstract measurable space.

Fortunately, Theorem 2 in \cite{Oxt} ensures the existence of a homeomorphism between a $G_\delta$ subset of $\mathcal{X}$ (with full $\mu$-measure) and the interval $(0,1)$ with the Lebesgue measure. This suggests that the analysis of the evaluation operator, and in particular of Lebesgue points, can be transferred to $(0,1)$ via this homeomorphism, since it is continuous in both directions. 

We believe this intuition supports the conjectured extension, though several technical challenges remain. A detailed treatment of these issues lies beyond the scope of the present work. %As mentioned, such an extension would be particularly valuable for studying statistical inference for Brownian-driven processes in the Wiener space under privacy constraints, and we intend to pursue this direction in future work.

\end{rem}









\section{Proof of main results}{\label{s: proof main}}

This section is dedicated to the proofs of our main results, specifically the bounds on the Fisher information in both the discrete and continuous settings, as presented in Sections \ref{ss: finite} and \ref{ss: infinite}, respectively.

\subsection{Discrete case: proof of Theorem \ref{th : info cas fini}}

Let us begin by introducing some notation that will be useful throughout the proof.
For each $x \in \mathcal{X}$, define the index sets
\begin{equation}{\label{eq: set Qx+-}}
Q^+_x = \{ \beta \in \{0, \dots, 2^d - 1\} \mid r_\beta(x) = e^\alpha \}, \quad 
Q^-_x = \{ \beta \in \{0, \dots, 2^d - 1\} \mid r_\beta(x) = 1 \}.
\end{equation}
With this notation, we can express the mechanism $q^{(1)}$ as
\begin{equation}{\label{eq: mechanism q1}}
q^{(1)}_x(\beta) = \omega_\beta r_\beta(x) = \omega_\beta \left( e^\alpha \indi{Q^+_x}(\beta) + \indi{Q^-_x}(\beta) \right).
\end{equation}

We are now ready to proceed with the proof of Theorem \ref{th : info cas fini}. The result follows directly from the upper and lower bounds stated in Propositions \ref{prop : borne sup Fisher fini} and \ref{prop : borne inf Fisher fini}, respectively.

We begin with the more involved part—establishing the upper bound—as detailed in Proposition \ref{prop : borne sup Fisher fini}.



\begin{prop}\label{prop : borne sup Fisher fini}
	Assume that $\Theta \subset \R$ and that $(p_{\theta})_\theta$ 
	is DQM at $\theta_0 \in \overset{\circ}{\Theta}$ with score function $s_{\theta_0}$. Let $q : \cX \to \cZ$ be any $\alpha$-LDP mechanism taking values in a finite set $\cZ$
	and $(\widetilde{p}_\theta)_\theta$ the model for the public variable.
	 Assume that $s_{\theta_0}(x)\neq 0$ for all $x\in \cX$.
	Then, there exists $\overline{\alpha}$ such that if $\alpha \le \overline{\alpha}$,
	\begin{equation} \label{eq :  uppper bound Fisher finite exacte}
		\cI_{\theta_0}((\widetilde{p}_\theta)_\theta)=\E\left[\abs*{t_{\theta_0}(Z)}^2\right] \le 
		\cI^{\text{max},\alpha}_{\theta_0}((p_\theta)_\theta),
	\end{equation}
	where we recall that
	\begin{equation*}
\cI^{\text{max},\alpha}_{\theta_0}((p_\theta)_\theta)
=\frac{(e^\alpha-1)^2}{4}% 
\times \frac{\E\left[\abs{s_{\theta_0}(X)}\right]^2}{\left[ (1-n_\text{max})+e^\alpha n_\text{max}  \right]
\left[n_\text{max}+(1-n_\text{max})e^\alpha  \right]}
	\end{equation*}
	and $n_\text{max}=\int_{x: s_{\theta_0}(x)>0} p_{\theta_0}(x) \mu(dx)$.
	  The value of $\overline{\alpha}$ is independent of the channel $q$, but depends on the model $(p_{\theta})_\theta$.
\end{prop}	

\begin{proof}
As \( q \in \cQ_\alpha \), we use the factorisation Lemma~\ref{lem : factorisation cas fini} to write \( q = q^{(2)} \circ q^{(1)} \), where \( q^{(1)} \) is an extremal \(\alpha\)-LDP channel described in the statement of Lemma~\ref{lem : factorisation cas fini}. From \cite{steinbergerEfficiencyLocalDifferent}, and as recalled in~\eqref{eq : decroissance Fisher}, we have
\[
\cI_{\theta_0}\big((\widetilde{p}_\theta)_\theta\big)
= \cI_{\theta_0}(q \circ \cP)
= \cI_{\theta_0}\big((q^{(2)} \circ q^{(1)}) \circ \cP\big)
= \cI_{\theta_0}\big(q^{(2)} \circ (q^{(1)} \circ \cP)\big)
\le \cI_{\theta_0}(q^{(1)} \circ \cP).
\]
Hence, it is sufficient to prove~\eqref{eq :  uppper bound Fisher finite exacte} for \( q = q^{(1)} \). To lighten notation, we will denote this quantity as \( q \) in place of \( q^{(1)} \) in the remainder of the proof. Accordingly, we consider the extremal mechanism \( q : \cX \to \cE = \{0,\dots,2^d-1\} \), with probability coefficients
\begin{equation} \label{eq: qx beta 7.25}
q_x(\beta) = \omega_\beta r_\beta(x),
\end{equation}
as in Lemma~\ref{lem : factorisation cas fini}. Recall from~\eqref{eq: mass omega beta} and~\eqref{eq: r beta x} that
\[
\sum_{\beta=0}^{2^d-1} \omega_\beta \in [e^{-\alpha}, 1],
\quad \text{and} \quad
r_\beta(x) = e^{\alpha} \indi{F^+_\beta}(x) + \indi{F^-_\beta}(x) = 1 + (e^\alpha - 1)\indi{F^+_\beta}(x).
\]
In particular, the expression for the probability \( \widetilde{p}_\theta \) on \( \cE \) is given by
\[
\widetilde{p}_\theta(\beta)
= \int_\cX p_\theta(x) q_x(\beta) \, \mu(dx)
= \int_\cX p_\theta(x) \omega_\beta r_\beta(x) \, \mu(dx)
= \omega_\beta \int_\cX p_\theta(x) \left[1 + (e^\alpha - 1)\indi{F^+_\beta}(x)\right] \mu(dx),
\]
where in the last equality we replaced \( r_\beta(x) \) using~\eqref{eq: r beta x}. Since \( p_\theta \) is a probability density, we have \( \int_\cX p_\theta(x) \mu(dx) = 1 \), so that
\begin{equation} \label{eq: ptilde 7.5}
\widetilde{p}_\theta(\beta) = \omega_\beta \left(1 + (e^\alpha - 1) n^+_\beta\right),
\end{equation}
with
\begin{equation} \label{eq : def nplus}
n^+_\beta = \int_{F^+_\beta} p_\theta(x) \mu(dx).
\end{equation}

Recall that, according to~\cite{steinbergerEfficiencyLocalDifferent} and as stated in~\eqref{eq : score public}, the score function of the model \( q \circ \cP \) is given by
\[
t_\theta(\beta) = \frac{\int_\cX s_\theta(x) q_x(\beta) p_\theta(x) \mu(dx)}{\widetilde{p}_\theta(\beta)}.
\]
We already obtained an expression for the denominator \( \widetilde{p}_\theta(\beta) \) in~\eqref{eq: ptilde 7.5}. For the numerator, using~\eqref{eq: qx beta 7.25} and the definition of \( r_\beta(x) \), we get
\[
\int_\cX s_\theta(x) q_x(\beta) p_\theta(x) \mu(dx)
= \int_\cX s_\theta(x) \omega_\beta r_\beta(x) p_\theta(x) \mu(dx)
= \omega_\beta \int_\cX s_\theta(x) \left[1 + (e^\alpha - 1)\indi{F^+_\beta}(x)\right] p_\theta(x) \mu(dx).
\]
Since the score function \( s_\theta \) of \( \mathcal{P} \) is centered, we have \( \int_\cX s_\theta(x) p_\theta(x) \mu(dx) = 0 \). Thus,
\[
\int_\cX s_\theta(x) q_x(\beta) p_\theta(x) \mu(dx)
= (e^\alpha - 1) \omega_\beta \int_\cX s_\theta(x) \indi{F^+_\beta}(x) p_\theta(x) \mu(dx).
\]
This yields
\begin{equation} \label{eq: t beta 8.5}
t_\theta(\beta) = (e^\alpha - 1)
\frac{\int_\cX s_\theta(x) \indi{F^+_\beta}(x) p_\theta(x) \mu(dx)}{1 + (e^\alpha - 1) n^+_\beta},
\end{equation}
for all \( \beta \) such that \( \omega_\beta > 0 \). Substituting~\eqref{eq: ptilde 7.5} and~\eqref{eq: t beta 8.5}, we obtain the following expression for the Fisher information of \( q \circ \cP \):
\begin{align*}
\cI_\theta(q \circ \cP)
&= \mathbb{E}[t_\theta(Y)^2]
= \int_\mathcal{E} \widetilde{p}_\theta(\beta) \, t_\theta(\beta)^2 \, \nu(d\beta) \\
&= (e^\alpha - 1)^2 \int_\mathcal{E} \omega_\beta
\frac{\left(\int_\cX s_\theta(x) \indi{F^+_\beta}(x) p_\theta(x) \mu(dx)\right)^2}{1 + (e^\alpha - 1) n^+_\beta} \, \nu(d\beta),
\end{align*}
where \( \nu \) is the counting measure on \( \mathcal{E} = \{0, \dots, 2^d - 1\} \). Let us recall
$$i_\beta := \frac{\left(\int_\cX s_\theta(x) \mathbf{1}_{F^+_\beta}(x) p_\theta(x) \mu(dx)\right)^2}{1 + (e^\alpha - 1) n^+_\beta},$$
and consider the maximization problem:
\[
M^* = \sup_{\omega} \int_\mathcal{E} \omega_\beta i_\beta \, \nu(d\beta)
\quad \text{subject to} \quad
\int_\mathcal{E} \omega_\beta r_\beta(x) \, \nu(d\beta) = 1, \quad \forall x \in \cX,
\]
where the supremum is over sub-probability weights \( (\omega_\beta)_{\beta = 0, \dots, 2^d - 1} \) on \( \{0, \dots, 2^d - 1\} \). The constraint arises from~\eqref{eq: constraint factorisation 4.55}. We conclude that
\[
\cI_\theta(q \circ \cP) \le (e^\alpha - 1)^2 M^*.
\]
Thus, the proof of Proposition~\ref{prop : borne sup Fisher fini} is complete once we invoke Lemma~\ref{lem :  determination max fini} below, whose proof can be found in Section \ref{s: proof technical}.

\end{proof}


\begin{lem}\label{lem :  determination max fini}
Under the assumptions of Proposition~\ref{prop : borne sup Fisher fini}, consider the maximization problem
\begin{equation} \label{eq : optimization pb finite}
	M^* = \sup_{\omega \in (\R_+)^{2^d}} \int_{\mathcal{E}} \omega_\beta i_\beta \, \nu(d\beta) \quad
	\text{subject to} \quad \int_{\mathcal{E}} \omega_\beta r_\beta(x) \, \nu(d\beta) = 1, \quad \forall x \in \cX.
\end{equation}
Then, there exists $\overline{\alpha} > 0$ such that for all $\alpha \leq \overline{\alpha}$,
\begin{equation} \label{eq : def M star}
	M^* = \frac{1}{4} \cdot \frac{\E\left[ |s_{\theta_0}(X)| \right]^2}{\left[(1 - n_{\max}) + e^\alpha n_{\max} \right] \left[ n_{\max} + (1 - n_{\max}) e^\alpha \right]},
\end{equation}
where $n_{\max}$ is as defined in~\eqref{eq: def nmax} in the statement of Theorem~1. Moreover, the maximum is attained for a vector $\overline{\omega}$ with only two non-zero components.
\end{lem}
	
This concludes the part of the proof dedicated to the upper bound in Theorem~\ref{th : info cas fini}.  
We now proceed with the proof of the lower bound, as stated in Proposition~\ref{prop : borne inf Fisher fini} below.  
Note that this proof heavily relies on the notation and arguments already introduced in the upper bound analysis.





\begin{prop}\label{prop : borne inf Fisher fini}
	Assume that $\Theta \subset \R$ and that $\cP=(p_{\theta})_\theta$ 
	is DQM at $\theta_0 \in \overset{\circ}{\Theta}$ with score function $s_{\theta_0}$.
	Then for all $\alpha>0$, there exists $q \in \cQ_\alpha$, such that $\cI_{\theta_0}(q\circ \cP)=\cI_{\theta_0}^{\text{max},\alpha}$.
\end{prop}
\begin{proof}
The randomization $q$ achieving the maximum is the one discussed in Remark \ref{rk: choice Fmax} and used in the proof of Lemma \ref{lem :  determination max fini},  where the maximization problem \eqref{eq : optimization pb finite} is solved.
%Proposition~\ref{prop : borne sup Fisher fini}, where the maximization problem \eqref{lem :  determination max fini} is solved. 
For the sake of completeness, we detail it here using simplified notation.

Since the support of the optimal measure $\overline{\omega}$ in Lemma~\ref{lem :  determination max fini} has cardinality $2$, we let $\cZ = \{z_1, z_2\}$. Define
\[
F^+ := \{x \in \cX \mid s_{\theta_0}(x) > 0\}.
\]
We now define the randomization:
\begin{equation}\label{eq: optimal q 49.5}
q_x(z) = P(Z = z \mid X = x) =
\begin{cases}
\displaystyle \frac{1}{1 + e^\alpha} \left( e^\alpha \indi{z_1}(z) + \indi{z_2}(z) \right) & \text{if } x \in F^+, \\
\displaystyle \frac{1}{1 + e^\alpha} \left( \indi{z_1}(z) + e^\alpha \indi{z_2}(z) \right) & \text{if } x \notin F^+.
\end{cases}
\end{equation}
After computation, for any $\theta \in \Theta$, the marginal probability of $Z$ is given by
\begin{align*}
\widetilde{p}_{\theta}(z_1) 
&= P_\theta(Z = z_1) = \int_\cX q_x(z_1) p_\theta(x) \mu(dx) \\
&= \frac{1}{1 + e^\alpha} \int_\mathcal{X} \left( e^\alpha \indi{F^+}(x) + \indi{(F^+)^c}(x) \right) p_\theta(x) \mu(dx) \\
&= \frac{1}{1 + e^\alpha} + \frac{e^\alpha - 1}{1 + e^\alpha} \int_{F^+} p_\theta(x) \mu(dx),
\end{align*}
while $\widetilde{p}_{\theta}(z_2) = 1 - \widetilde{p}_{\theta}(z_1)$. From \eqref{eq : score public} and using the definition \eqref{eq: optimal q 49.5} of $q_x(z)$, the score function of the public data is:
\begin{align*}
t_{\theta_0}(z_1) 
&= \frac{1}{\widetilde{p}_{\theta_0}(z_1)} \int_\cX s_{\theta_0}(x) \, q_x(z_1) \, p_{\theta_0}(x) \mu(dx) \\
&= \frac{1}{\widetilde{p}_{\theta_0}(z_1)} \cdot \frac{1}{1 + e^\alpha} \int_\cX s_{\theta_0}(x) \left( e^\alpha \indi{F^+}(x) + \indi{(F^+)^c}(x) \right) p_{\theta_0}(x) \mu(dx) \\
&= \frac{1}{\widetilde{p}_{\theta_0}(z_1)} \cdot \frac{e^\alpha - 1}{1 + e^\alpha} \int_{F^+} s_{\theta_0}(x) p_{\theta_0}(x) \mu(dx).
\end{align*}
Similarly,
\[
t_{\theta_0}(z_2) = \frac{1}{\widetilde{p}_{\theta_0}(z_2)} \cdot \frac{e^\alpha - 1}{1 + e^\alpha} \int_{(F^+)^c} s_{\theta_0}(x) p_{\theta_0}(x) \mu(dx).
\]
Recalling that 
\[
n_{\max} := \int_{F^+} p_{\theta_0}(x) \mu(dx), \quad 
\int_{F^+} s_{\theta_0}(x) p_{\theta_0}(x) \mu(dx) = \frac{1}{2} \int_\cX \abs{s_{\theta_0}(x)} p_{\theta_0}(x) \mu(dx),
\]
we obtain, after straightforward algebra,
\begin{align}\label{eq: end lower}
\cI_{\theta_0}((\widetilde{p}_\theta)_\theta) 
&= \sum_{i = 1, 2} \widetilde{p}_{\theta_0}(z_i) \, t_{\theta_0}(z_i)^2 \notag \\
&= \frac{(e^\alpha - 1)^2}{(1 + e^\alpha)^2} \cdot \frac{1}{4} \left( \int_\cX \abs{s_{\theta_0}(x)} p_{\theta_0}(x) \mu(dx) \right)^2 \left( \frac{1}{\widetilde{p}_{\theta_0}(z_1)} + \frac{1}{\widetilde{p}_{\theta_0}(z_2)} \right) \notag \\
&= \frac{(e^\alpha - 1)^2}{4} \cdot \frac{\E[\abs{s_{\theta_0}(X)}]^2}{(1 + e^\alpha)^2} \cdot \frac{1}{\widetilde{p}_{\theta_0}(z_1)(1 - \widetilde{p}_{\theta_0}(z_1))}.
\end{align}
Moreover, observe that
\[
\widetilde{p}_{\theta_0}(z_1)(1 - \widetilde{p}_{\theta_0}(z_1)) = 
\frac{1}{1 + e^\alpha} \left( 1 + (e^\alpha - 1) n_{\max} \right)
\left( 1 - \frac{1}{1 + e^\alpha} (1 + (e^\alpha - 1) n_{\max}) \right),
\]
which simplifies to
\[
\frac{1}{(1 + e^\alpha)^2} \left( 1 + (e^\alpha - 1) n_{\max} \right) \left( e^\alpha - (e^\alpha - 1) n_{\max} \right).
\]
Plugging this into the denominator of \eqref{eq: end lower} yields the desired expression \eqref{eq : def I max caf fini}, thus concluding the proof.
\end{proof}	

\subsection{Continuous case: proof of results in Section \ref{ss: infinite}}{\label{ss: proof cont}}

As in the proof of the discrete case, we begin by introducing some notation that will be useful throughout the proofs in this subsection.

 For each $x \in \mathcal{X}$, define the sets
\[
Q^+_x := e_x^{-1}(\{e^\alpha\}) = \{ r \in \mathcal{E} \mid e_x(r) = e^\alpha \}, \quad
Q^-_x := e_x^{-1}(\{1\}) = \{ r \in \mathcal{E} \mid e_x(r) = 1 \}.
\]
Since the map $r \mapsto e_x(r)$ is measurable for each $x$, both $Q^+_x$ and $Q^-_x$ are Borel sets. These sets are the continuous counterparts of those defined in \eqref{eq: set Qx+-} in the previous subsection. With this notation, the kernel defined in \eqref{eq : def random generic extr ctn} can be written as
\[
q^{(\mu)}_x(dr) = \left( e^\alpha \indi{Q^+_x}(r) + \indi{Q^-_x}(r) \right) \mu(dr), \quad \text{for all } x \in \mathcal{X},
\]
and the normalization condition \eqref{eq : cond norm mu continue} becomes
\[
e^\alpha \mu(Q^+_x) + \mu(Q^-_x) = 1 \quad \text{for almost every } x \in \mathcal{X}.
\]
This construction plays the same role as the mechanism $q^{(1)}$ in the discrete case, defined in \eqref{eq: mechanism q1}.

With these definitions in place, we can now proceed to establish the upper and lower bounds stated in Theorem~\ref{thm : equivalent alpha petit ctn}.


 \subsubsection{Proof of Theorem \ref{thm : equivalent alpha petit ctn}}
 \begin{proof}
We begin by proving the upper bound on $\cJ^{\text{max},\alpha}_{\theta_0}$. Let $q \in \cQ_\alpha$. By the factorization result in Proposition~\ref{prop : factor ctn}, any such $q$ can be written as a composition $q = q^{(2)} \circ q^{(\mu)}$, where $\mu$ is a sub-probability measure on $\cE$ satisfying the normalization condition \eqref{eq : cond norm mu continue}, and $q^{(2)}$ is a Markov kernel from $\cE$ to $\cZ$. According to \eqref{eq: fact on Fisher cont}, the Fisher information of the model $q \circ \cP$ is upper bounded by the Fisher information of $q^{(\mu)} \circ \cP$. Therefore, it is sufficient to prove the desired upper bound for $\cI_{\theta_0}(q^{(\mu)} \circ \cP)$. Using \eqref{eq : Fisher ctn general_preli}, we obtain
\begin{equation} \label{eq : majo I in prop alpha small ctn}
	\cI_{\theta_0}(q \circ \cP) \le (e^\alpha - 1)^2
	\int_\cE \frac{\left( \int_{F^+_r} s_{\theta_0}(x) p_{\theta_0}(x) dx \right)^2}
	{\widetilde{p}_{\theta_0}(r)} \mu(dr).
\end{equation}
Now, define the sets $F^+_{\text{max}} := \{ x \in \mathcal{X} \mid s_{\theta_0}(x) > 0 \}$ and $F^{\prime +}_{\text{max}} := \{ x \in \mathcal{X} \mid s_{\theta_0}(x) \le 0 \}$, so that $\mathcal{X} = F^+_{\text{max}} \cup F^{\prime +}_{\text{max}}$. As in the discrete case (see in particular Equation~\eqref{eq : iF smaller iFmax} of the proof of Lemma~\ref{lem :  determination max fini}), for any Borel set $F \subset \mathcal{X}$ we have
\[
\left| \int_F s_{\theta_0}(x) p_{\theta_0}(x) dx \right|
\le \frac{1}{2} \int_{\mathcal{X}} |s_{\theta_0}(x)| p_{\theta_0}(x) dx
= \left| \int_{F^+_{\text{max}}} s_{\theta_0}(x) p_{\theta_0}(x) dx \right|
= \left| \int_{F^{\prime +}_{\text{max}}} s_{\theta_0}(x) p_{\theta_0}(x) dx \right|.
\]
Inserting this bound into \eqref{eq : majo I in prop alpha small ctn}, we find
\begin{equation*}
	\cI_{\theta_0}(q \circ \cP) \le
	(e^\alpha - 1)^2
	\frac{ \left( \int_{\mathcal{X}} |s_{\theta_0}(x)| p_{\theta_0}(x) dx \right)^2 }{4}
	\int_\cE \frac{1}{\widetilde{p}_{\theta_0}(r)} \mu(dr).
\end{equation*}
Recalling from \eqref{eq : tilde p ctn} that $\widetilde{p}_{\theta_0}(r) \ge 1$ for all $r \in \cE$ (since $r(x) \ge 1$), and using that $\mu$ is a sub-probability measure, we get
\[
\int_\cE \frac{1}{\widetilde{p}_{\theta_0}(r)} \mu(dr) \le \mu(\cE) \le 1.
\]
This establishes the upper bound in \eqref{eq :  encadrement Fisher ctn small alpha}.

\medskip

We now turn to the lower bound. It is sufficient to construct an $\alpha$-LDP mechanism whose Fisher information exceeds the lower bound in \eqref{eq : encadrement Fisher ctn small alpha}. Define the two extremal elements
\[
r_1 := 1 + (e^\alpha - 1) \indi{F^+_{\text{max}}}, \qquad
r_2 := 1 + (e^\alpha - 1) \indi{F^{\prime +}_{\text{max}}},
\]
which clearly belong to $\mathcal{E}$. Mimicking the construction in the discrete case (see Equation~\eqref{eq: optimal q 49.5}), we define the randomized mechanism $q^{(\mu)}$ with
\begin{equation} \label{eq: qmu 28.5}
	\mu := \frac{1}{1 + e^\alpha} \left( \delta_{r_1} + \delta_{r_2} \right).
\end{equation}
We verify that this mechanism satisfies the normalization condition \eqref{eq : cond norm mu continue}. Since $e_x(r_j) = r_j(x)$ a.e.\ for $j=1,2$ (as established in Subsection~\ref{ss:point eval}), it suffices to check that
\[
\int_\cE e_x(r) \mu(dr)=\frac{1}{1 + e^\alpha} \sum_{j=1}^2 r_j(x) = 1 \quad \text{a.e.},
\]
which follows directly from the disjointness of $F^+_{\text{max}}$ and $F^{\prime +}_{\text{max}}$ and their union being the whole space $\mathcal{X}$ up to null sets. Note that for $j=1,2$,
\[
\int_{\mathcal{X}} s_{\theta_0}(x) r_j(x) p_{\theta_0}(x) dx
= \int_{\mathcal{X}} s_{\theta_0}(x) p_{\theta_0}(x) dx +
(e^\alpha - 1) \int_{F^{+}_{\text{max}} \text{ or } F^{\prime +}_{\text{max}}} s_{\theta_0}(x) p_{\theta_0}(x) dx,
\]
and the first term is zero by centering of $s_{\theta_0}$. The Fisher information of this mechanism therefore is, by \eqref{eq : Fisher ctn general_preli},
\begin{align*}
	\cI_{\theta_0}(q^{(\mu)} \circ \cP)
	&= \frac{(e^\alpha - 1)^2}{1 + e^\alpha}
	\left[
	\frac{\left( \int_{F^+_{\text{max}}} s_{\theta_0}(x) p_{\theta_0}(x) dx \right)^2}
	{\widetilde{p}_{\theta_0}(r_1)} +
	\frac{\left( \int_{F^{\prime +}_{\text{max}}} s_{\theta_0}(x) p_{\theta_0}(x) dx \right)^2}
	{\widetilde{p}_{\theta_0}(r_2)}
	\right].
\end{align*}
 Moreover, $r_j(x) \le e^\alpha$ implies $\widetilde{p}_{\theta_0}(r_j) \le e^\alpha$ for $j = 1,2$. Since the two integrals in the numerator are equal in absolute value and each equals
\[
\frac{1}{2} \int_{\mathcal{X}} |s_{\theta_0}(x)| p_{\theta_0}(x) dx,
\]
we deduce the lower bound
\begin{equation*}
	\cI_{\theta_0}(q^{(\mu)} \circ \cP) \ge
	\frac{(e^\alpha - 1)^2}{1 + e^\alpha} \cdot \frac{2}{e^\alpha}
	\left( \frac{1}{2} \int_{\mathcal{X}} |s_{\theta_0}(x)| p_{\theta_0}(x) dx \right)^2.
\end{equation*}
This concludes the proof of the lower bound in the theorem.

\end{proof}



\subsubsection{Proof of Proposition \ref{prop : maximisation pb convex}}
\begin{proof}
Since the mapping $r \mapsto j(r)$ is bounded by a constant, we deduce that $J^\star < \infty$. Moreover, the set of Radon sub-probability measures $\mu$ satisfying the constraint in \eqref{eq : optimisation pb ctn J} is non-empty. One example is $\mu$ being the Dirac mass on the function $r_0$ defined by $r_0(x)=1$ for all $x\in \cX$. 
% of such a measure is given in the proof of Theorem \ref{thm : equivalent alpha petit ctn}, where it is shown that $q^{(\mu)} = e_x(r) \mu(dr)$ satisfies \eqref{eq : cond norm mu continue}, with $q^{(\mu)}$ defined as in \eqref{eq: qmu 28.5}. 
Hence, there exists a sequence of sub-probabilities $\mu_n$ on $\mathcal{E}$ satisfying the constraint in \eqref{eq : optimisation pb ctn J} and such that
\[
\int_\mathcal{E} j(r) \mu_n(dr) \xrightarrow{n \to \infty} J^\star.
\]
Since each $\mu_n$ is supported on $\mathcal{E}$, we have $\mu_n(\theconvex) = \mu_n(\mathcal{E}) \leq 1$. As $\theconvex$ is compact for the weak-$\star$ topology, the sequence $(\mu_n)_n$ is tight and thus admits a converging subsequence. Let $k : \mathbb{N} \to \mathbb{N}$ be a strictly increasing function such that $(\mu_{k(n)})_n$ converges weakly to some Radon sub-probability measure $\mu^\star$ on $\theconvex$.


The main difficulty in the proof is that although $\mu_n$ are supported on $\mathcal{E}$, the limit $\mu^\star$ is not necessarily supported on $\mathcal{E}$, since $\mathcal{E}$ is not closed. We will show that the existence of an optimal measure $\overline{\mu}$ supported on $\mathcal{E}$ follows from a convexity argument.


First, we show that $\mu^\star$ inherits properties from the linear constraints in \eqref{eq : optimisation pb ctn J}. Let $h \in L^1(\mathcal{X})$. Then, the map 
\[
r \in \theconvex \mapsto \int_\mathcal{X} h(x) r(x) dx
\]
is continuous for the weak-$\star$ topology and bounded. Hence, the weak convergence of $\mu_{k(n)}$ towards $ \mu^\star$ implies
\begin{equation} \label{eq : proof closeness contraint}
\int_\theconvex \left( \int_\mathcal{X} h(x) r(x) dx \right) \mu_{k(n)}(dr) \xrightarrow{n \to \infty} \int_\theconvex \left( \int_\mathcal{X} h(x) r(x) dx \right) \mu^\star(dr).
\end{equation}
Using that $e_x(r) = r(x)$ a.e., the left-hand side becomes by Fubini theorem,
\[
\int_\mathcal{X} h(x) \left( \int_\theconvex e_x(r) \mu_{k(n)}(dr) \right) dx.
\]
By the constraint $\int_\mathcal{E} e_x(r) \mu_{k(n)}(dr) = 1$ a.e.\ in $x$, this simplifies to $\int_\mathcal{X} h(x) dx$. Similarly, the right-hand side of \eqref{eq : proof closeness contraint} becomes
\[
\int_\mathcal{X} h(x) \left( \int_\theconvex e_x(r) \mu^\star(dr) \right) dx.
\]
Therefore,
\[
\int_\mathcal{X} h(x) dx = \int_\mathcal{X} h(x) \left( \int_\theconvex e_x(r) \mu^\star(dr) \right) dx, \quad \text{for all } h \in L^1(\mathcal{X}),
\]
implying that
\[
\int_\theconvex e_x(r) \mu^\star(dr) = 1 \quad \text{a.e.\ in } x.
\]
Since $r \mapsto j(r)$ is continuous and bounded, we have
\[
J^\star = \lim_{n \to \infty} \int_\theconvex j(r) \mu_{k(n)}(dr) = \int_\theconvex j(r) \mu^\star(dr).
\]
Consequently, $\mu^\star$ is a solution to the optimization problem \eqref{eq : optimisation pb ctn J} when integrals are taken over the convex set $\theconvex$ instead of its extremal points $\mathcal{E}$.

\medskip

We now construct a measure $\overline{\mu}$ supported on $\mathcal{E}$. The proof uses results from the proof of Choquet’s theorem in \cite{Choquet}. Denote
\[
\mathfrak{M}^+ := \{ \mu \mid \text{positive Radon measures on } \theconvex \}.
\]
Following \cite{Choquet}, we endow $\mathfrak{M}^+$ with the partial order: for $\mu_1, \mu_2 \in \mathfrak{M}^+$, we write $\mu_1 \leq \mu_2$ if
\begin{equation} \label{eq : def rel ordre Choquet}
\left\{
\begin{aligned}
&\mu_1(\theconvex) = \mu_2(\theconvex), \\
&\int_\theconvex r \mu_1(dr) = \int_\theconvex r \mu_2(dr), \\
&\int_\theconvex f(r) \mu_1(dr) \leq \int_\theconvex f(r) \mu_2(dr), \quad \text{for all continuous convex } f : \theconvex \to \mathbb{R}_+.
\end{aligned}
\right.
\end{equation}
Lemma 1 in \cite{Choquet} shows that $(\mathfrak{M}^+, \leq)$ is inductive, meaning any chain admits an upper bound. Corollary 1 in \cite{Choquet}, which is a consequence of Zorn's lemma, then implies that any $\mu \in \mathfrak{M}^+$ is upper bounded by a maximal element $\nu \in \mathfrak{M}^+$. Applying this to $\mu^\star$, there exists a maximal element $\overline{\mu} \in \mathfrak{M}^+$ such that $\mu^\star \leq \overline{\mu}$. By the first point in the definition 
 $\mu^\star \leq \overline{\mu}$, we have $1 \geq \mu^\star(\theconvex) = \overline{\mu}(\theconvex)$, so $\overline{\mu}$ is a sub-probability. Moreover,
\[
\int_\theconvex r \mu^\star(dr) = \int_\theconvex r \overline{\mu}(dr),
\]
and thus, applying $e_x$ by Proposition \ref{p: evaluation},
\[
\int_\theconvex e_x(r) \overline{\mu}(dr) = \int_\theconvex e_x(r) \mu^\star(dr) = 1 \quad \text{a.e.\ in } x.
\]
Applying the third property in \eqref{eq : def rel ordre Choquet} with the convex function $f = j$, we obtain
\[
J^\star = \int_\theconvex j(r) \mu^\star(dr) \leq \int_\theconvex j(r) \overline{\mu}(dr).
\]
By Lemma 5 in \cite{Choquet}, any maximal element in $\mathfrak{M}^+$ is supported on the set of extremal points of $\theconvex$, which is $\mathcal{E}$. Consequently,
\[
\int_\mathcal{E} e_x(r) \overline{\mu}(dr) = \int_\theconvex e_x(r) \overline{\mu}(dr) = 1 \quad \text{a.e.\ in } x,
\]
and
\[
J^\star \leq \int_\mathcal{\theconvex} j(r) \overline{\mu}(dr)=\int_\mathcal{E} j(r) \overline{\mu}(dr).
\]
By definition of $J^\star$ as the supremum, we conclude
\[
J^\star = \int_\mathcal{E} j(r) \overline{\mu}(dr).
\]
Hence, $\overline{\mu}$ is a solution to the optimization problem \eqref{eq : optimisation pb ctn J}.
\end{proof}







\subsubsection{Proof of Theorem \ref{thm: main opt cont}}
The proof of the theorem is based on the fact that we can apply Proposition \ref{prop : maximisation pb convex} to the function \( i \) defined as in \eqref{eq : def i r ctn convex}. This is justified by the following lemma.
\begin{lem}\label{lem : i ctn cvx}
The function 
\[
i: \begin{cases} 
& \theconvex \to \R_+ \\ 
&r \mapsto i(r)
\end{cases}
\]
is bounded, continuous, and convex.
\end{lem}
\begin{proof}
Since \( s_{\theta_0} \in L^2(p_{\theta_0}) \) and \( 1 \le r \le e^{\alpha} \) for \( r \in \theconvex \), we deduce that
\[
i(r) \le e^{2\alpha} \left( \int_\cX \abs{s_{\theta_0}(x)} p_{\theta_0}(x) \, dx \right)^2
\le e^{2\alpha} \norm{s_{\theta_0}}_{ L^2(p_{\theta_0})}^2, ~
\quad \text{for all } r \in \theconvex.
\]
Hence, \( i \) is bounded. Also, the functions \( x \mapsto s_{\theta_0}(x) p_{\theta_0}(x) \) and \( x \mapsto p_{\theta_0}(x) \) belong to \( L^1(\cX) \), thus the mapping \( r \mapsto i(r) \) is continuous in the weak-\(\star\) topology
recalling the expression \eqref{eq : def i r ctn convex}. The convexity of \( i \) follows from the algebraic properties:
\begin{align*}
i(\lambda r) &= \abs{\lambda} \, i(r), \quad \text{for all } \lambda \in \R,\, r \in \theconvex, \\
i(r_1 + r_2) &\le i(r_1) + i(r_2), \quad \text{for all } r_1, r_2 \in \theconvex.
\end{align*}
These properties are verified through straightforward computations using the definition \eqref{eq : def i r ctn convex}.
\end{proof}

\begin{comment}
\begin{prop} \label{prop : existence mu sup pb info ctn}
There exists a Radon sub-probability \( \overline{\mu} \) on \( \cE \) such that
\[
I^* := \int_\cE i(r) \, \overline{\mu}(dr), \quad
\text{and } \quad  \int_\cE e_x(r) \, \overline{\mu}(dr) = 1 \quad \text{for a.e. } dx.
\]
\end{prop}

\begin{proof}
It is a direct application of Lemma \ref{lem : i ctn cvx} and Proposition \ref{prop : maximisation pb convex}.
\end{proof}
\end{comment}

We can now invoke Lemma \ref{lem : i ctn cvx} to prove Theorem \ref{thm: main opt cont}, as detailed below.


\begin{proof}
For any \( q \in \cQ_\alpha \), by the factorization result in Proposition \ref{prop : factor ctn}, we can write \( q = q^{(2)} \circ q^{(\mu)} \). Using the fact that Fisher information is reduced by randomization, we have
\[
\cI_{\theta_0}(q \circ \cP) \le \cI_{\theta_0}(q^{(\mu)} \circ \cP).
\]
Thus,
\[
\cJ^{\max,\alpha}_{\theta_0} \le \sup_{\mu} \cI_{\theta_0}(q^{(\mu)} \circ \cP),
\]
where the supremum is over all Radon sub-probabilities satisfying the condition \eqref{eq : cond norm mu continue}. 
From \eqref{eq : Fisher ctn general_preli}, we know that
\[
\cI_{\theta_0}(q^{(\mu)} \circ \cP) = \int_\cE i(r) \, \mu(dr),
\]
where \( i \) is defined in \eqref{eq : def i r ctn convex}. By Proposition \ref{prop : maximisation pb convex} (applicable by Lemma \ref{lem : i ctn cvx}), there exists a Radon sub-probability \( \overline{\mu} \) on \( \cE \) such that
\[
\sup_{\mu} \int_\cE i(r) \, \mu(dr) = \int_\cE i(r) \, \overline{\mu}(dr) = \cI_{\theta_0}(q^{(\overline{\mu})} \circ \cP).
\]
It follows that
\[
\cJ^{\max,\alpha}_{\theta_0} \le \cI_{\theta_0}(q^{(\overline{\mu})} \circ \cP).
\]

We now show the reverse inequality \( \cI_{\theta_0}(q^{(\overline{\mu})} \circ \cP) \le \cJ^{\max,\alpha}_{\theta_0} \). As the mechanism \( q^{(\overline{\mu})} \) is not valued in a countable space, this is not straightforward. From Lemma \ref{lem : expression tilde p et score ctn}, the experiment \( q^{(\overline{\mu})} \circ \cP \) is DQM with continuous density \( \widetilde{p}_{\theta_0}(r) \) and score function \( t_{\theta_0}(r) \). Let \( \varepsilon > 0 \). Since \( \theconvex \) is compact for the \( d_\star \) distance, we can write
\[
\theconvex = \bigsqcup_{i=1}^N C_i,
\]
where \( C_1, \dots, C_N \) are measurable sets forming a partition with \( d_\star \)-radius at most \( \varepsilon \). Define the map \( T^\varepsilon: \theconvex \to \{1, \dots, N\} \) by
\[
T^\varepsilon(r) = \sum_{i=1}^N \indi{C_i}(r),
\]
which is a randomization. Then set \( q^{(\overline{\mu}), \varepsilon} := T^\varepsilon \circ q^{(\overline{\mu})} \), which defines an \( \alpha \)-LDP randomization from \( \cX \) to $\{1, \dots, N\}$.
By Lemma 3.1 of \cite{steinbergerEfficiencyLocalDifferent}, the experiment 
\[
q^{(\overline{\mu}),\varepsilon} \circ \cP 
= T^\varepsilon \circ (q^{(\overline{\mu})} \circ \cP)
\]
is DQM with density
\[
p^\varepsilon_{\theta_0}(i) = \int_{C_i} \widetilde{p}_{\theta_0}(r) \, \overline{\mu}(dr),
\quad i \in \{1, \dots, N\},
\]
and score function
\[
t^\varepsilon_{\theta_0}(i) = \frac{ \int_{C_i} t_{\theta_0}(r) \widetilde{p}_{\theta_0}(r) \, \overline{\mu}(dr) }{ p^\varepsilon_{\theta_0}(i) }.
\]
Again from Lemma 3.1 in \cite{steinbergerEfficiencyLocalDifferent}, the Fisher information satisfies:
\begin{equation} \label{eq : loss Info quantif}
\cI_{\theta_0}(T^\varepsilon \circ (q^{(\overline{\mu})} \circ \cP)) =
\cI_{\theta_0}(q^{(\overline{\mu})} \circ \cP) -
\int_\cE \left( t_{\theta_0}(r) - t^\varepsilon_{\theta_0}(T^\varepsilon(r)) \right)^2
\widetilde{p}_{\theta_0}(r) \, \overline{\mu}(dr).
\end{equation}
Since \( t_{\theta_0} \) and \( \widetilde{p}_{\theta_0} \) are uniformly continuous on the compact set \( \theconvex \), we can make
\[
\sup_{r \in \cE} \abs{ t_{\theta_0}(r) - t^\varepsilon_{\theta_0}(T^\varepsilon(r)) }
\]
arbitrarily small by choosing \( \varepsilon \) small enough. Consequently, the integral on the RHS of \eqref{eq : loss Info quantif} can be made arbitrarily small.

Thus, \( \cI_{\theta_0}(q^{(\overline{\mu})} \circ \cP) \) can be approximated by the Fisher information of finite-valued \( \alpha \)-LDP mechanisms, implying:
\[
\cI_{\theta_0}(q^{(\overline{\mu})} \circ \cP) \le \sup_{q \in \cQ_\alpha} \cI_{\theta_0}(q \circ \cP) = \cJ^{\max,\alpha}_{\theta_0}.
\]
This concludes the proof.
\end{proof}




\section{Proofs of technical results}{\label{s: proof technical}}
This section presents the detailed proofs of the properties stated in Section \ref{S: tools}, as well as the proof of some technical lemmas needed to prove our main results.
\subsection{Proofs of technical results in functional analysis}
We begin by proving that the set $\theconvex$ is compact with respect to the weak-$\star$ topology, as stated in Lemma \ref{l: B closed}. We then proceed to Lemma \ref{l: B separable}, which asserts that $\widehat{B}$, equipped with the weak-$\star$ topology, is separable. This separability is crucial to ensuring that the integral $\int_{\widehat{B}} f \, \mu(df)$, introduced in Lemma \ref{l: Bochner}, is well-defined in the sense of Bochner. Establishing this will, in turn, allow us to rigorously justify the pointwise evaluation given in Proposition \ref{p: evaluation}. The proof of this proposition, which is the most technically involved part of this section, will conclude our exposition of the technical tools.


\subsubsection{Proof of Lemma \ref{l: B closed}}
\begin{proof}
Let us begin by proving that the set \( \theconvex \) is closed with respect to the weak-\(\star\) topology. Let \( (f_n)_n \subset C \) be a sequence such that \( f_n \xrightharpoonup{\star} f \). Since \( C \subset \widehat{B}=B_{\|\cdot\|_\infty}(0, e^{\alpha}) \), and the closed ball $\widehat{B}$ %\( B_{\|\cdot\|_\infty}(0, e^{\alpha}) \) 
is compact in the weak-\(\star\) topology (by the Banach–Alaoglu Theorem, as stated in Theorem~\ref{th: BanAla}, together with the fact that dilations are continuous mappings in the strong topology, hence also in the weak and weak-\(\star\) topologies), it follows that $ f \in \widehat{B}$. % B_{\|\cdot\|_\infty}(0, e^{\alpha}) \). 
To conclude that \( f \in \theconvex \), it remains to show that \( f(x) \geq 1 \) almost everywhere on \( \mathcal{X} \). 

Recall that, by the definition of weak-\(\star\) convergence, for all \( \varphi \in L^1(\mathcal{X}) \), we have:
\[
\langle \varphi, f_n \rangle_{L^1, L^\infty} \longrightarrow \langle \varphi, f \rangle_{L^1, L^\infty}, \quad \text{as } n \to \infty.
\]
Since each \( f_n \in \theconvex \), we have \( f_n(x) \geq 1 \) almost everywhere, and thus for $\varphi \in L^1(\cX)$, $\varphi \ge0$ :
\[
\int_{\mathcal{X}} \varphi(x) \cdot 1 \, dx \leq \int_{\mathcal{X}} \varphi(x) f_n(x) \, dx.
\]
Passing to the limit, we obtain:
\[
\int_{\mathcal{X}} \varphi(x) \, dx \leq \int_{\mathcal{X}} \varphi(x) f(x) \, dx,
\]
which implies:
\[
\int_{\mathcal{X}} \varphi(x) [f(x) - 1] \, dx \geq 0, \quad \forall \varphi \in L^1(\mathcal{X}) \text{ with $\varphi \ge 0$.}
\]
This inequality shows that \( f(x) \geq 1 \) almost everywhere, as desired.

We conclude that \( f \in \theconvex \), hence \( \theconvex \) is closed in the weak-\(\star\) topology. Since \( \theconvex \) is also contained in a weak-\(\star\) compact set, it follows that \( \theconvex \) is compact in the weak-\(\star\) topology.
% by Banach–Alaoglu, using again that dilations and translations are continuous in the strong topology, and therefore also in the weak and weak-\(\star\) topologies.
\end{proof}


\subsubsection{Proof of Lemma \ref{l: B separable}}
\begin{proof}
The proof of this lemma consists in explicitly constructing a countable sequence that is dense in \( \widehat{B} \). It is well known that \( L^1(\cX) \) is separable, so we may fix a countable dense sequence \( (f_n)_{n \in \mathbb{N}} \subset L^1(\cX) \). Define the truncated functions \( \tilde{f}_n := -e^\alpha \lor f_n \land e^\alpha \), which belong to \( \widehat{B} \). Our goal is to show that the sequence \( (\tilde{f}_n)_{n \in \mathbb{N}} \) is dense in \( \widehat{B} \) with respect to the weak-\(\star\) topology.

To this end, let \( g \in \widehat{B} \). Consider the truncations \( g_M := g \cdot \indi{\{|x| \le M\}} \). It satisfies $(-e^\alpha) \vee g_M \wedge e^\alpha$
$\xrightharpoonup{\star} g$  as \( M \to \infty \). Indeed, the sequence $(-e^\alpha)\vee g_M(x) \wedge e^\alpha$ converges to  $(-e^\alpha) \vee g(x) \wedge e^\alpha=g(x)$ for almost every $x \in \cX$.
%, which satisfy \( g_M \xrightharpoonup{\star} g \) as \( M \to \infty \). 
%Indeed, since both \( g_M \) and \( g \) lie in \( L^\infty(\mathbb{R}^d) \), 
Hence, the dominated convergence theorem can be applied, ensuring that for any \( f \in L^1(\mathcal{X}) \),
\[
\int_{\cX} f(x) (-e^\alpha)\vee g_M(x) \wedge e^\alpha \, dx \to \int_{\cX} f(x) g(x) \, dx \qquad \mbox{ for } M\rightarrow \infty.
\]
Thus, it suffices to approximate each \( (-e^\alpha )\vee g_M \wedge e^\alpha \) by elements of the sequence \( (\tilde{f}_n)_{n \in \mathbb{N}} \).

Observe that \( g_M \in L^1(\cX) \). Since \( (f_n) \) is dense in \( L^1(\cX) \), there exists a strictly increasing map \( j : \mathbb{N} \to \mathbb{N} \) such that \( f_{j(n)} \to g_M \) in \( L^1(\cX) \). Up to passing to a subsequence, we may also assume \( f_{j(n)}(x) \to g_M(x) \) for almost every \( x \in \cX \). %Given that definitively \( g_M \) satisfies \( 1 \le g_M(x) \le e^\alpha \), 
It follows that
\[
\tilde{f}_{j(n)}(x) = (-e^\alpha) \lor f_{j(n)}(x) \land e^\alpha \to (-e^\alpha)\vee g_M(x) \wedge e^\alpha  \quad \text{a.e. in } \cX,
\]
and \( \tilde{f}_{j(n)} \in B \) for all \( n \). Applying the dominated convergence theorem once more, we conclude that for all \( \varphi \in L^1(\cX) \),
\[
\int_{\cX} \varphi(x) \tilde{f}_{j(n)}(x) \, dx \to \int_{\cX} \varphi(x)  \left( (-e^\alpha)\vee g_M(x) \wedge e^\alpha \right) \, dx,
\]
so that \( \tilde{f}_{j(n)} \xrightharpoonup{\star}  (-e^\alpha)\vee g_M \wedge e^\alpha \). This proves the lemma.  
%Since \( g_M \xrightharpoonup{\star} g \), we conclude that the sequence \( (\tilde{f}_n) \) is weak-\(\star\) dense in \( B \), as desired.
\end{proof}


\subsubsection{Proof of Lemma \ref{l: Bochner}}
\begin{proof}
%We prove the claim first for the set \( \hat{B} := B_{\|\cdot\|_\infty}(0, e^{\alpha}) \), and then deduce the result for \( B \) by observing that \( B \subset \hat{B} \). As a preliminary step, we note that the proof of Lemma \ref{l: B separable} applies verbatim to \( \hat{B} \), which immediately implies that \( \hat{B} \), endowed with the weak-* topology, is separable. Moreover, by Lemma \ref{l: B metrizable}, this topology is metrizable with the compatible distance \( d_\star \) defined in \eqref{eq: def dstar}.

Using Lemma \ref{l: B separable}, we let \( (f_n)_{n \in \mathbb{N}} \) be a countable sequence dense in \( \widehat{B} \) for the $d_\star$ distance. Fix \( M \geq 1 \) and define a sequence of subsets of \( \widehat{B} \) as follows:
\[
\begin{aligned}
B_0^{(M)} & := \left\{ f \in \widehat{B} \,:\, d_\star(f, f_0) < \frac{1}{M} \right\}, \\
B_1^{(M)} & := \left\{ f \in \widehat{B} \,:\, d_\star(f, f_1) < \frac{1}{M} \right\} \setminus B_0^{(M)}, \\
& \ \, \vdots \\
B_n^{(M)} & := \left\{ f \in \widehat{B} \,:\, d_\star(f, f_n) < \frac{1}{M} \right\} \setminus \bigcup_{j=0}^{n-1} B_j^{(M)},
\end{aligned}
\]
and so on, for all \( n \in \mathbb{N} \). By construction, these sets are disjoint and, since \( (f_n)_{n \in \mathbb{N}} \) is dense in \( \widehat{B} \), we have \( \bigcup_{n=0}^\infty B_n^{(M)} = \widehat{B} \). 
Moreover, by compactness of $\widehat{B}$, we know that any covering of $\widehat{B}$ by open balls of radius $1/M$ admits a finite covering. Thus, for $n$ large enough, the set $B^{(M)}_n$ are empty. In turn,
the union $\bigcup_{n=0}^\infty B_n^{(M)}$ is finite partition of \( \widehat{B} \), which can be used to approximate integrals on \( \widehat{B} \) via simple functions. 

For each \( f \in \widehat{B} \), define the function
\begin{equation}{\label{eq: Id f}}
I_d^{(M)}(f) := \sum_{n = 0}^\infty f_n \indi{B_n^{(M)}}(f),
\end{equation}
where the above sum is actually on a finite number of indexes, depending on $M>0$.
By construction, \( I_d^{(M)}(f) \in \widehat{B} \), and
\[
d_\star\left( I_d^{(M)}(f), f \right) \le \frac{1}{M}.
\]
Furthermore, the integral of the simple map \(  f \mapsto I_d^{(M)}(f) \) is 
%a Bochner integral, 
given explicitly by a sum :
\begin{equation}{\label{eq: int Id f}}
\int_{\widehat{B}} I_d^{(M)}(f) \, \mu(df) = \sum_{n = 0}^\infty f_n \mu\left(B_n^{(M)}\right).
\end{equation}
Since each \( f_n \in \widehat{B} \), we have
\[
\left\| \int_{\widehat{B}} I_d^{(M)}(f) \, \mu(df) \right\|_\infty \le \sup_n \|f_n\|_\infty \sum_{n = 0}^\infty \mu(B_n^{(M)}) \le e^\alpha,
\]
which implies that \( \int_{\widehat{B}} I_d^{(M)}(f) \, \mu(df) \in \widehat{B} \).

Thus, we have constructed a sequence of simple functions that approximate the integral in the sense of Bochner:
\[
\lim_{M \to \infty} \int_{\widehat{B}} d_\star\left( I_d^{(M)}(f), f \right) \, \mu(df) = 0,
\]
as required by the definition of the Bochner integral. This completes the proof.
\end{proof}
\begin{rem}{\label{rk: Cauchy}}
For $M > 0$ and $H > 0$, we have $d_\star\left( I_d^{(M)}(f), I_d^{(M+H)}(f) \right) \le \frac{2}{M}.$
From the definition of $d_\star$ and the fact that $\mu(\widehat{B}) \le 1$, it is possible to check that
\[
d_\star \left( \int_{\widehat{B}} I_d^{(M)}(f) \, \mu(df), \int_{\widehat{B}} I_d^{(M+H)}(f) \, \mu(df) \right) \le \frac{2}{M},
\]
which shows that the sequence of approximating integrals is Cauchy. Moreover, we obtain
\begin{equation}\label{eq : approx Bochner d star}
d_\star \left( \int_{\widehat{B}} I_d^{(M)}(f) \, \mu(df), \int_{\widehat{B}} f \, \mu(df) \right) \le \frac{2}{M},
\end{equation}
a bound that will be useful in the next proof.
\end{rem}
\subsubsection{Proof of Proposition \ref{p: evaluation}}
\begin{proof}
	To prove the desired result, we proceed through a structured argument split into four key steps, which we outline below and then establish rigorously.
%	We begin by recalling that, according to the second point of Choquet’s theorem, any element \( b \in B \) can be represented as
%	\[
%	b = \int_B f \, d\mu(f),
%	\]
%	where \( \mu \) is a probability measure supported on a subset \( \mathcal{E} \subset B \), and vanishes on \( B \setminus E \). This means that the integral effectively reduces to one over the set \( \mathcal{E} \), and \( b \) can be seen as a Bochner integral of functions in \( B \) thanks to Lemma \ref{l: Bochner}.	
%	Recall also the definition of the evaluation operator
%	$e_x(f) = f(x)$
%	as introduced in \eqref{eq: evaluation}. Since \( b \in B \subset L^\infty(\mathbb{R}^d) \), and almost every point \( x \in \mathbb{R}^d \) is a Lebesgue point of \( b \), it follows that
%	\[
%	b(x) = e_x(b) \quad \text{for a.e.\ } x \in \mathbb{R}^d.
%	\]
%	In order to rigorously establish the pointwise representation of \( b \), we now define the set
%	\[
%	C := \left\{ (f,x) \in B \times \mathbb{R}^d : \limsup_n u_n(f,x) = \liminf_n u_n(f,x) \right\},
%	\]
%	where we have denoted
%	\[
%	u_n(f,x) := \frac{1}{|B(x,1/n)|} \int_{B(x,1/n)} f(\omega) \, d\omega.
%	\]
%	The set \( C \) therefore captures all pairs \( (f,x) \) for which the sequence \( (u_n(f,x))_n \) converges, i.e., where \( x \) is a Lebesgue point for the function \( f \). Using a standard characterization of convergence, we may rewrite \( C \) as
%	\[
%	C = \bigcap_{q=1}^{\infty} \bigcup_{n=1}^{\infty} \bigcap_{p_1,p_2 \geq n} \left\{ (f,x) : |u_{p_1}(f,x) - u_{p_2}(f,x)| \leq \frac{1}{q} \right\}.
%	\]
%	We now present the strategy in four steps:\\
%	
\\
	{Step 1:} We will prove that the map $(f,x) \mapsto e_x(f)$ is bi-measurable, as in Point 1. 
	%We will prove that the set \( C \subset B \times \mathbb{R}^d \) is bi-measurable, i.e., measurable with respect to the product \( \sigma \)-algebra \( \mathcal{B}(B) \otimes \mathcal{B}(\mathbb{R}^d) \). \\
%	{Step 2:} We will use the measurability of \( C \) to apply Fubini’s theorem and construct a full-measure set \( \mathcal{L}_x \subset \mathbb{R}^d \), for which the map \( f \mapsto e_x(f) \) is well-defined and integrable with respect to \( \mu \). \\
	\\
	{Step 2:} We will prove the identity \eqref{eq : evalutation rigorous} for simple functions, specifically for the approximations \( I_d(M)(f) \) used to define the Bochner integral. 
	\\
	{Step 3:} We will pass to the limit in these approximations, applying dominated convergence to conclude the desired pointwise identity, concluding the proof of Point 2. 
	\\
	{Step 4:} We will obtain the negligibility of the set \eqref{eq : neg by Fubini}, as stated in Point 3, by using Fubini theorem.
	
	\textbf{Step 1:} Let us  denote
		\begin{equation} \label{eq : def approx of evaluation}
		u_n(f,x) := \frac{1}{|B(x,1/n)|} \int_{B(x,1/n)\cap\cX} f(y) \, dy.
		\end{equation}
	We have $e_x(f)=\limsup_{n \to \infty} u_n(f)$ for all $(f,x) \in \widehat{B}\times\cX$.  Hence, it is sufficient to prove that $(r,x) \mapsto u_n(f,x)$ is measurable.	
%	 In order to prove the desired bi-measurability of the map
%	\[
%	u_n : B \times X \to \mathbb{R}, \qquad (f,x) \mapsto u_n(f,x) := \frac{u_n^{(1)}(f,x)}{u_n^{(2)}(x)},
%	\]
	We analyze separately the continuity properties of the numerator \( u_n^{(1)}(r,x) \) and the denominator \( u_n^{(2)}(x)= |B(x,1/n)|\). We begin with
	\[
	u_n^{(1)}(f,x) = \int_{B(x, \frac{1}{n})\cap\cX} f(y) \, dy = \left\langle \mathbbm{1}_{B(x, \frac{1}{n})}, r \right\rangle_{L^1, L^\infty}.
	\]
	To show continuity of \( u_n^{(1)} \), consider a sequence \( (f_p, x_p) \to (r,x) \) as \( p \to \infty \), where \( f_p \to f \) in the \( d_\star \)-topology on \( \widehat{B} \), and \( x_p \to x \) in the Euclidean topology on \( \mathbb{R}^d \). We estimate
	\[
	|u_n^{(1)}(f,x) - u_n^{(1)}(f_p,x_p)| 
	\le \int_{\cX} \left| \mathbbm{1}_{B(x, \frac{1}{n})} - \mathbbm{1}_{B(x_p, \frac{1}{n})} \right| |f_p(y)| \, dy + \left| \int_{B(x, \frac{1}{n})\cap\cX} [f(y) - f_p(y)] \, dy \right| 
	= I_1 + I_2.
	\]
	Since \( f_p \xrightharpoonup{\star} f \) in \( L^\infty \), we have
	\[
	I_2 = \left\langle \mathbbm{1}_{B(x, \frac{1}{n})}, f - f_p \right\rangle \longrightarrow 0
	\quad \text{as } p \to \infty.
	\]
	To control \( I_1 \), observe that for all \( p \), \( f_p \in \widehat{B} \) and thus satisfy \( \|f_p\|_{L^\infty} \le e^{\alpha} \). Hence,
	\[
	I_1 \le e^{\alpha} \int_{\mathbb{R}^d} \left| \mathbbm{1}_{B(x, \frac{1}{n})} - \mathbbm{1}_{B(x_p, \frac{1}{n})} \right| \, dy,
	\]
	which tends to zero as \( x_p \to x \), since the symmetric difference of the balls \( B(x, \frac{1}{n}) \) and \( B(x_p, \frac{1}{n}) \) vanishes in measure.
	
	A similar argument applies to \( u_n^{(2)}(x) := |B(x, \frac{1}{n})| \), and we omit the details. We thus conclude that both \( u_n^{(1)} \) and \( u_n^{(2)} \) are continuous, and therefore \( (f,x)\mapsto u_n(f,x) \) is continuous as a quotient of continuous functions (with strictly positive denominator).
	
	Since \( u_n \) is continuous on \( \widehat{B} \times \cX \), it is measurable with respect to the Borel \( \sigma \)-algebra \( \mathcal{B}(\widehat{B} \times \cX) \). As both \( \widehat{B} \) and \( \cX \) are metric spaces, we have
	\[
	\mathcal{B}(\widehat{B} \times \mathbb{R}^d) = \mathcal{B}(\widehat{B}) \otimes \mathcal{B}(\mathbb{R}^d),
	\]
	so \( u_n \) is measurable with respect to the product \( \sigma \)-algebra. 
	This step proves the point 1) of the proposition.
	% This implies that the set \( C \subset B \times \mathbb{R}^d \), defined in terms of convergence of the sequence \( u_n \), is bi-measurable, as required. \\
	\\
	\\
	\textbf{Step 2:}  Let $\mu$ be a sub-probability on the space $\widehat{B}$ and assume that $b=\int_{\widehat{B}} f \mu(df)$. 	
	In this step, we aim to establish \eqref{eq : evalutation rigorous} but with $I_d^{(M)}(f)$, as defined in \eqref{eq: Id f} instead of $f$. 
	%Note that this time we are working on the space \( B \), rather than on \( \hat{B} \), and that \( I_d^{(M)}(f) \) is constructed using a countable dense sequence \( (f_n)_n \subset B \). Observe moreover that, since \( f \in B \), and \( I_d^{(M)}(f) \) is built from elements in \( B \), it follows that \( I_d^{(M)}(f) \in B \) as well.
	
	Recall from \eqref{eq: int Id f} that
	\[
	\int_{\widehat{B}} I_d^{(M)}(f) \, \mu(df) = \sum_{n=0}^\infty f_n \, \mu\left(B_n^{(M)}\right),
	\]
	where the equality holds in \( L^\infty(\cX) \). This implies that for almost every \( x \in \cX \),
	\[
	\left( \int_{\widehat{B}} I_d^{(M)}(f) \, \mu(df) \right)(x) = \sum_{n=0}^\infty f_n(x) \, \mu\left(B_n^{(M)}\right).
	\]
	Moreover, since for almost every \( x \in \cX \) it holds that
	\[
	I_d^{(M)}(f)(x) = \sum_{n=0}^\infty f_n(x) \, \mathbbm{1}_{B_n^{(M)}}(f),
	\]
	we have
	\[
	\int_{\widehat{B}} I_d^{(M)}(f)(x) \, \mu(df) = \sum_{n=0}^\infty f_n(x) \, \mu\left(B_n^{(M)}\right),
	\]
	so that
	\[
	\left( \int_{\widehat{B}} I_d^{(M)}(f) \, \mu(df) \right)(x) = \int_{\widehat{B}} I_d^{(M)}(f)(x) \, \mu(df).
	\]
	This is precisely the result we intended to prove, but for the case of simple functions. The goal now is to extend this identity to the Bochner integral
	\[
	\int_{\widehat{B}} f \, \mu(df),
	\]
	which will be the focus of the next step. %(and final) step. \\
	\\
	\\
	\textbf{Step 3.} Let us introduce the following notations:
	\[
	u^{(M)} := \int_{\widehat{B}} I_d^{(M)}(f) \, \mu(df) ,  \quad u:= \int_{\widehat{B}} f \, \mu(df).
	\]
	From equation \eqref{eq : approx Bochner d star} following the proof of Lemma \ref{l: Bochner}, we have $u^{(M)} \xrightharpoonup{\star} u$ as $M \rightarrow \infty$.
%	
%	\[
%	u^{(M)} := \int_{\hat{B}} I_d^{(M)}(f) \, \mu(df) \xrightharpoonup{\star} \int_{\hat{B}} f \, \mu(df) =: u
%	\quad \text{as } M \rightarrow \infty.
%	\]
%	

Thanks to this convergence, we have, for any \( h \in L^1(\cX) \cap L^\infty(\cX) \),
	\begin{align}\nonumber
		\int_{\cX} h(x) u(x) \, dx 
		&= \lim_{M \rightarrow \infty} \int_{\cX} h(x) u^{(M)}(x) \, dx \\
	\nonumber
		&= \lim_{M \rightarrow \infty} \int_{\cX} h(x) \left( \int_{\widehat{B}} I_d^{(M)}(f) \, \mu(df) \right)(x) \, dx \\
		\label{eq : h scalaire u proof prop eval}
		&= \lim_{M \rightarrow \infty} \left( \int_{\cX} \int_{\widehat{B}} h(x) I_d^{(M)}(f)(x) \, \mu(df)\right) \, dx,
	\end{align}
	where in the last equality we used Step 2.
	
	From the definition of \( I_d^{(M)}(f) \), it follows that the function
	\[
	(x, f) \mapsto h(x) I_d^{(M)}(f)(x) = \sum_{n=1}^\infty h(x) f_n(x) \, \mathbbm{1}_{B_n^{(M)}}(f)
	\]
	is bi-measurable and bounded. Therefore, we can apply Fubini's theorem to obtain from \eqref{eq : h scalaire u proof prop eval}:
	\[
	\int_{\cX} h(x) u(x) \, dx=	\lim_{M \rightarrow \infty} \int_{\widehat{B}} \left( \int_{\cX } h(x) I_d^{(M)}(f)(x) \, dx \right) \mu(df).
	\]
	Now, for each \( f \in \widehat{B} \), we have \( I_d^{(M)}(f) \xrightharpoonup{\star} f \) as \( M \to \infty \), which implies
	\[
	\int_{\cX} h(x) I_d^{(M)}(f)(x) \, dx \longrightarrow \int_{\cX} h(x) f(x) \, dx.
	\]
	Moreover, since \( I_d^{(M)}(f) \in \widehat{B} \), this quantity is uniformly bounded by \( \|h\|_{L^1(\cX)} e^\alpha \). On the other hand, the right-hand side above equals \( \int_{\cX} h(x) e_x(f) \, dx \), because \( f(x) = e_x(f) \) almost everywhere.
	%, and this expression is well-defined by Step 2.
	 Hence, we may apply the dominated convergence theorem to deduce:
	\begin{align*}
		\lim_{M \rightarrow \infty} \int_{\widehat{B}} \left( \int_{\cX} h(x) I_d^{(M)}(f)(x) \, dx \right) \mu(df) 
		&= \int_{\widehat{B}} \left( \int_{\cX} h(x) f(x) \, dx \right) \mu(df) \\
		&= \int_{\widehat{B}} \left( \int_{\cX} h(x) e_x(f) \, dx \right) \mu(df).
	\end{align*}
	Since the map \( (x, f) \mapsto h(x) e_x(f) \) is bi-measurable and bounded, we can again apply Fubini's theorem to conclude:
	\[
	\int_{\cX} h(x) u(x) \, dx = \int_{\cX} h(x) \left( \int_{\widehat{B}} e_x(f) \, \mu(df) \right) dx.
	\]
	As this equality holds for every \( h \in L^1(\cX) \cap 
	L^\infty(\cX) \), we deduce that for almost every \( x \in  \cX \),
	\[
	u(x) = \int_{\widehat{B}} e_x(f) \, \mu(df),
	\]
	which is precisely the pointwise identity we set out to rigorously establish. The point 2) of the proposition is obtained.
	\\
	\\
	\\
	\textbf{Step 4:} In this final step, we show the third point of the proposition.
	%In order to rigorously establish the pointwise representation of \( b \), 
	We define the set
		\[
		R := \left\{ (f,x) \in \widehat{B} \times \cX : \limsup_n u_n(f,x) = \liminf_n u_n(f,x) \right\},
		\]
		where we recall that $u_n$ is defined in \eqref{eq : def approx of evaluation}.
	%	\[
	%	u_n(f,x) := \frac{1}{|B(x,1/n)|} \int_{B(x,1/n)} f(\omega) \, d\omega.
	%	\]
		The set \( R \) therefore captures all pairs \( (f,x) \) for which the sequence \( (u_n(f,x))_n \) converges, which is for instance the case when \( x \) is a Lebesgue point for the function \( f \). Using a standard characterization of convergence, we may rewrite \( R \) as
		\[
		R = \bigcap_{q=1}^{\infty} \bigcup_{n=1}^{\infty} \bigcap_{p_1,p_2 \geq n} \left\{ (f,x) : |u_{p_1}(f,x) - u_{p_2}(f,x)| \leq \frac{1}{q} \right\},
		\]
	which establishes the measurability of the set $R$ from those of $u_n$.
	
	We analyze the integral
	\[
	\int_{\cX} \int_{\widehat{B}} \mathbbm{1}_{R^c}(f,x) \, \mu(df) \, dx.
	\]
	Thanks to the measurability,  we may apply Fubini’s theorem to exchange the order of integration:
	\[
	\int_{\cX} \int_{\widehat{B}} \mathbbm{1}_{R^c}(f,x) \, \mu(df) \, dx = \int_{\widehat{B}} \int_{\cX} \mathbbm{1}_{R^c}(f,x) \, dx \, \mu(df).
	\]
	As already discussed, for every \( f \in \widehat{B} \), almost every \( x \in \cX  \) is a Lebesgue point of \( f \). This implies
	\[
	\int_{\cX} \mathbbm{1}_{R^c}(f,x) \, dx \le \int_{\cX} \mathbbm{1}_{\{x \colon x \text{ is not a Lebesgue point of } f\}} \, dx = 0.
	\]
	Thus, we conclude that
	\[
	\int_{\cX} \int_{\widehat{B}}\mathbbm{1}_{R^c}(f,x) \, \mu(df) \, dx = 0.
	\]
	From the definition of the set $R$, this proves the third point of the proposition.
%	
%	Now, define the function
%	\[
%	n(x) := \int_{\widehat{B}} \mathbbm{1}_{R^c}(f,x) \, \mu(df).
%	\]
%	This function is measurable, non-negative, and satisfies \( \int_{\mathbb{R}^d} n(x) \, dx = 0 \), which implies \( n(x) = 0 \) for almost every \( x \in \mathbb{R}^d \).
%	We then introduce the set
%	\[
%	\mathcal{L}_x := \left\{ x \in \mathbb{R}^d : n(x) = 0 \right\}.
%	\]
%	This set \( \mathcal{L}_x \) is precisely the one referred to at the beginning of the proof, in the description of Step 2. Indeed, for any \( x \in \mathcal{L}_x \), the definition of \( n(x) \) gives
%	\[
%	\mu\left( \left\{ f \in B : (u_n(f,x))_n \text{ does not converge} \right\} \right) = 0.
%	\]
%	That is, for almost every \( f \in B \), the sequence \( u_n(f,x) \) converges.
%	
%	Finally, we verify that for all \( x \in \mathcal{L}_x \), the integral
%	\[
%	\int_B e_x(f) \, \mu(df)
%	\]
%	is well defined. This follows from the fact that the map
%	\[
%	e_x : B \to \mathbb{R}, \qquad f \mapsto e_x(f) = f(x)
%	\]
%	is measurable and bounded on \( B \). Hence, the Bochner integral \( \int_B e_x(f) \, \mu(df) \) exists for all \( x \in \mathcal{L}_x \), concluding Step 2. \\
%	\\
%	
	\medskip
	
	\noindent This completes the proof.
	
\end{proof}




%\subsubsection{Proof of Proposition \ref{p: evaluation}}
%\begin{proof}
%To prove the desired result, we proceed through a structured argument split into four key steps, which we outline below and then establish rigorously.
%
%We begin by recalling that, according to the second point of Choquet’s theorem, any element \( b \in B \) can be represented as
%\[
%b = \int_B f \, d\mu(f),
%\]
%where \( \mu \) is a probability measure supported on a subset \( \mathcal{E} \subset B \), and vanishes on \( B \setminus E \). This means that the integral effectively reduces to one over the set \( \mathcal{E} \), and \( b \) can be seen as a Bochner integral of functions in \( B \) thanks to Lemma \ref{l: Bochner}.
%
%Recall also the definition of the evaluation operator
%$e_x(f) = f(x)$
%as introduced in \eqref{eq: evaluation}. Since \( b \in B \subset L^\infty(\mathbb{R}^d) \), and almost every point \( x \in \mathbb{R}^d \) is a Lebesgue point of \( b \), it follows that
%\[
%b(x) = e_x(b) \quad \text{for a.e.\ } x \in \mathbb{R}^d.
%\]
%In order to rigorously establish the pointwise representation of \( b \), we now define the set
%\[
%C := \left\{ (f,x) \in B \times \mathbb{R}^d : \limsup_n u_n(f,x) = \liminf_n u_n(f,x) \right\},
%\]
%where we have denoted
%\[
%u_n(f,x) := \frac{1}{|B(x,1/n)|} \int_{B(x,1/n)} f(\omega) \, d\omega.
%\]
%The set \( C \) therefore captures all pairs \( (f,x) \) for which the sequence \( (u_n(f,x))_n \) converges, i.e., where \( x \) is a Lebesgue point for the function \( f \). Using a standard characterization of convergence, we may rewrite \( C \) as
%\[
%C = \bigcap_{q=1}^{\infty} \bigcup_{n=1}^{\infty} \bigcap_{p_1,p_2 \geq n} \left\{ (f,x) : |u_{p_1}(f,x) - u_{p_2}(f,x)| \leq \frac{1}{q} \right\}.
%\]
%We now present the strategy in four steps:\\
%\\
%{Step 1:} We will prove that the set \( C \subset B \times \mathbb{R}^d \) is bi-measurable, i.e., measurable with respect to the product \( \sigma \)-algebra \( \mathcal{B}(B) \otimes \mathcal{B}(\mathbb{R}^d) \). \\
%\\
%{Step 2:} We will use the measurability of \( C \) to apply Fubini’s theorem and construct a full-measure set \( \mathcal{L}_x \subset \mathbb{R}^d \), for which the map \( f \mapsto e_x(f) \) is well-defined and integrable with respect to \( \mu \). \\
%\\
%{Step 3:} We will prove the desired identity for simple functions, specifically for the approximations \( I_d(M)(f) \) used to define the Bochner integral. \\
%\\
%{Step 4:} We will pass to the limit in these approximations, applying dominated convergence to conclude the desired pointwise identity. \\
%\\
%\textbf{Step 1:} In order to prove the desired bi-measurability of the map
%\[
%u_n : B \times X \to \mathbb{R}, \qquad (f,x) \mapsto u_n(f,x) := \frac{u_n^{(1)}(f,x)}{u_n^{(2)}(x)},
%\]
%we analyze separately the continuity properties of the numerator \( u_n^{(1)}(f,x) \) and the denominator \( u_n^{(2)}(x) \). We begin with
%\[
%u_n^{(1)}(f,x) = \int_{B(x, \frac{1}{n})} f(\omega) \, d\omega = \left\langle \mathbbm{1}_{B(x, \frac{1}{n})}, f \right\rangle_{L^1, L^\infty}.
%\]
%To show continuity of \( u_n^{(1)} \), consider a sequence \( (f_p, x_p) \to (f,x) \) as \( p \to \infty \), where \( f_p \to f \) in the \( d_* \)-topology on \( B \), and \( x_p \to x \) in the Euclidean topology on \( \mathbb{R}^d \). We estimate
%\[
%|u_n^{(1)}(f,x) - u_n^{(1)}(f_p,x_p)| 
%\le \int_{\mathbb{R}^d} \left| \mathbbm{1}_{B(x, \frac{1}{n})} - \mathbbm{1}_{B(x_p, \frac{1}{n})} \right| |f_p(\omega)| \, d\omega + \left| \int_{B(x, \frac{1}{n})} [f(\omega) - f_p(\omega)] \, d\omega \right| 
%= I_1 + I_2.
%\]
%Since \( f_p \xrightharpoonup{\star} f \) in \( L^\infty \), we have
%\[
%I_2 = \left\langle \mathbbm{1}_{B(x, \frac{1}{n})}, f - f_p \right\rangle \longrightarrow 0
%\quad \text{as } p \to \infty.
%\]
%To control \( I_1 \), observe that for all \( p \), \( f_p \in B \) and thus satisfy \( \|f_p\|_{L^\infty} \le e^{\alpha} \). Hence,
%\[
%I_1 \le e^{\alpha} \int_{\mathbb{R}^d} \left| \mathbbm{1}_{B(x, \frac{1}{n})} - \mathbbm{1}_{B(x_p, \frac{1}{n})} \right| \, d\omega,
%\]
%which tends to zero as \( x_p \to x \), since the symmetric difference of the balls \( B(x, \frac{1}{n}) \) and \( B(x_p, \frac{1}{n}) \) vanishes in measure.
%
%A similar argument applies to \( u_n^{(2)}(x) := |B(x, \frac{1}{n})| \), and we omit the details. We thus conclude that both \( u_n^{(1)} \) and \( u_n^{(2)} \) are continuous, and therefore \( u_n(f,x) \) is continuous as a quotient of continuous functions (with strictly positive denominator).
%
%Since \( u_n \) is continuous on \( B \times \mathbb{R}^d \), it is measurable with respect to the Borel \( \sigma \)-algebra \( \mathcal{B}(B \times \mathbb{R}^d) \). As both \( B \) and \( \mathbb{R}^d \) are metric spaces, we have
%\[
%\mathcal{B}(B \times \mathbb{R}^d) = \mathcal{B}(B) \otimes \mathcal{B}(\mathbb{R}^d),
%\]
%so \( u_n \) is measurable with respect to the product \( \sigma \)-algebra. This implies that the set \( C \subset B \times \mathbb{R}^d \), defined in terms of convergence of the sequence \( u_n \), is bi-measurable, as required. \\
%\\
%\textbf{Step 2:} We now proceed to the second step, where we analyze the integral
%\[
%\int_{\mathbb{R}^d} \int_B \mathbbm{1}_{C^c}(f,x) \, \mu(df) \, dx.
%\]
%Thanks to the measurability established in Step 1, we may apply Fubini’s theorem to exchange the order of integration:
%\[
%\int_{\mathbb{R}^d} \int_B \mathbbm{1}_{C^c}(f,x) \, \mu(df) \, dx = \int_B \int_{\mathbb{R}^d} \mathbbm{1}_{C^c}(f,x) \, dx \, \mu(df).
%\]
%As already discussed, for every \( f \in B \), almost every \( x \in \mathbb{R}^d \) is a Lebesgue point of \( f \). This implies
%\[
%\int_{\mathbb{R}^d} \mathbbm{1}_{C^c}(f,x) \, dx \le \int_{\mathbb{R}^d} \mathbbm{1}_{\{x \colon x \text{ is not a Lebesgue point of } f\}} \, dx = 0.
%\]
%Thus, we conclude that
%\[
%\int_{\mathbb{R}^d} \int_B \mathbbm{1}_{C^c}(f,x) \, \mu(df) \, dx = 0.
%\]
%Now, define the function
%\[
%n(x) := \int_B \mathbbm{1}_{C^c}(f,x) \, \mu(df).
%\]
%This function is measurable, non-negative, and satisfies \( \int_{\mathbb{R}^d} n(x) \, dx = 0 \), which implies \( n(x) = 0 \) for almost every \( x \in \mathbb{R}^d \).
%We then introduce the set
%\[
%\mathcal{L}_x := \left\{ x \in \mathbb{R}^d : n(x) = 0 \right\}.
%\]
%This set \( \mathcal{L}_x \) is precisely the one referred to at the beginning of the proof, in the description of Step 2. Indeed, for any \( x \in \mathcal{L}_x \), the definition of \( n(x) \) gives
%\[
%\mu\left( \left\{ f \in B : (u_n(f,x))_n \text{ does not converge} \right\} \right) = 0.
%\]
%That is, for almost every \( f \in B \), the sequence \( u_n(f,x) \) converges.
%
%Finally, we verify that for all \( x \in \mathcal{L}_x \), the integral
%\[
%\int_B e_x(f) \, \mu(df)
%\]
%is well defined. This follows from the fact that the map
%\[
%e_x : B \to \mathbb{R}, \qquad f \mapsto e_x(f) = f(x)
%\]
%is measurable and bounded on \( B \). Hence, the Bochner integral \( \int_B e_x(f) \, \mu(df) \) exists for all \( x \in \mathcal{L}_x \), concluding Step 2. \\
%\\
%\textbf{Step 3:} In this step, we aim to establish the desired result for \( I_d^{(M)}(f) \), as defined in \eqref{eq: Id f}. Note that this time we are working on the space \( B \), rather than on \( \hat{B} \), and that \( I_d^{(M)}(f) \) is constructed using a countable dense sequence \( (f_n)_n \subset B \). Observe moreover that, since \( f \in B \), and \( I_d^{(M)}(f) \) is built from elements in \( B \), it follows that \( I_d^{(M)}(f) \in B \) as well.
%
%Recall from \eqref{eq: int Id f} that
%\[
%\int_{\hat{B}} I_d^{(M)}(f) \, \mu(df) = \sum_{n=0}^\infty f_n \, \mu\left(B_n^{(M)}\right),
%\]
%where the equality holds in \( L^\infty \). This implies that for almost every \( x \in \mathbb{R}^d \),
%\[
%\left( \int_{\hat{B}} I_d^{(M)}(f) \, \mu(df) \right)(x) = \sum_{n=0}^\infty f_n(x) \, \mu\left(B_n^{(M)}\right).
%\]
%Moreover, since for almost every \( x \in \mathbb{R}^d \) it holds that
%\[
%I_d^{(M)}(f)(x) = \sum_{n=0}^\infty f_n(x) \, \mathbbm{1}_{B_n^{(M)}}(f),
%\]
%we have
%\[
%\int_{\hat{B}} I_d^{(M)}(f)(x) \, \mu(df) = \sum_{n=0}^\infty f_n(x) \, \mu\left(B_n^{(M)}\right),
%\]
%so that
%\[
%\left( \int_{\hat{B}} I_d^{(M)}(f) \, \mu(df) \right)(x) = \int_{\hat{B}} I_d^{(M)}(f)(x) \, \mu(df).
%\]
%This is precisely the result we intended to prove, but for the case of simple functions. The goal now is to extend this identity to the Bochner integral
%\[
%\int_{\hat{B}} f \, \mu(df),
%\]
%which will be the focus of the next (and final) step. \\
%\\
%\textbf{Step 4.} Let us introduce the following notation:
%\[
%u^{(M)} := \int_{\hat{B}} I_d^{(M)}(f) \, \mu(df) \xrightharpoonup{\star} \int_{\hat{B}} f \, \mu(df) =: u
%\quad \text{as } M \rightarrow \infty.
%\]
%Thanks to this convergence, we have, for any \( h \in L^1(\mathbb{R}^d) \cap L^\infty(\mathbb{R}^d) \),
%\begin{align*}
%\int_{\mathbb{R}^d} h(x) u(x) \, dx 
%&= \lim_{M \rightarrow \infty} \int_{\mathbb{R}^d} h(x) u^{(M)}(x) \, dx \\
%&= \lim_{M \rightarrow \infty} \int_{\mathbb{R}^d} h(x) \left( \int_{\hat{B}} I_d^{(M)}(f) \, \mu(df) \right)(x) \, dx \\
%&= \lim_{M \rightarrow \infty} \int_{\mathbb{R}^d} \int_{\hat{B}} h(x) I_d^{(M)}(f)(x) \, \mu(df) \, dx,
%\end{align*}
%where in the last equality we used Step 3.
%
%From the definition of \( I_d^{(M)}(f) \), it follows that the function
%\[
%(x, f) \mapsto h(x) I_d^{(M)}(f)(x) = \sum_{n=1}^\infty h(x) f_n(x) \, \mathbbm{1}_{B_n^{(M)}}(f)
%\]
%is bi-measurable and bounded. Therefore, we can apply Fubini's theorem to obtain:
%\[
%\lim_{M \rightarrow \infty} \int_{\hat{B}} \left( \int_{\mathbb{R}^d} h(x) I_d^{(M)}(f)(x) \, dx \right) \mu(df).
%\]
%Now, for each \( f \in B \), we have \( I_d^{(M)}(f) \xrightharpoonup{\star} f \) as \( M \to \infty \), which implies
%\[
%\int_{\mathbb{R}^d} h(x) I_d^{(M)}(f)(x) \, dx \longrightarrow \int_{\mathbb{R}^d} h(x) f(x) \, dx.
%\]
%Moreover, since \( I_d^{(M)}(f) \in B \), this quantity is uniformly bounded by \( \|h\|_{L^1(\mathbb{R}^d)} e^\alpha \). On the other hand, the right-hand side above equals \( \int_{\mathbb{R}^d} h(x) e_x(f) \, dx \), because \( f(x) = e_x(f) \) almost everywhere, and this expression is well-defined by Step 2. Hence, we may apply the dominated convergence theorem to deduce:
%\begin{align*}
%\lim_{M \rightarrow \infty} \int_{\hat{B}} \left( \int_{\mathbb{R}^d} h(x) I_d^{(M)}(f)(x) \, dx \right) \mu(df) 
%&= \int_{\hat{B}} \left( \int_{\mathbb{R}^d} h(x) f(x) \, dx \right) \mu(df) \\
%&= \int_{\hat{B}} \left( \int_{\mathbb{R}^d} h(x) e_x(f) \, dx \right) \mu(df).
%\end{align*}
%Since the map \( (x, f) \mapsto h(x) e_x(f) \) is bi-measurable and bounded, we can again apply Fubini's theorem to conclude:
%\[
%\int_{\mathbb{R}^d} h(x) u(x) \, dx = \int_{\mathbb{R}^d} h(x) \left( \int_{\hat{B}} e_x(f) \, \mu(df) \right) dx.
%\]
%As this equality holds for every \( h \in L^1(\mathbb{R}^d) \cap L^\infty(\mathbb{R}^d) \), we deduce that for almost every \( x \in \mathbb{R}^d \),
%\[
%u(x) = \int_{\hat{B}} e_x(f) \, \mu(df),
%\]
%which is precisely the pointwise identity we set out to rigorously establish.
%
%\medskip
%
%\noindent This completes the proof.
%
%\end{proof}

\subsection{Proof of other lemmas}
We conclude by proving Lemma~\ref{lem : determination max fini}, as stated in Section~\ref{s: proof main} within the proofs of the main results.

\subsubsection{Proof Lemma \ref{lem :  determination max fini}}
\begin{proof}
By definition \eqref{eq : def nplus}, we have $n_\beta^+ \in [0,1]$ for all $\beta$, so we can upper and lower bound the denominator of $i_\beta$, obtaining:
\begin{equation} \label{eq : deno not so important fini}
	e^{-\alpha} \left( \int_\cX s_{\theta}(x) \indi{F^+_\beta}(x) p_\theta(x) \mu(dx) \right)^2
	\le i_\beta \le 
	\left( \int_\cX s_{\theta}(x) \indi{F^+_\beta}(x) p_\theta(x) \mu(dx) \right)^2.
\end{equation}
Regarding the numerator of $i_\beta$, let us recall
\[
F_\text{max}^+ = \{x \in \mathcal{X} \mid s_\theta(x) > 0\}, \quad
F_\text{max}^{\prime +} = \{x \in \mathcal{X} \mid s_\theta(x) < 0\}.
\]
Since $s_\theta(x) \neq 0$ for all $x \in \cX$ by assumption, we have 
\[
F_\text{max}^+ \cup F_\text{max}^{\prime +} = \cX.
\]
Recalling that the score function is centered, i.e., 
\[
\int_\cX s_\theta(x) p_\theta(x) \mu(dx) = 0,
\]
we deduce that:
\begin{equation} \label{eq: symmetry 12.5}
	\int_\cX s_{\theta}(x) \indi{F^+_\text{max}}(x) p_\theta(x) \mu(dx)
	= - \int_\cX s_{\theta}(x) \indi{F^{\prime +}_\text{max}}(x) p_\theta(x) \mu(dx)
	= \frac{1}{2} \int_\cX \abs{s_{\theta}(x)} p_\theta(x) \mu(dx).
\end{equation}
Since $\{F_\text{max}^+, F_\text{max}^{\prime +}\}$ is a partition of $\cX$, for any $F \subset \cX$, we can write:
\[
\int_F s_\theta(x) p_\theta(x) \mu(dx) 
= \int_{F \cap F^+_\text{max}} s_\theta(x) p_\theta(x) \mu(dx)
+ \int_{F \cap F^{\prime +}_\text{max}} s_\theta(x) p_\theta(x) \mu(dx).
\]
This implies:
\[
\int_F s_\theta(x) p_\theta(x) \mu(dx) \in 
\left[ -\frac{1}{2} \int_\cX \abs{s_{\theta}(x)} p_\theta(x) \mu(dx), 
        \frac{1}{2} \int_\cX \abs{s_{\theta}(x)} p_\theta(x) \mu(dx) \right],
\]
and consequently:
\begin{equation} \label{eq : iF smaller iFmax}
	\left( \int_F s_\theta(x) p_\theta(x) \mu(dx) \right)^2
	\le \frac{1}{4} \left( \int_\cX \abs{s_{\theta}(x)} p_\theta(x) \mu(dx) \right)^2.
\end{equation}
Moreover, by assumption $s_\theta(x) \neq 0$ for all $x \in \cX$ and from \eqref{eq: symmetry 12.5}, equality in \eqref{eq : iF smaller iFmax} holds if and only if $F = F_\text{max}^+$ or $F = F_\text{max}^{\prime +}$.

Since the collection of sets $\{F^+_\beta\}_{\beta=0,\dots,2^d-1}$ introduced in \eqref{eq: Fbeta+ 4.8} describes all subsets of $\cX$, there exist indices $\beta_\text{max}, \beta'_\text{max}$ such that
\[
F_{\beta_\text{max}}^+ = F_\text{max}^+, \quad 
F_{\beta'_\text{max}}^+ = F_\text{max}^{\prime +}.
\]
Therefore, the function
\[
\beta \mapsto \left( \int_\cX s_\theta(x) \indi{F^+_\beta}(x) p_\theta(x) \mu(dx) \right)^2
\]
reaches its maximum value only at $\beta = \beta_\text{max}$ and $\beta = \beta'_\text{max}$.
As a consequence, using the bounds in \eqref{eq : deno not so important fini}, for $\alpha$ sufficiently small (say, $\alpha < \overline{\alpha}$), the function $\beta \mapsto i_\beta$ reaches its maximum on $\{\beta_\text{max}, \beta'_\text{max}\}$.
An explicit constraint on $\overline{\alpha}$ is given by:
\[
e^{\overline{\alpha}} 
< \sup_{\beta \notin \{\beta_\text{max}, \beta'_\text{max}\}} 
\frac{
	\left( \int_{F^+_{\beta_\text{max}}} s_\theta(x) p_\theta(x) \mu(dx) \right)^2}{
	\left( \int_{F^+_\beta} s_\theta(x) p_\theta(x) \mu(dx) \right)^2}.
\]
Recalling the definition \eqref{eq : def i beta}, and the fact that $n_\beta^+$ appears in the denominator of $i_\beta$, we observe that:
\begin{itemize}
    \item If $n^+_{\beta_{\text{max}}} \leq n^+_{\beta'_{\text{max}}}$, then $i_{\beta_{\text{max}}} \geq i_{\beta'_{\text{max}}}$,
    \item If $n^+_{\beta_{\text{max}}} \geq n^+_{\beta'_{\text{max}}}$, then $i_{\beta_{\text{max}}} \leq i_{\beta'_{\text{max}}}$,
\end{itemize}
since the numerators are equal for both indices. Assume without loss of generality that $n^+_{\beta_{\text{max}}} \le n^+_{\beta'_{\text{max}}}$. Then we have:
\[
i_{\beta_{\text{max}}} \ge i_{\beta'_{\text{max}}} > i_\beta, \quad \forall \beta \notin \{\beta_\text{max}, \beta'_\text{max}\}.
\]
From the definition \eqref{eq : def nplus} of $n_\beta^+$, the explicit expressions for the values at the maxima are:
\begin{equation} \label{eq : explicit i iprime}
	\begin{aligned}
		i_{\beta_{\text{max}}}
		&= \frac{\frac{1}{4} \left( \int_\cX \abs{s_{\theta}(x)} p_\theta(x) \mu(dx) \right)^2}
		{1 + (e^\alpha - 1) \int_{F^+_\text{max}} p_\theta(x) \mu(dx)}, \\
		i_{\beta'_{\text{max}}}
		&= \frac{\frac{1}{4} \left( \int_\cX \abs{s_{\theta}(x)} p_\theta(x) \mu(dx) \right)^2}
		{1 + (e^\alpha - 1) \int_{F^{\prime +}_\text{max}} p_\theta(x) \mu(dx)}.
	\end{aligned}
\end{equation}
 
 Now, the rest of the proof consists in showing that the maximization problem \eqref{eq : optimization pb finite} has a maximum that is attained for a sub-probability measure $(\omega_\beta)_{\beta}$ supported on $\{\beta_\text{max},\beta'_\text{max}\}$.

Let us denote by $R$ the staircase pattern matrix of size $d \times 2^d$ as defined in Definition~\ref{def: staircase matrix}. It satisfies
\[
R_{j,\beta} = (r_\beta)_j = e^\alpha \indi{F^+_\beta}(x_j) + \indi{F^-_\beta}(x_j).
\]
This matrix has columns given by the vectors defined in \eqref{eq : R matrix by slice}. Then, the optimization problem \eqref{eq : optimization pb finite} can be rewritten as
\[
M^* = \max_{\omega \in \R^{2^d}_+} \omega^T i \quad \text{subject to} \quad R \omega = \bm{1},
\]
where $i = (i_\beta)_\beta$ and $\bm{1} = [1, \dots, 1]^T \in \R^d$. This problem admits a solution since the constraint defines a compact set. Let $\overline{\omega}$ be a maximizing solution. Define 
\[
q = \#\left\{ \beta \in \{0, \dots, 2^d - 1\} \mid \overline{\omega}_\beta > 0 \right\}.
\]
We distinguish three cases depending on the value of $q$.

\paragraph{$\bullet$ Case $q = 1$.}
In this case, 
%the constraint $R\overline{\omega} = \bm{1}$ implies that
 $\overline{\omega}$ is
 supported on a single index $\bar{\beta}$. The constraint becomes:
\begin{equation}\label{eq: system q=1}
\begin{cases}
1 = (r_{\bar{\beta}})_1 \bar{\omega}_{\bar{\beta}} \\
\vdots \\
1 = (r_{\bar{\beta}})_d \bar{\omega}_{\bar{\beta}}.
\end{cases}
\end{equation}
From the definition of $r_\beta$, for $\bar{\beta} = 0$ we have $(r_{\bar{\beta}})_j = 1$ for all $j$, and for $\bar{\beta} = 2^d - 1$, we have $(r_{\bar{\beta}})_j = e^\alpha$ for all $j$. For other values of $\bar{\beta}$, the vector $r_{\bar{\beta}}$ has mixed entries in $\{1, e^\alpha\}$. Thus, the system~\eqref{eq: system q=1} has a solution only if $\bar{\beta} \in \{0, 2^d - 1\}$. However, for these values, the objective coefficients vanish:
\[
i_\beta = \frac{\left( \int_{F^+_\beta} s_{\theta}(x) p_\theta(x)\mu(dx)\right)^2}{1 + (e^\alpha - 1) n^+_\beta} = 0,
\]
since $F^+_0 = \emptyset$ and $F^+_{2^d - 1} = \mathcal{X}$, and the score function integrates to $0$ over the full space. Hence, $M^* = 0$. We will show that for $q = 2$, $M^* > 0$, so the maximum cannot be attained with $q = 1$.

\paragraph{$\bullet$ Case $q = 2$.}
Assume there exist $0 \le \beta_1 < \beta_2 \le 2^d - 1$ such that $\overline{\omega}_{\beta_1}, \overline{\omega}_{\beta_2} > 0$ and $\overline{\omega}_\beta = 0$ for all other $\beta$. Exclude the case $\beta_1 = 0$, $\beta_2 = 2^d - 1$, since that would again imply $M^* = 0$.

Because $r_{\beta_1} \ne r_{\beta_2}$, there exists $x \in \mathcal{X}$ such that $R_{x,\beta_1} \ne R_{x,\beta_2}$. Without loss of generality, assume $R_{x,\beta_1} = 1$, $R_{x,\beta_2} = e^\alpha$. Since $\beta_1 \ne 0$, we can find $y \in \mathcal{X}$ such that $R_{y,\beta_1} = e^\alpha$ (as $r_0$ is the only vector with all entries equal to $1$). The constraint $R \overline{\omega} = \bm{1}$ yields:
\begin{equation}\label{eq: system 14.5}
\left\{
\begin{aligned}
\overline{\omega}_{\beta_1} + e^\alpha \overline{\omega}_{\beta_2} &= 1 \\
e^\alpha \overline{\omega}_{\beta_1} + R_{y,\beta_2} \overline{\omega}_{\beta_2} &= 1.
\end{aligned}
\right.
\end{equation}
To ensure a solution with both $\overline{\omega}_{\beta_1}, \overline{\omega}_{\beta_2} > 0$, it must be that $R_{y,\beta_2} = 1$. Then the solution is
\[
\overline{\omega}_{\beta_1} = \overline{\omega}_{\beta_2} = \frac{1}{1 + e^\alpha},
\]
and the objective function yields
\[
M^* = \frac{i_{\beta_1} + i_{\beta_2}}{1 + e^\alpha}.
\]
This value is maximal when $i_{\beta_1}$ and $i_{\beta_2}$ are the two largest among the $(i_\beta)_\beta$, hence:
\begin{equation}\label{eq : maximume q is 2}
M^* = \frac{i_{\beta_{\max}} + i_{\beta'_{\max}}}{1 + e^\alpha}.
\end{equation}
Using the explicit form \eqref{eq : explicit i iprime} and the identity
\[
\int_{F^{\prime +}_{\max}} p_\theta(x)\mu(dx) = 1 - \int_{F^{+}_\text{max}} p_\theta(x)\mu(dx),
\]
we deduce that
\[
i_{\beta_{\max}} + i_{\beta'_{\max}} = \frac{1}{4} \mathbb{E}[|s_\theta(x)|]^2 \cdot \frac{1 + e^\alpha}{\left(1 + (e^\alpha - 1)n_{\max}\right)\left(1 + (e^\alpha - 1)(1 - n_{\max})\right)}.
\]
Hence,
\[
M^* = \frac{1}{4} \mathbb{E}[|s_\theta(x)|]^2 \cdot \frac{1}{\left(1 + (e^\alpha - 1)n_{\max}\right)\left(1 + (e^\alpha - 1)(1 - n_{\max})\right)},
\]
which gives the desired expression \eqref{eq : def M star}, noting that
\[
1 + (e^\alpha - 1)(1 - n_{\max}) = e^\alpha(1 - n_{\max}) + n_{\max}.
\]

\paragraph{$\bullet$ Case $q \ge 3$.}
We will show that if $\alpha$ is small enough, then the maximum cannot be attained, leading to a contradiction and implying that the assumption $q \ge 3$ is false.
Let $\beta_{1}, \dots, \beta_{q} \in \{0,\dots,2^d-1\}$ be the indices such that $\overline{\omega}_{\beta_{l}} > 0$ for $l = 1,\dots,q$.

We divide the analysis into several subcases:

\subparagraph{Subcase 1:} Assume that $\{\beta_{\max}, \beta'_{\max}\} \cap \{\beta_{1},\dots,\beta_{q}\} = \emptyset$.  
That is, the support of $\overline{\omega}$ is disjoint from the two indices where $i_\beta$ is maximal. Define
\[
i_{(3)} := \max\left\{ i_{\beta} \mid \beta \in \{0,\dots,2^d-1\} \setminus \{\beta_{\max}, \beta'_{\max} \} \right\}.
\]
Since $i_{\beta_{\max}} \ge i_{\beta'_{\max}} > i_{(3)}$, we obtain
\[
\overline{\omega}^T i = \sum_{l=1}^q \overline{\omega}_{\beta_l} i_{\beta_l} \le i_{(3)} \sum_{l=1}^q \overline{\omega}_{\beta_l} \le i_{(3)},
\]
where we used that $\sum_{\beta=0}^{2^d-1} \overline{\omega}_\beta \le 1$ due to the constraint $R\overline{\omega} = \bm{1}$.  
If $\alpha$ is small enough so that
\[
i_{(3)} \le \frac{i_{\beta_{\max}} + i_{\beta'_{\max}}}{1 + e^\alpha},
\]
then this contradicts the optimality of this subcase, as the value is strictly less than the value achieved in the case $q = 2$ (see \eqref{eq : maximume q is 2}).

\subparagraph{Subcase 2:} Assume $i_{\beta_{\max}} > i_{\beta'_{\max}}$ and $\beta_{\max} \notin \{\beta_1, \dots, \beta_q\}$.  
Then all $\beta_l$ satisfy $i_{\beta_l} \le i_{\beta'_{\max}}$, since the only one potentially exceeding it is excluded. Thus,
\[
\overline{\omega}^T i \le i_{\beta'_{\max}} \sum_{l=1}^q \overline{\omega}_{\beta_l} \le i_{\beta'_{\max}}.
\]
We compare with the maximum value for $q = 2$:
\[
i_{\beta'_{\max}} - \frac{i_{\beta_{\max}} + i_{\beta'_{\max}}}{1 + e^\alpha} 
= \frac{e^\alpha i_{\beta'_{\max}} - i_{\beta_{\max}}}{1 + e^\alpha},
\]
which is negative when $\alpha < \log\left( \frac{i_{\beta_{\max}}}{i_{\beta'_{\max}}} \right)$.  
Hence, for such small $\alpha$, the maximum cannot be attained in this subcase.

\subparagraph{Subcase 3:} Assume that both $\beta_{\max}$ and $\beta'_{\max}$ belong to $\{\beta_1, \dots, \beta_q\}$.  
Without loss of generality, relabel such that $\beta_1 = \beta_{\max}$ and $\beta_2 = \beta'_{\max}$.  

Since $F^+_{\max} \cup F^{\prime +}_{\max} = \mathcal{X}$, there exist points $x \in F^+_{\max} = F^+_{\beta_1}$ and $y \in F^{\prime +}_{\max} = F^+_{\beta_2} = \mathcal{X} \setminus F^+_{\beta_1}$.  
From the definition of the sets $F^+_\beta$, we have:
\[
R_{x,\beta_1} = e^\alpha, \quad R_{x,\beta_2} = 1, \quad R_{y,\beta_1} = 1, \quad R_{y,\beta_2} = e^\alpha.
\]
Writing the corresponding rows of $R \overline{\omega} = \bm{1}$ at points $x$ and $y$ gives:
\begin{equation} \label{eq : syst 2x2 subcase3}
\left\{
\begin{aligned}
e^\alpha \overline{\omega}_{\beta_1} + \overline{\omega}_{\beta_2} + \sum_{l=3}^q R_{x,\beta_l}\overline{\omega}_{\beta_l} &= 1, \\
\overline{\omega}_{\beta_1} + e^\alpha \overline{\omega}_{\beta_2} + \sum_{l=3}^q R_{y,\beta_l}\overline{\omega}_{\beta_l} &= 1.
\end{aligned}
\right.
\end{equation}

Define $s_x := \sum_{l=3}^q R_{x,\beta_l}\overline{\omega}_{\beta_l}$ and $s_y := \sum_{l=3}^q R_{y,\beta_l}\overline{\omega}_{\beta_l}$.  
Solving \eqref{eq : syst 2x2 subcase3} for $\overline{\omega}_{\beta_1}$ and $\overline{\omega}_{\beta_2}$ gives:
\begin{equation} \label{eq : expression omega subcase 3}
\overline{\omega}_{\beta_1} = \frac{1}{1 + e^\alpha} - \frac{s_x e^\alpha - s_y}{e^{2\alpha} - 1}, \quad
\overline{\omega}_{\beta_2} = \frac{1}{1 + e^\alpha} - \frac{s_y e^\alpha - s_x}{e^{2\alpha} - 1}.
\end{equation}

Define $\sigma := \sum_{l=3}^q \overline{\omega}_{\beta_l} > 0$.  
Since all entries of $R$ are either $1$ or $e^\alpha$, we have:
\[
\sigma \le s_x \le e^\alpha \sigma, \quad \sigma \le s_y \le e^\alpha \sigma.
\]
In particular, both $s_x e^\alpha - s_y \ge 0$ and $s_y e^\alpha - s_x \ge 0$. We now compute:
\begin{align*}
\overline{\omega}^T i &= \overline{\omega}_{\beta_1} i_{\beta_1} + \overline{\omega}_{\beta_2} i_{\beta_2} + \sum_{l=3}^q \overline{\omega}_{\beta_l} i_{\beta_l} \\
&= \overline{\omega}_{\beta_1} i_{\beta_{\max}} + \overline{\omega}_{\beta_2} i_{\beta'_{\max}} + \sum_{l=3}^q \overline{\omega}_{\beta_l} i_{\beta_l} \\
&= \frac{i_{\beta_{\max}} + i_{\beta'_{\max}}}{1 + e^\alpha}
- \frac{s_x e^\alpha - s_y}{e^{2\alpha} - 1} i_{\beta_{\max}} 
- \frac{s_y e^\alpha - s_x}{e^{2\alpha} - 1} i_{\beta'_{\max}} 
+ \sum_{l=3}^q \overline{\omega}_{\beta_l} i_{\beta_l} \\
&\le \frac{i_{\beta_{\max}} + i_{\beta'_{\max}}}{1 + e^\alpha}
- \frac{s_x + s_y}{e^\alpha + 1} i_{\beta'_{\max}} 
+ \sum_{l=3}^q \overline{\omega}_{\beta_l} i_{\beta_l}, \quad 
\text{as $i_{\beta_{\max}} \ge i_{\beta'_{\max}}$,} \\
&\le \frac{i_{\beta_{\max}} + i_{\beta'_{\max}}}{1 + e^\alpha}
+ \sigma \left( i_{(3)} - 2 \frac{i_{\beta'_{\max}}}{e^\alpha + 1} \right),
\end{align*}
where we used $s_x + s_y \ge 2 \sigma$ and $i_{\beta_l} \le i_{(3)}$ for $l \ge 3$.

As $i_{(3)} < i_{\beta'_{\max}}$, choosing $\alpha$ small enough ensures that 
\[
i_{(3)} - 2 \frac{i_{\beta'_{\max}}}{e^\alpha + 1} < 0,
\]
thus the final term is negative and the value of $\overline{\omega}^T i$ is strictly less than $\frac{i_{\beta_{\max}} + i_{\beta'_{\max}}}{1 + e^\alpha}$, hence the maximum is not attained in this subcase.

\subparagraph{Subcase 4:} The complement of the previous subcases.  
By exclusion of the subcases 1 and 3, exactly one of $\beta_{\max}$ or $\beta'_{\max}$ appears among $\{\beta_1, \dots, \beta_q\}$.
%Assume that exactly one of $\beta_{\max}$ or $\beta'_{\max}$ appears among $\{\beta_1, \dots, \beta_q\}$ and the assumptions of Subcases 1–3 do not hold.  
From exclusion of the subcase 2, there exists an index $l_1$ such that
\[
i_{\beta_{l_1}} = \max_{\beta} i_\beta = i_{\beta_{\max}}, \quad \text{and } i_{\beta_l} \le i_{(3)} \text{ for all } l \neq l_1.
\]

The proof proceeds analogously to Subcase 3. One extracts a $2 \times 2$ system from $R \overline{\omega} = \bm{1}$ involving $\beta_{l_1}$ and some $\beta_{l_2} \neq \beta_{l_1}$, provided the geometric condition
\begin{equation} \label{eq : subcase 4 finite cond on F}
F^+_{\beta_{l_1}} \cap \left(F^+_{\beta_{l_2}}\right)^c \neq \emptyset, \quad 
\left(F^+_{\beta_{l_1}}\right)^c \cap F^+_{\beta_{l_2}} \neq \emptyset,
\end{equation}
holds. This is stated and proven in Lemma~\ref{lem : for subcase 4 finite cond on F} below.


In particular, let $x \in \left( F^{+}_{\beta_{l_2}} \right)^c \cap F^+_{\beta_{l_1}}$ and $y \in F^+_{\beta_{l_2}} \cap \left( F^{+}_{\beta_{l_1}} \right)^c$.  
By the definition of the sets $F^+_\beta$, we have:
\[
R_{x,\beta_{l_1}} = e^\alpha, \quad R_{x,\beta_{l_2}} = 1, \quad R_{y,\beta_{l_1}} = 1, \quad R_{y,\beta_{l_2}} = e^\alpha.
\]

We now write the two equations of the system $R\overline{\omega} = \bm{1}$ corresponding to the rows indexed by $x$ and $y$. Similarly to \eqref{eq : syst 2x2 subcase3}, we obtain:
\begin{equation*}
\left\{
\begin{aligned}
e^\alpha \, \overline{\omega}_{\beta_{l_1}} + \overline{\omega}_{\beta_{l_2}} 
+ \sum_{\substack{l=1 \\ l \notin \{l_1,l_2\}}}^q R_{x,\beta_l} \, \overline{\omega}_{\beta_l} &= 1, \\
\overline{\omega}_{\beta_{l_1}} + e^\alpha \, \overline{\omega}_{\beta_{l_2}} 
+ \sum_{\substack{l=1 \\ l \notin \{l_1,l_2\}}}^q R_{y,\beta_l} \, \overline{\omega}_{\beta_l} &= 1.
\end{aligned}
\right.
\end{equation*}

Let us denote the remaining sums by:
\[
\sigma_x := \sum_{\substack{l=1 \\ l \notin \{l_1,l_2\}}}^q R_{x,\beta_l} \, \overline{\omega}_{\beta_l}, \qquad 
\sigma_y := \sum_{\substack{l=1 \\ l \notin \{l_1,l_2\}}}^q R_{y,\beta_l} \, \overline{\omega}_{\beta_l}.
\]

Then, as in \eqref{eq : expression omega subcase 3}, we deduce:
\begin{equation*}
\overline{\omega}_{\beta_{l_1}} = \frac{1}{1 + e^\alpha} - \frac{\sigma_x e^\alpha - \sigma_y}{e^{2\alpha} - 1}, \qquad
\overline{\omega}_{\beta_{l_2}} = \frac{1}{1 + e^\alpha} - \frac{\sigma_y e^\alpha - \sigma_x}{e^{2\alpha} - 1}.
\end{equation*}

We now compute:
\begin{align*}
\overline{\omega}^T i &= \overline{\omega}_{\beta_{l_1}} i_{\beta_{l_1}} + \overline{\omega}_{\beta_{l_2}} i_{\beta_{l_2}} 
+ \sum_{\substack{l=1 \\ l \notin \{l_1,l_2\}}}^q \overline{\omega}_{\beta_l} i_{\beta_l} \\
&\le \overline{\omega}_{\beta_{l_1}} i_{\beta_{\max}} + \overline{\omega}_{\beta_{l_2}} i_{\beta'_{\max}} 
+ \sum_{\substack{l=1 \\ l \notin \{l_1,l_2\}}}^q \overline{\omega}_{\beta_l} i_{\beta_l},
\end{align*}
where we used that $i_{\beta_{l_1}} = i_{\beta_{\max}}$ and $i_{\beta_{l_2}} \le i_{(3)} \le i_{\beta'_{\max}}$. Then, exactly as in Subcase~3, we deduce that
\[
\overline{\omega}^T i < \frac{i_{\beta_{\max}} + i_{\beta'_{\max}}}{1 + e^\alpha}.
\]
Consequently, the maximum $M^*$ cannot be attained in Subcase~4.
\medskip 

In all cases, we obtain a contradiction to the assumption that $q \ge 3$ yields the maximum. Hence, the maximum is reached only when $q = 2$.

\medskip

This concludes the proof, up to showing that the geometric condition \eqref{eq : subcase 4 finite cond on F} is satisfied for some $\beta_{l_2}$.  
This is verified in Lemma~\ref{lem : for subcase 4 finite cond on F} below.




\end{proof}




\begin{lem}\label{lem : for subcase 4 finite cond on F}
	There exists some $l_2 \in \{1,\dots,q\}$, $l_2\neq l_1$ such that,
	\begin{equation*} 
		\left( F^{+}_{\beta_{l_2}}\right)^c \cap  F^+_{\beta_{l_1}}  \neq \emptyset \quad \text{and} \quad 	F^+_{\beta_{l_2}} \cap \left( F^{+}_{\beta_{l_1}}\right)^c \neq \emptyset.
	\end{equation*}
\end{lem}
\begin{proof}
We proceed by contradiction. Suppose the statement does not hold. Then, for all $l \neq l_1$, we must have:
\begin{equation*} 
    \left( F^{+}_{\beta_l} \right)^c \cap F^+_{\beta_{l_1}} = \emptyset 
    \quad \text{or} \quad 
    F^+_{\beta_l} \cap \left( F^+_{\beta_{l_1}} \right)^c = \emptyset,
\end{equation*}
which implies:
\begin{equation} \label{eq: F subset 18.5}
    F^+_{\beta_{l_1}} \subset F^+_{\beta_l}
    \quad \text{or} \quad
    F^+_{\beta_l} \subset F^+_{\beta_{l_1}}.
\end{equation}
Let us choose some $x \in F^+_{\beta_{l_1}}$ and $y \in \left( F^+_{\beta_{l_1}} \right)^c$.  
This is possible because $F^+_{\beta_{l_1}} \neq \emptyset$ and $F^+_{\beta_{l_1}} \neq \cX$,  
otherwise we would have $i_{\beta_{\max}} = i_{\beta_{l_1}} = 0$, which contradicts $M^* > 0$.

By definition of the matrix $R$, we know that:
\[
R_{x,\beta_{l_1}} = e^\alpha, \qquad R_{y,\beta_{l_1}} = 1.
\]
Now fix any $l \in \{1, \dots, q\} \setminus \{l_1\}$. Then, one of the two alternatives in \eqref{eq: F subset 18.5} must occur:

- Case 1: If $F^+_{\beta_{l_1}} \subset F^+_{\beta_l}$, then since $x \in F^+_{\beta_{l_1}}$, we also have $x \in F^+_{\beta_l}$, hence:
\[
R_{x,\beta_l} = e^\alpha \ge R_{y,\beta_l} \in \{1, e^\alpha\}.
\]

- Case 2: If $F^+_{\beta_l} \subset F^+_{\beta_{l_1}}$, then since $y \notin F^+_{\beta_{l_1}}$, it follows that $y \notin F^+_{\beta_l}$, hence:
\[
R_{y,\beta_l} = 1 \le R_{x,\beta_l} \in \{1, e^\alpha\}.
\]
In both cases, we conclude:
\[
R_{y,\beta_l} \le R_{x,\beta_l} \quad \text{for all } l \in \{1, \dots, q\} \setminus \{l_1\}, \qquad
R_{y,\beta_{l_1}} = 1 < e^\alpha = R_{x,\beta_{l_1}}.
\]
This leads to a contradiction with the strict positivity of the weights $\overline{\omega}_{\beta_l}$ and the fact that:
\begin{equation*}
    \sum_{l=1}^q R_{y,\beta_l} \, \overline{\omega}_{\beta_l} = 1 = 
    \sum_{l=1}^q R_{x,\beta_l} \, \overline{\omega}_{\beta_l}.
\end{equation*}
The proof is therefore complete.

\end{proof}



\appendix


 
\begin{thebibliography}{9}

\bibitem{1LDP} Abadi, M., Chu, A., Goodfellow, I., McMahan, H. B., Mironov, I., Talwar, K., \& Zhang, L. (2016, October). Deep learning with differential privacy. In Proceedings of the 2016 ACM SIGSAC conference on computer and communications security (pp. 308-318).

\bibitem{AmoGlo23} Amorino, C., \& Gloter, A. (2025). Minimax rate for multivariate data under componentwise local differential privacy constraints. The Annals of Statistics, 53(3), 1176-1202.

\bibitem{AmoGloHal} Amorino, C., Gloter, A., \& Halconruy, H. (2025). Evolving privacy: drift parameter estimation for discretely observed iid diffusion processes under LDP. Stochastic Processes and their Applications, 181, 104557.

\bibitem{5LDP} Apple Differential Privacy Team. (2017). Learning with privacy at scale. Apple Mach. Learn. J, 1(8), 1-25. Available at https://machinelearning.apple.com/2017/12/06/learning-withprivacy-at-scale.html.

\bibitem{Asi22} Asi, H., Feldman, V., \& Talwar, K. (2022, June). Optimal algorithms for mean estimation under local differential privacy. In International Conference on Machine Learning (pp. 1046-1056). PMLR.

\bibitem{Avella} Avella-Medina, M., Bradshaw, C., \& Loh, P. L. (2023). Differentially private inference via noisy optimization. The Annals of Statistics, 51(5), 2067-2092.

\bibitem{BarChe20} Barnes, L. P., Chen, W. N., \& Özgür, A. (2020). Fisher information under local differential privacy. IEEE Journal on Selected Areas in Information Theory, 1(3), 645-659.

\bibitem{4CLDP} Berrett, T., \& Butucea, C. (2020). Locally private non-asymptotic testing of discrete distributions is faster using interactive mechanisms. Advances in Neural Information Processing Systems, 33, 3164-3173.

\bibitem{6CLDP} Berrett, T., \& Yu, Y. (2021). Locally private online change point detection. Advances in Neural Information Processing Systems, 34, 3425-3437.

\bibitem{Bic93} Bickel, P. et al (1993). Efficient and adaptive estimation for semiparametric models (Vol. 4). Baltimore: Johns Hopkins University Press.

\bibitem{Bre83} Brezis, H. (1983). Analyse fonctionnelle. Théorie et applications.

\bibitem{BunSte16} Bun, M., \& Steinke, T. (2016, October). Concentrated differential privacy: Simplifications, extensions, and lower bounds. In Theory of cryptography conference (pp. 635-658). Berlin, Heidelberg: Springer Berlin Heidelberg.

\bibitem{10CLDP} Butucea, C., Dubois, A., Kroll, M., \& Saumard, A. (2020). Local differential privacy: Elbow effect in optimal density estimation and adaptation over Besov ellipsoids.

\bibitem{ButKlo25} Butucea, C., Klockmann, K., \& Krivobokova, T. (2025). Nonparametric spectral density estimation using interactive mechanisms under local differential privacy. arXiv preprint arXiv:2504.00919.


\bibitem{11CLDP} Butucea, C., Rohde, A., \& Steinberger, L. (2023). Interactive versus noninteractive locally differentially private estimation: Two elbows for the quadratic functional. The Annals of Statistics, 51(2), 464-486.


\bibitem{Choquet} Choquet, G. (1960). Le théorème de représentation intégrale dans les ensembles convexes compacts. In Annales de l'institut Fourier (Vol. 10, pp. 333-344).

%\bibitem{Dub20} Dubost, C., Humbert, P., Oudre, L., Labourdette, C., Vayatis, N., \& Vidal, P. P. (2021). Quantitative assessment of consciousness during anesthesia without EEG data. Journal of Clinical Monitoring and Computing, 35, 993-1005.

\bibitem{DieUhl_Book} Diestel, J. and Uhl, J.J. (1977). Vector Measures. American Mathematical Soc. ISBN 978-0-8218-1515-1


\bibitem{Martin} Duchi, J. C., Jordan, M. I., \& Wainwright, M. J. (2018). Minimax optimal procedures for locally private estimation. Journal of the American Statistical Association, 113(521), 182-201.

\bibitem{DucRog19} Duchi, J., \& Rogers, R. (2019, June). Lower bounds for locally private estimation via communication complexity. In Conference on Learning Theory (pp. 1161-1191). PMLR.


\bibitem{DucRua24} Duchi, J. C., \& Ruan, F. (2024). The right complexity measure in locally private estimation: It is not the fisher information. The Annals of Statistics, 52(1), 1-51.


\bibitem{Dwo06} Dwork, C., McSherry, F., Nissim, K., \& Smith, A. (2006). Calibrating noise to sensitivity in private data analysis. In Theory of Cryptography: Third Theory of Cryptography Conference, TCC 2006, New York, NY, USA, March 4-7, 2006. Proceedings 3 (pp. 265-284). Springer Berlin Heidelberg.

\bibitem{DwoRot16} Dwork, C., \& Rothblum, G. N. (2016). Concentrated differential privacy. arXiv preprint arXiv:1603.01887.

\bibitem{23LDP} Erlingsson, Ú., Pihur, V., \& Korolova, A. (2014, November). Rappor: Randomized aggregatable privacy-preserving ordinal response. In Proceedings of the 2014 ACM SIGSAC conference on computer and communications security (pp. 1054-1067).

\bibitem{Gen15} Geng, Q., Kairouz, P., Oh, S., \& Viswanath, P. (2015). The staircase mechanism in differential privacy. IEEE Journal of Selected Topics in Signal Processing, 9(7), 1176-1184.

\bibitem{GenVis12} Geng, Q., \& Viswanath, P. (2014). The optimal mechanism in differential privacy. In 2014 IEEE international symposium on information theory (pp. 2371-2375). IEEE.

\bibitem{GenVis13} Geng, Q., \& Viswanath, P. (2015). Optimal noise adding mechanisms for approximate differential privacy. IEEE Transactions on Information Theory, 62(2), 952-969.

\bibitem{Ghosh12} Ghosh, A., Roughgarden, T., \& Sundararajan, M. (2009, May). Universally utility-maximizing privacy mechanisms. In Proceedings of the forty-first annual ACM symposium on Theory of computing (pp. 351-360).

\bibitem{Jos19} Joseph, M., Kulkarni, J., Mao, J., \& Wu, S. Z. (2019). Locally private gaussian estimation. Advances in Neural Information Processing Systems, 32

%\bibitem{Jun21} Jung, S., Oudre, L., Truong, C., Dorveaux, E., Gorintin, L., Vayatis, N., \& Ricard, D. (2021, November). Adaptive change-point detection for studying human locomotion. In 2021 43rd Annual International Conference of the IEEE Engineering in Medicine \& Biology Society (EMBC) (pp. 2020-2024). IEEE.

\bibitem{kairouzExtremalMechanismsLocal2015} Kairouz, P., Oh, S., \& Viswanath, P. (2016). Extremal mechanisms for local differential privacy. Journal of Machine Learning Research, 17(17), 1-51.

\bibitem{KalSte} Kalinin, N. P., \& Steinberger, L. (2024). Efficient estimation of a Gaussian mean with Local Differential Privacy. arXiv preprint arXiv:2402.04840.

\bibitem{KeaRot19} Kearns, M., \& Roth, A. (2019). The ethical algorithm: The science of socially aware algorithm design. Oxford University Press.

\bibitem{puff} Kifer, D., \& Machanavajjhala, A. (2014). Pufferfish: A framework for mathematical privacy definitions. ACM Transactions on Database Systems (TODS), 39(1), 1-36. 

\bibitem{Kro24} Kroll, M. (2024). Nonparametric spectral density estimation under local differential privacy. Statistical Inference for Stochastic Processes, 27(3), 725-759.

\bibitem{Kul23} Kulesza, A., Suresh, A. T., \& Wang, Y. (2023). Mean estimation in the add-remove model of differential privacy. arXiv preprint arXiv:2312.06658.


\bibitem{Lang_Book} Lang, S. (1993). Real and Functional Analysis. Graduate Texts in Mathematics. Springer New York. ISSN: 0072-5285
	
\bibitem{Lal23} Lalanne, C. (2023). On the tradeoffs of statistical learning with privacy (Doctoral dissertation, Ecole normale supérieure de lyon-ENS LYON).

\bibitem{29CLDP} Lam-Weil, J., Laurent, B., \& Loubes, J. M. (2022). Minimax optimal goodness-of-fit testing for densities and multinomials under a local differential privacy constraint. Bernoulli, 28(1), 579-600.

\bibitem{31CLDP} Li, M., Berrett, T. B., \& Yu, Y. (2023). On robustness and local differential privacy. The Annals of Statistics, 51(2), 717-737.

\bibitem{NamLee22} Nam, S. H., \& Lee, S. H. (2022). A tighter converse for the locally differentially private discrete distribution estimation under the one-bit communication constraint. IEEE Signal Processing Letters, 29, 1923-1927.

\bibitem{Nar08} Narayanan, Arvind, and Vitaly Shmatikov. Robust de-anonymization of large sparse datasets. 2008 IEEE Symposium on Security and Privacy (sp 2008). IEEE, 2008.

\bibitem{Oxt} Oxtoby, J. C. (1970). Homeomorphic measures in metric spaces. Proceedings of the American Mathematical Society, 24(3), 419-423.

\bibitem{Smi08} Smith, A. (2008). Efficient, differentially private point estimators. arXiv preprint arXiv:0809.4794.

\bibitem{Smi11} Smith, A. (2011, June). Privacy-preserving statistical estimation with optimal convergence rates. In Proceedings of the forty-third annual ACM symposium on Theory of computing (pp. 813-822).

\bibitem{49LDP} Tang, J., Korolova, A., Bai, X., Wang, X., \& Wang, X. (2017). Privacy loss in apple's implementation of differential privacy on macos 10.12. arXiv preprint arXiv:1709.02753.

%\bibitem{TruOud22} Truong, C., \& Oudre, L. (2022, December). Supervised change-point detection with dimension reduction, applied to physiological signals. In NeurIPS 2022 Workshop on Learning from Time Series for Health.

\bibitem{steinbergerEfficiencyLocalDifferent} Steinberger, L. (2024). Efficiency in local differential privacy. The Annals of Statistics, 52(5), 2139-2166.

\bibitem{Van07} Van der Vaart, A. W. (2000). Asymptotic statistics (Vol. 3). Cambridge university press.

\bibitem{WassermanZhou10} Wasserman, L., \& Zhou, S. (2010). A statistical framework for differential privacy. Journal of the American Statistical Association, 105(489), 375-389.

\end{thebibliography}

\end{document}
